% $Id$
% $Log$
% Revision 1.7  2012/08/24 16:20:13  dstotler
% Corrected error in Eq. (2.59): n_e factor missing in He line
% radiation rate.  Add link to the Bateman distribution report.
% Revised the discussion of transfers between BGK species in
% connection with the CMod-1D example to reflect the current
% (not yet corrected) state of the BGK data files.  Updated
% run time information for Couette flow example. Added ucd_plot
% in same places where geomtesta appears.  Noted problem with
% ncgen and netcdf 4.1.1.
%
% Revision 1.6  2009/11/25 20:50:52  dstotler
% Extensive update to cover current operating systems and compilers,
% as well as DG.  Add the two new definegeometry2d /
% defineback examples.
%
% Revision 1.5  2007/06/22 16:32:02  dstotler
% Update particle balance analysis for CMod_1D example for changes
% associated with switch to ehr3.dat.
%
% Revision 1.4  2007/06/19 16:54:10  dstotler
% Fixed small bug in CMod-1D example.
%
% Revision 1.3  2006/01/19 17:42:54  dstotler
% Few small changes with implementation of 1-D geometry to Couette flow
% problem. More detailed discussion of geometry options, including 3-D.
% Add Historical Background section.  Update some neutral-neutral
% scattering text.  Revise "atomic physics" section to discuss both
% hydrogen and helium CR models.
%
% Revision 1.2  2001/11/16 16:41:18  dstotler
% Extensive revisions, particularly the inclusion of elastic scattering
% background and examples.
%
% Revision 1.1  2001/07/16 20:30:17  dstotler
% Adding documentation directory to CVS.
%
% To process this file:
% 
% Check DEGAS 2 version tag below and update accordingly.
% (in tex directory), do the following with 
% codefile = ratecalc, readgeometry, random, classes, geometry, tallysetup,
% elementsetup, speciesetup, reactionsetup, materialsetup, pmisetup,
% problemsetup, outputbrowser, definegeometry2d, defineback, geomtesta,
% ucd_plot, postdetector, gpicamera:
%   gmake codefile.tex
%   pdflatex codefile.tex (may need to be done two or three times)
%   mv codefile.pdf ../Doc
% update page numbers in href's to these files; 
% use the actual numbers printed
% on the page for this, not the page number in the PDF.  
% Also, the DEGAS 2
% web page now has links to pages in this manual; those also
% need to be checked.
% rm degas2_all.html
% Remove old {theindex} from this file
% tth -e0 -Ldegas2_all < degas2_all.tex > degas2_all.html
% (this reads but does not update the degas2_all.aux file, so should
%  have that synched with the .tex file before running tth)
% In html file, replace ">Figure with .jpg">Figure
% (also cleaned up title of html file)
% pdflatex degas2_all.tex
% makeindex degas2_all
%
% Insert degas2_all.ind at end of this file. How about before TOC?
%
% pdflatex degas2_all.tex (will now have a hyperlinked index)
%
%  acroread -toPostScript degas2_all.pdf
%  (pdf2ps DOES NOT appear to work well)
%  In Acrobat, was able to do Save As... / PostScript.
%
%
\documentclass[12pt]{report}
% Macro for http reference inclusion, per hypertex:
\def\href#1#2{\special{html:<a href="#1">}{#2}\special{html:</a>}}
\makeindex
\usepackage{times}
% Use this line with pdflatex. The problem with this is that it
% doesn't use eps or ps formats with includegraphics. It
% will allow pdf or jpg.
% this is different from the usage with latex!
% Nonetheless, ``distill file.ps'' produces a pdf file that looks
% much better than a jpeg from xv.  Can use ``convert file.pdf file.jpg''
% (e.g., on tern) to nicely convert these pdf's to jpeg's.
%
\usepackage[pdftex]{hyperref}
\usepackage[pdftex]{graphicx}
% Do we need:
%
% \hyperbaseurl{http://w3.pppl.gov/degas2/Doc/} ??
%
% Use the following with latex, dvips -Pwww, and then distill:
% The problem with this is that the \tt titles choke distill.
% \usepackage[dvips]{hyperref}
% \usepackage{graphicx}
\begin{document}
% \input epsf
% The version number WAS a cvs command, but decide to update manually.
% If we're going to do that, might as well use the actual tag
% for the code.
\title{\includegraphics{degas2d}
User's Guide for DEGAS 2 \\
Release V.\ 4.9\thanks{This file is written in \TeX. 
A hyper-linked PDF
file is generated from that source using {\tt pdflatex}. This document
is distributed as part of the DEGAS 2 code. The user can be best assured
of concurrence between code and manual by referring to the 
version of the manual
that came with the code. 
The DEGAS 2 code itself is licensed under the terms spelled out
in the \href{http://theorycodes.pppl.wikispaces.net/file/view/TheoryCodeLicenseReleaseForm.pdf/249471929/TheoryCodeLicenseReleaseForm.pdf}{user agreement}.
Permission is granted to make and distribute verbatim copies of this manual
provided the copyright and permission notice is preserved on all copies.
Permission is granted to copy and distribute modified versions of this
manual under the conditions for verbatim copying, provided also that
the entire resulting derived work is distributed under the terms of a
permission notice identical to this one.
Permission is granted to copy and distribute translations of this manual
into another language, under the above conditions for modified versions.}}
\author{\vspace*{-0.5in} \\ Daren Stotler, Charles Karney, and George Wilkie \\ {\em Princeton Plasma Physics 
Laboratory} \\ Randall Kanzleiter \\ {\em Rensselaer Polytechnic Institute} \\
Jaishankar S. \\ {\em Institute for Plasma Research, India}}
\begin{abstract}
This is the user's manual for DEGAS 2 - A
Monte Carlo code for the study of neutral atom and molecular transport 
in confined plasmas. It is intended to provide an
introduction to DEGAS 2 from the
user's point-of-view: obtaining
the files, compiling them, setting up the input files, executing the
code, and interpreting the output. 
References will be provided for the underlying physics, but the
essential aspects will be highlighted. 
More detailed documentation of the procedures used to
create input files is contained in the typeset
source code of the relevant pre-processors; links
to those files are provided.
\end{abstract}
\maketitle

\chapter{Introduction}

\section{Purpose of Monte Carlo Neutral Transport}
\index{Neutral Transport}
\index{Transport codes}
\index{Theory, transport}
\index{Boltzmann equation}

The interaction of neutral atoms and molecules with the background 
plasma is an important component in the physics of fusion research
plasmas. These
neutrals affect both the particle and energy balance  of the plasma,
providing a source of new plasma and a channel for heat transport across
the magnetic field that confines the charged particles. In addition, hot
neutrals which leave the plasma volume can interact with the wall of the
plasma chamber, sputtering impurities into the plasma and (in reactors)
possibly damaging the wall. Finally, these hot neutrals can also be used
for plasma diagnostics. An accurate
description of the neutral behaviour would, thus, have be useful
in all of these areas.

The theory of neutral particle kinetics\cite{Tend87}
treats the transport of mass, momentum, and energy in a plasma due to
neutral particles that are themselves unaffected by magnetic
fields.  Though analytic models (solving the
Boltzmann equation for the neutrals or using the diffusion approximation),
both single- and multi-species, have been used to study neutral transport,
a variety of approximations must be made. Some numerical methods for solving
the Boltzmann equation employ simplified geometries and atomic physics. On
the other hand, the Monte Carlo approach to integrating the Boltzmann
equation can treat in great detail three-dimensional geometries, as well as complex
atomic physics and wall interactions. Because of the similarity in the
problems, Monte Carlo neutral transport codes can build on the techniques
developed for neutron transport in reactor design\cite{Span69}.

\section{Historical Background}
\index{DEGAS}

The predecessor to DEGAS 2 was  DEGAS\cite{Heif82}. The origin of the
``DEGAS'' name can best be explained by pointing out that the first
Monte Carlo neutral transport code written by D.\ B.\ Heifetz was
named SEURAT. Its name was a direct reference to the artist Georges
Seurat. This appellation was inspired by the analogy between the
artist's pointillist style and the nature of the Monte Carlo algorithm.
Up close, one sees that Seurat's paintings consist of a huge number of very
tiny dots of paint of a variety of colors.  Only when viewed from a distance
is the image and its coloring clearly apparent. In much the same way,
the track of an individual neutral flight in the Monte Carlo code
gives relatively little insight into the problem being studied. Yet, when the
results of many flights are compiled together, an accurate solution can
be obtained.

Heifetz' fondness for Impressionist artists was also manifested
in succeeding codes like MONET and DEGAS. The latter name was further
motivated by two associated puns. One is that ``DEGAS'' is an
abbreviation of ``DEuterium GAS''; DEGAS was written to satisfy a
request for a code that could treat deuterium as well as hydrogen.
The other is that ``DEGAS'' is just short for ``degassing''; another
proposed application for the code was in assessing measurements of
wall degassing in fusion devices. 

For those who insist that a code's name must be an acronym, we will
reply that ``DEGAS'' stands for Divertor and Edge Gas Analysis System.

By the time the authors began contemplating the concepts that led
to DEGAS 2, the original DEGAS code was in widespread use. It also
enjoyed a very good reputation as being easy to use and well documented.
For this reason, we were reluctant to abandon the DEGAS name, even
though we knew we would be rewriting the code largely from scratch.
Nonetheless, much of the atomic and surface physics from
the original code has been carried over. Likewise, many of the
numerical techniques of DEGAS served as inspiration for
those in DEGAS 2.

\section{Need for DEGAS 2}
\index{DEGAS}

Two events pointed out limitations in the original DEGAS code\cite{Heif82}.
\begin{itemize}
\item The combined DEGAS-B2 code was too slow 
(about 100 Cray hours)\cite{Stot92}. 
The stronger the coupling between neutral and plasma species, the more 
difficult it is to arrive at a converged solution. In this case which 
looked at a high recycling
divertor for a compact, high field reactor, hundreds or thousands of complete
Monte Carlo profiles must be carried out in order to couple effectively
to a fluid code. 
\item The problem physics was difficult to modify. In the radiative and 
detached divertor
regimes being investigated for future reactors, it is believed that the list
of important atomic processes will grow to include elastic collisions and
molecular hydrogenic species other than ground state H$_{2}$.
\end{itemize}

Hence, it was decided that a new code should be written which draws
on the extensive experience at PPPL and that from Garching and
J\"{u}lich.
The philosophy behind DEGAS 2 was defined during 1992-1993:
\begin{itemize}
\item  High speed, 
\item  Flexibility,
\item  Ease-of-use,
\item  Well documented.
\end{itemize}


\section{DEGAS 2 Features}
\index{Object oriented programming}
\index{Allocation, memory}
\index{Macros}
\index{FWEB}

An effort was made to realise the above mentioned objectives of
DEGAS 2 by the following means:
\begin{itemize}
\item Improve speed relative to DEGAS by utilizing the 
``track length'' estimator.
\item Ensure flexibility by designing the code using quasi-object oriented
programming techniques. The modifier ``quasi'' refers to the fact that DEGAS 2
is written in FORTRAN (compatible with FORTRAN 77 and FORTRAN 90). Even though
the basic ``classes'' in the problem have been abstracted, the degree of
encapsulation actually achieved (via pre-processor macros) in DEGAS 2 is 
limited. Furthermore, there is no provision for inheritance within this
programming environment.
\item Provide an easy-to-use and well documented code by writing it with the
FWEB\cite{FWEB} package. FWEB provides the powerful preprocessing and
macro facility which underlies the ``quasi-object oriented''  implementation
and internal \TeX\ documentation.
\end{itemize}

\index{NetCDF}
\index{HDF}
\index{CVS}
\index{Processing, parallel}
Other significant features of DEGAS 2 include:
\begin{itemize}
\item Dynamic memory allocation for optimal performance in large problems,
\item Operation on a number of UNIX platforms, compilable with both 
FORTRAN 77 and FORTRAN 90,
\item Input and output data stored in self-describing, platform-independent 
binary files.
Presently, the main code uses exclusively
netCDF\cite{nCDF} (Network Common Data
Format) files for input and output. However, a simple post-processing
utility is included which generates graphics and data in  
HDF\cite{HDF} (Hierarchial Data Format) or Silo\cite{Silo} files. 
Hopefully, a single format will be settled upon eventually.
\item Parallel execution on distributed workstations and massively
parallel computers. This is the most promising
technique for improving the speed of DEGAS 2.  The reason is that most of
the future advancements in computer speed will result from increases in the
number of processors. Parallel processing is practical to Monte Carlo since
the amount of communication required per computation is small.
\item High quality coding standards.
Good programming practices such as use of the {\tt implicit none} statement
and an assertion facility help to ensure that the code works properly. 
This document will eventually describe a basic test suite which actually
demonstrates that the code is correct. In the code design process, clarity
has been given a high priority as well, facilitating readability and
maintainability.
\item Code versions numbered and documented with the ``CVS'' (concurrent
version system) package.
\end{itemize}

\chapter{Background}

\section{Generic Monte Carlo Algorithms}
\index{Monte Carlo}
\index{Tracking}
\index{Random}
\index{Scoring}
\index{Distribution, probability}
\index{Generators, random number}

Monte Carlo methods generally involve sampling from some ``parent population''.
The sampling process is usually facilitated by pseudo-random numbers generated
via a numerical algorithm. The quantity or quantities sampled are then 
used as input to a subsequent (deterministic) calculation which results in the
desired output. This process is repeated a number of times, so as to 
thoroughly sample the parent population, and a distribution of the output
data is compiled. In some cases, only the mean of this distribution is of 
interest (as well as some estimate of its error), while in other cases, the
distribution itself is the objective.

Note that the process being simulated need not be random. For example, Monte
Carlo techniques may be the only way of computing the volume of some complex
three-dimensional object. This is just one example of the versatility of
Monte Carlo techniques. But, Monte Carlo is not always the most efficient
solution for a problem because of the significant computational
effort required. Two reasons for using Monte Carlo are: 
\begin{enumerate}
\item No other solution is available.
\item The details of the problem are sufficiently important that no 
approximations
can be tolerated.
\end{enumerate}

In the case of particle transport, the latter reason is usually invoked.
Examples of Monte Carlo particle transport include:
\begin{enumerate}
\item Neutron tranport in fission devices,
\item Radiation transport,
\item Electron transport in semiconductors,
\item Neutral particle transport in plasmas.
\end{enumerate}

\subsection{Sampling a Distribution}

The most fundamental process of a Monte Carlo simulation is the sampling
of probability distributions. Typical examples for a particle simulations
would be:
\begin{enumerate}
\item Choosing a velocity vector from a Maxwellian distribution,
\item Picking a particle from a spatially varying source,
\item Determining the distance to the next collision.
\end{enumerate}

\subsection{Sampling Distance to Next Collision}
\label{SamplingDistance}

To illustrate the third case, consider particles with a total, 
macroscopic cross
section $\Sigma_{t}$. We write the  {\em probability density function} 
for having 
the next collision between $x$ and $x + dx$, $p(x)  \: dx$, as:
\begin{equation}
p(x) dx = \Sigma_{t} \exp(-\Sigma_{t} x)   \: dx.
\end{equation}
Then, the  {\em cumulative probability distribution} $P(y)$,
\begin{equation}
P(y) = \int_{0}^{y} p(x)  \: dx, 
\end{equation}
is the probability of having a collision for any $x \leq y$. Note that
$P(0) = 0$ and $P(\infty) = 1$.

\begin{figure}
% \epsfverbosetrue
\begin{center}
\leavevmode
% \epsfxsize=3in
% \epsfbox{MC_Sampling2.eps}
% \includegraphics[width=3in]{MC_Sampling2.eps}
\includegraphics[width=4in]{MC_Sampling2}
\caption{A uniformly sampled $\xi$ value is translated into a sample of $y$
distributed according to the ``Probability Density'' curve using the integrated
``Probability'' curve.}
\label{MCSampling}
\end{center}
\end{figure}

As suggested by the above figure, $P(y)$ can be sampled by choosing a 
uniformly distributed random number $\xi$, with $0 \leq \xi < 1$. For
this example, $P(y)$ is easily computed and inverted so that
\begin{equation}
 y = - \frac{\ln \xi}{\Sigma_{t}},
\end{equation}
where we have surreptitiously replaced $1 - \xi$ with $\xi$ since both
are uniform over the unit interval. Again the logic of this process,
as depicted in the figure, is that if $\xi$ is
uniformly distributed, then numbers between $\xi$ to $\xi + \Delta \xi$ 
are chosen with the same frequency as particles having collisions between
$y$ and $y + \Delta y$. Thus, $y$ is the sampled
  {\em distance to next collision}.

\subsection{Sampling More Complicated Distributions}

Because of the simplicity of the functional form of the collision probability
density function, the distance to the next collision can be easily sampled
using a single, uniform random number. If this were not the case, more
general methods, such as this \emph{rejection technique}, would have 
to be employed.

As an example, consider the sampling of a velocity $v$ from a Maxwellian 
distribution,
\begin{equation}
p(v) = \frac{1}{\sqrt{\pi} v_{\rm th}} \exp \left[ - \frac{(v -
v_{0})^{2}}{v_{\rm th}^{2}} \right].
\end{equation}
For simplicity, limit the $v$ values to $|v - v_{0}| < M v_{\rm th}$ for some
appropriate $M$.

\begin{figure}
% \epsfverbosetrue
\begin{center}
\leavevmode
\includegraphics[width=4in]{MC_Rejection2}
\caption{Values of $x$ distributed according to $p(x)$ can be sampled by
uniformly selecting $(x,\xi)$ pairs and tossing out ones with $\xi > p(x)$.}
\label{MCRejection}
\end{center}
\end{figure}

Now, we choose a pair of random numbers $(x, \xi)$ with $x$ uniformly 
distributed over the interval
$v_{0} - M v_{\rm th} \leq x < v_{0} + M v_{\rm th}$ and $\xi$ uniform over 
$0 \leq \xi < p_{\rm max}$, where $p_{\rm max}$ is the maximum value of $p(v)$.
If $\xi < p(x)$, accept this $x$ as the sampled velocity. Otherwise, if 
$\xi > p(x)$, reject these $(x, \xi)$ and sample another pair. That the
resulting set of $x$ values is distributed as required can be seen by studying
the figure above. This figure also makes the case that the rejection approach
may not be very efficient for sampling a Maxwellian (indeed, DEGAS 2 
uses the well-known Box-Mueller method for this task).

In general, a rejection technique factors the probability density $p(x)$ into 
two functions, $p(x) = g(x) h(x)$, where it is known how to sample $x$ from
the function $g(x)$. The acceptance condition is just $\xi < h(x)$. The
Maxwellian example simply sets $g(x) = 1$.

\section{Particle Transport}

One of the main benefits of doing Monte Carlo particle transport is that 
most of the time the algorithm can be understood directly from a physical
picture. Namely, we think of the code as tracking particles through a
medium and having collisions with probabilities determined from the cross
sections of the various collision processes. The outcome of those collisions
is the same as for the actual physical particles, and so on. This is 
known as the \emph{analog} description of Monte Carlo particle transport.

In this section, we want to give a rough idea of how this picture relates
to the integral equations which provide the actual mathematical basis
for Monte Carlo techniques. Only from these equations can we understand
the {\em non-analog} techniques (the rejection technique described in
the previous subsection is a non-analog technique) used to make the Monte 
Carlo process more efficient. We hope to familiarize the reader with
the integral equations to the point where more detailed 
references\cite{Span69,Reit92} can be readily digested.

From an analog point-of-view, we imagine a Monte Carlo particle simulation
as consisting of four parts:
\begin{enumerate}
\item A way to sample an initial particle distribution,
\item A way to track particles from one point in space to another,
\item Tools to treat ``interesting things'' that might happen along the way,
\item Some useful quantities to compute and record at each stage of the 
process.
\end{enumerate}
The last of these, referring to ``scores'' and their ``estimators'' will
be addressed in Sec.~\ref{Estimators}.

\subsection{Motivating the Integral Equation}

The initial particle distribution, or source, is given by a probability 
density function $Q(x) > 0$. We take $x$ here to represent a point in 
phase space so that this function provides both the velocity and spatial
distribution of the source. 

The tracking of particles is described in the integral equation by a function
$L(y \rightarrow x)$ which gives the probability that a particle starting 
at a point $y$
in phase space ends up at point $x$. This function includes the ``interesting
things'' (collisions) 
that might happen along the way, as we will describe in the next 
subsection.

Using these two functions, we can write the {\em density of particles
arriving at $x$}, $\chi(x)$ as the sum of two pieces:
\begin{enumerate}
  \item Arrivals directly from source: $Q(x)$.
  \item Particles transported from elsewhere: 
\begin{equation}
\int \chi(y) L(y \rightarrow x)  \: dy
\end{equation}
   \end{enumerate}
And we have
\begin{equation}
\chi(x) = Q(x) + \int \chi(y) L(y \rightarrow x)  \: dy.
\end{equation}

\subsection{Transport and Collision Kernels}

In order to describe the kernel $L(y \rightarrow x)$ in more detail, substitute
for $y$ and $x$ the phase space variables $(\vec{r},\vec{v})$ and a species
index $i$. The kernel $L$ can then be broken into two pieces:
\begin{equation}
L(\vec{r}^{\prime}, \vec{v}^{\prime}, i^{\prime} \rightarrow \vec{r}, 
\vec{v}, i) = T(\vec{r}^{\prime} \rightarrow \vec{r}; \vec{v}, i) C(\vec{v}^{\prime}, 
i^{\prime} \rightarrow \vec{v}, i; \vec{r}),
\end{equation}
where $T(\vec{r}^{\prime} \rightarrow \vec{r}; \vec{v}, i)$ 
is the (probability) density of particles leaving a 
collision at $\vec{r}^{\prime}$ and having a subsequent collision at $\vec{r}$. 
The ``;'' is used in
the arguments of $T$ and $C$ to separate changing parameters (on the left) from
those held fixed (on the right). Basically, sampling from $T$ is as 
described in
Section~\ref{SamplingDistance}.

Collisions are described by $C(\vec{v}^{\prime}, i^{\prime} \rightarrow 
\vec{v}, i; \vec{r})$ which provides the probability
that a particle of velocity $\vec{v}^{\prime}$, species $i^{\prime}$ 
will have a collision resulting in 
a particle of velocity $\vec{v}$, species $i$. In practice, $C$ is a 
sum over the physical
reactions in which species $i^{\prime}$ participates.

\subsection{Integral Equation Details}
\label{IntegralEquationDetails}

One of the problems with reading Monte Carlo particle references is that
there are several slightly different ``quantities of interest'' equivalent
to $\chi(x)$. Furthermore, there is no well-established notation for those
quantities. In this subsection, we attempt to list these various quantities
and to provide a detailed form of an integral equation for later reference.

The quantity $\chi(x)$ is called the {\em post-collision density}. The 
origin of this designation should be clear from the previous discussion.
Equivalently, the integral equation can be written in terms of the 
{\em pre-collision density} $\psi(x)$. The two are related by
\begin{equation}
\psi(\vec{r},\vec{v},i) = \int d^{3}r^{\prime} 
\chi(\vec{r}^{\prime},\vec{v},i) 
T(\vec{r}^{\prime} \rightarrow \vec{r}; \vec{v},i).
\label{post2pre}
\end{equation}

It is this pre-collision density which is most closely related to the 
more familiar particle flux $\Phi$ and particle
distribution function $f$ by the definitions:
\begin{equation}
\psi(\vec{r},\vec{v},i) = \Sigma_{t}(\vec{r},\vec{v},i) \Phi(\vec{r},\vec{v},i)
 = \Sigma_{t}(\vec{r},\vec{v},i) |\vec{v}| f(\vec{r},\vec{v},i).
\end{equation}

The integral equation for $\psi$ bears a strong resemblance to that for
$\chi$. We write it out here in detail to illustrate the details as well as 
to provide a reference point for the rest of the document and other
reading:
\begin{eqnarray}
\psi(\vec{r},\vec{v},i) & = & S(\vec{r},\vec{v},i) \nonumber \\
& & \quad + \int_{0}^{\infty} dR \:
\int d^{3}v \: 
C(\vec{v^{\prime}},i^{\prime} \rightarrow \vec{v},i; \vec{r} - \vec{\Omega} R) 
\nonumber \\
& & \qquad \quad \times T(\vec{r} - \vec{\Omega} R \rightarrow \vec{r}; \vec{v},i)
\psi(\vec{r} - \vec{\Omega} R, \vec{v^{\prime}},i^{\prime}),
\end{eqnarray}
where
\begin{equation}
S(\vec{r},\vec{v},i) = \int_{0}^{\infty} dR \: 
Q(\vec{r} - \vec{\Omega} R, \vec{v}, i)
T(\vec{r} - \vec{\Omega} R \rightarrow \vec{r}; \vec{v},i),
\end{equation}
and
\begin{equation}
T(\vec{r} - \vec{\Omega} R \rightarrow \vec{r}; \vec{v},i)
= \Sigma_{t}(\vec{r},\vec{v},i) \exp \left[ - \int_{0}^{R} dR^{\prime} \:
\Sigma_{t}(\vec{r} - \vec{\Omega} R^{\prime},\vec{v},i) \right].
\end{equation}
The actual form of the collision kernel $C$ need not be written down
here. Instead, in this document we'll only talk about its implementation
from an analog point of view. The only additional piece of information
about it that we could add to the above equations is that $C$ consists
of a sum over the relevant collision processes involving species $i^{\prime}$.

The approach of writing the path integrals as above instead of as general
phase-space volume integrals with the path enforced via 
delta-functions\cite{Reit92}
was copied from Case and Zweifel\cite{Case67}
We find that this leads to expressions which are easier to work
with and think about.

\section{Estimators}
\label{Estimators}

The objective of the Monte Carlo calculation is to evaluate integrals over the
distribution function:
\begin{equation}
I(g) \equiv \int d^{3}r d^{3}v f(\vec{r},\vec{v})
g(\vec{r},\vec{v}).
\label{gestimator}
\end{equation}
For example, if we wish to estimate a reaction rate 
$\rho \equiv \Sigma |\vec{v}|$,
\begin{equation}
I(\rho) = \int d^{3}r d^{3}v \Sigma \Phi(\vec{r},\vec{v}).
\end{equation}

Values of $\rho$ are accumulated along the random walk of the neutral flight
and used to estimate $I(\rho)$. The mathematical description of this
accumulation process is called the {\em estimator} of $I$. 

\subsection{Simple Estimators}
\label{SimpleEstimators}

Consider a random walk $\gamma$ which consists of $k$ steps,
$(x_{1}, x_{2}, \ldots, x_{k})$. I.e., the
particle is absorbed or otherwise terminated at the $k$th step.
The simplest estimator involves writing down the value of $\rho$
at the point of absorption. Thus, the absorption estimator is
\begin{equation}
\xi_{a}(\gamma) = \frac{\rho}{\rho_{a}} w \chi_{V},
\end{equation}
where $\rho_{a}$ is the absorption rate in the volume of interest
$V$ and $w$ is the statistical weight of the particle (constant for
this random walk process). The function $\chi_{V}$
\begin{equation}
\chi_{V} (x_{m}) = \left \{ \begin{array}{ll}
                        1 & \mbox{if $x_{m} \in V$} \\
                        0 & \mbox{if $x_{m} \not\in V$}
                       \end{array}
                         \right.,
\end{equation}
just permits the score to be localized to the volume of interest.

The expectation value of the estimator is just the sum over all
possible random walks $\gamma$ of the probability of each
walk $P(\gamma)$ times the corresponding estimator,
\begin{equation}
E(\xi_{a}) = \sum_{\gamma} P(\gamma) \xi_{a}(\gamma).
\end{equation}
Obviously, only one data point is collected for each random walk.

We can instead accumulate a score at each step (scattering collision)
of the random walk.
\begin{equation}
\xi_{c1}(\gamma) = \sum_{m=1}^{k} \left[ 
\frac{ \rho_{a} (x_{m}) }{ \rho_{t} (x_{m}) } \right] w
\chi_{V} (x_{m}), 
\end{equation}
where $\rho_{t} = \rho_{a} + \rho_{s}$ is the total reaction
rate for the particle, with $\rho_{s}$ being the scattering
rate. The subscript $c1$ denotes the first collision estimator
described here. The expectation value of the estimator is again
formally written as
\begin{equation}
E(\xi_{c1}) = \sum_{\gamma} P(\gamma) \xi_{c1}(\gamma).
\end{equation}
Because of the greater amount of data accumulated along each particle
track, the variance of this estimator will usually be smaller
than that of the absorption estimator.

Consider now a random walk process in which absorption is forbidden.
This is our first departure from the usual analog processes.
Instead of an absorption, the weight of the particle is reduced
at each scattering collision by its survival probability.
So, at the $m$th step of the walk the weight of the particle is
\begin{equation}
w_{m} = w  \prod_{j=1}^{m} \frac{ \rho_{s} (x_{j}) }{ \rho_{t} (x_{j}) }.
\end{equation}
The flight must be terminated in some manner; the usual practice
is to have it undergo Russian roulette\cite{Span69} when 
$w_{m} < w_{\rm min}$, where $w_{\rm min}$ is some specified 
minimum weight. The collision estimator for this nonanalog
random walk process is then
\begin{equation}
\xi_{c2}(\gamma) = \sum_{m=1}^{k} w_{m} \left[ 
\frac{ \rho_{a} (x_{m}) }{ \rho_{t} (x_{m}) } \right]
\chi_{V} (x_{m}).
\label{nonanalog}
\end{equation}

For formal proofs that these expectation values do indeed yield $I(\rho)$,
see Spanier and Gelbard\cite{Span69}.

\subsection{Generalized Collision and Tracklength Estimators}
\label{generalest}

Two principal
estimators are used in DEGAS 2. One is a generalization of the 
collision estimator described in the previous subsection; the other
is the tracklength estimator.

These estimators can be obtained by making a slight modification of the
derivation given by Macmillan\cite{Macm67} (who had improved upon an
earlier approach described by Spanier\cite{Span66}). To make the connection
to DEGAS 2 clearer, we replace the $\Sigma$ and $d$ variables of Macmillan
with $\rho$ and $t$.

Macmillan's paper, like the earlier one by Spanier, aims to derive the 
tracklength estimator. Collisions are described as either
scatterings, rate $\rho_{s}$, or absorptions $\rho_{a}$. The 
starting point is the nonanalog random walk process described in
Sec.~\ref{SimpleEstimators}. The probability of an event occurring in
a time interval $dt$ is $(\rho_{s} + \rho_{a}) \, dt$. Again, the weight
reduction factor applied at each scattering 
is $\rho_{s} / (\rho_{s} + \rho_{a})$.

The trick used by Spanier is to introduce a fictitious ``delta scattering''
that has no impact on the particle's trajectory or weight. These
``pseudocollisions'' are not unrelated to those of the original
DEGAS code\cite{Heif82}. Additional flexibility in the final results is 
obtained by associating some fraction of the absorptions $\alpha$ with
the delta scatterings. To obtain the estimators used in DEGAS 2, we
work with a total fictitious scattering reaction rate
\begin{equation}
\rho_{t,\delta} \equiv \rho_{s,\delta} + \alpha \rho_{a},
\end{equation}
where $\rho_{s,\delta}$ is equivalent to Macmillan's $\Sigma_{s,\delta}$.
Then, the probability of a delta scattering in a time interval $dt$ is
$\rho_{t,\delta} \, dt$. At each delta scattering, the weight is
reduced by a factor $(\rho_{t,\delta} - \alpha \rho_{a}) / \rho_{t,\delta}$.
The probability of a real scattering in the same interval is
$(\rho_{s} + (1 - \alpha) \rho_{a}) \, dt$; at each such event, the
weight is reduced by $\rho_{s} / [\rho_{s} + (1 - \alpha) \rho_{a}]$.

If a particle travels a time $t$ in a region of constant parameters
before having a real collision, the conditional probability of having
$n$ pseudocollisions is
\begin{equation}
P(n \mid d) = \frac{ ( \rho_{t,\delta} t)^{n}}{n!} \exp(- \rho_{t,\delta} d).
\end{equation}
At each pseudo or real collision, we use the standard collision estimator
for the nonanalog random walk, taken from Eq.~(\ref{nonanalog}):
\begin{equation}
\xi = \frac{\rho}{\rho_{t}^{\ast}} w,
\end{equation}
where $w$ is the current weight of the particle and
\begin{equation}
\rho_{t}^{\ast} = \rho_{t,\delta} + \rho_{s} + (1 - \alpha) \rho_{a}
\end{equation}
is the overall total reaction rate.

The expected value of this estimator over the $n$ collisions in the
region is then
\begin{equation}
E(\xi \mid t) = \sum_{n=0}^{\infty} P(n \mid t) \frac{\rho}{\rho_{t}^{\ast}}
\sum_{i=0}^{n} \left( 
\frac{\rho_{t,\delta} - \alpha \rho_{a}}{\rho_{t,\delta}} \right)^{i},
\end{equation}
where we have assumed that the initial weight of the particle is 1. The
last factor represents the subsequent weight reductions with each
pseudocollision.

Following Macmillan's evaluation of $E(\xi \mid t)$, we find that we
can write
\begin{equation}
E(\xi \mid t) = c^{\prime} \frac{\rho}{\rho_{s} + (1 - \alpha) \rho_{a}}
+ (1 - c^{\prime}) \rho \frac{1 - \exp(- \alpha \rho_{a} t)}{\alpha \rho_{a}},
\end{equation}
where
\begin{equation}
c^{\prime} \equiv 
\frac{\rho_{s} + (1 - \alpha) \rho_{a}}{\rho_{t,\delta} + \rho_{s} 
+ (1 - \alpha) \rho_{a}}
\end{equation}
should be compared with Macmillan's $c$.

Macmillan's expression for the weight reduction after time $t$
is exactly the same in our case:
\begin{equation}
E(w \mid t) = \exp( - \alpha \rho_{a} t) \frac{\rho_{s}}{\rho_{s} + (1 - \alpha) \rho_{a}}.
\end{equation}
The factor at the end represents the weight reduction by the real collision
at time $t$. If the particle travels a time $t$ through the region without
having a real collision, its weight should be reduced by just
$\exp( - \alpha \rho_{a} t)$.

In this latter situation in which there is not a real collision at time
$t$, only the pseudocollisions contribute to the score. An interesting
difference from MacMillan's approach is that if we consider the limit
$\rho_{t,\delta} \rightarrow 0$, we should get no score at all. In
contrast, MacMillan's pure ``collision estimator'' still compiles
a nonzero score as $\Sigma_{s,\delta} \rightarrow 0$ through the
$\alpha \Sigma_{a}$ contribution to the pseudocollision cross section.
Thus, MacMillan's ``collision estimator'' behaves somewhat like
a track-length estimator if there are no collisions in a region.
For simple scores, this behavior is desirable: lower variances will
result in regions of small cross section. However, a few scores are
sufficiently complex that averaging along the particle path cannot
be done, and scoring can only be done at collisions. A pure 
collision estimator is required for handling such scores.

In the case of no real collision at time $t$, the expectation value
for our estimator becomes:
\begin{equation}
E(\xi \mid t) = \frac{\rho}{\rho_{t}^{\ast}} \rho_{t,\delta}
\frac{1 - \exp(- \alpha \rho_{a} t)}{\alpha \rho_{a}}.
\label{nocollision}
\end{equation}
Clearly, $E(\xi \mid t) \rightarrow 0$ as $\rho_{t,\delta} \rightarrow 0$.
The limiting expressions above for $\rho_{t,\delta} \rightarrow 0$
can be obtained from MacMillan's with $\Sigma_{s,\delta} \rightarrow
- \alpha \Sigma_{a}$.

We now consider the expressions used in DEGAS 2. We start by using
$\alpha = 1$ so that absorptions occur only at the delta scatterings.
This is consistent with our use of the nonanalog random walk, i.e.,
suppressed absorption. The collision estimator is then given by
$\rho_{t,\delta} = 0$,
\begin{equation}
\xi_{c} \equiv \frac{\rho}{\rho_{s}} \exp(- \rho_{a} t) w \chi_{V},
\label{ce}
\end{equation}
where $w$ is the weight of the particle at the beginning of timestep
$t$ and we have inserted the characteristic function of $V$ $\chi_{V}$ 
for completeness.
At the end of the timestep (after the collision is scored) the weight
is reduced by
\begin{equation}
w^{\prime} = w \exp( - \rho_{a} t).
\label{weightfactor}
\end{equation}
The particle track continues from there, and multiple scores may be made
inside the volume $V$. If the flight travels through the volume $V$ without
a collision, no contribution to the score is made. However, the
weight is still reduced by the factor in Eq.~(\ref{weightfactor}).

The track-length estimator is obtained with $\rho_{t,\delta} \rightarrow
\infty$,
\begin{equation}
\xi_{t} \equiv \rho \frac{1 - \exp(- \rho_{a} t)}{\rho_{a}} w \chi_{V}.
\label{tle}
\end{equation}
Note that Eq.~(\ref{nocollision}) yields this same result in the case
of no collisions. The same weight reduction, Eq.~(\ref{weightfactor})
is applied at the end of $t$. Conveniently, this allows us to use the
same weight factor for both estimators (true in MacMillan's case as well).

Molecules and some other species are not treated with suppressed absorption.
Instead, an ionizing collision terminates the flight (typically changing
the species). The estimators for this case are obtained by setting
$\alpha = 0$ in the above expressions:
\begin{equation}
\xi_{c} \equiv \frac{\rho}{\rho_{s} + \rho_{a}} w \chi_{V},
\end{equation}
and
\begin{equation}
\xi_{t} \equiv \rho t w \chi_{V}.
\end{equation}
There is no weight reduction in this case, and $w$ retains its initial
value for the duration of the flight.

\subsection{Specification of Tallies}
\label{tallies}

This section briefly covers the practical implementation of the estimators
described in the previous subsection.
DEGAS 2 currently features three primary types of estimators. These estimators
are used in accumulating three types of tallies. The tally types are
\begin{enumerate}
  \item test,
  \item reaction,
  \item sector.
\end{enumerate}
These designations can also be viewed as characterizing the nature of the
information contained in the function $g$ of Eq.~(\ref{gestimator}).
The following discussion of the three estimators applies to the first
two of these; there is only one estimator for the sector tallies.
 
The first estimator, associated with the accumulating
test particle information, is just the track length estimator (see also
Sec.~\ref{generalest}),
\begin{equation}
\xi_{g}^{\rm TLE} = w \, \frac{1 - \exp(-\sigma_{i} t_{V})}{\sigma_{i}} g,
\end{equation}
where  $t_{V}$ is the time the particle track spent in volume V,
$w =$ weight at start of the time step. Compare this with Eq.~(\ref{tle}).
The function $g$ can be a test particle attribute, such as
mass, momentum, energy, or reaction-related, e.g., ion 
momentum source due to charge exchange. 
Volumetric
plasma sources, such as the latter, 
enter the track-length estimator as averages over the background 
distribution\cite{Reit93}.

The second type of estimator is the collision estimator,
\begin{equation}
\xi_{g}^{\rm CE} = \sum_{m_{s}} w_{m_{s}} 
\frac{g}{\sigma_{s}},
\end{equation}
where the sum is over $m_{s}$, all scattering collisions within $V$.
The quantity $w_{m_{s}}$ is the weight at $m_{s}$th collision (by using the
instantaneous weight rather than the initial weight, 
the exponential weight reduction factor that appears explicitly in
Eq.~(\ref{ce}) is incorporated automatically). 
Again, $g$ is a test particle attribute.

There is an analogous expression for the collision estimator associated
with a particular reaction,
\begin{equation}
\xi_{g}^{\rm CE} = \sum_{m_{j}} w_{m_{j}} \frac{g}{\sigma_{j}},
\end{equation}
where the sum is now over $m_{j}$ collisions of reaction $j$ within $V$,
The function $g$ is some quantity related to the reaction and is likely
something known only at collisions. One example of such a quantity would
be the energy source due to H$_{2}$ dissociation (i.e., the energy
exchange is determined only once the product velocities have been 
determined).

The third type of estimator arises from the observation that 
many integrals of interest can be written as:
\begin{equation}
I(g) \equiv \int d^{3}r g(\vec{r}) \int d^{3}v f(\vec{r},\vec{v}) = 
\int d^{3} r g(\vec{r}) n(\vec{r}).
\end{equation}
This leads to the post-processing estimator. The requirement is that 
$g$ not depend upon the test or background velocities. An example of
this is the H$_{\alpha}$ emission rate.

Information accumulated when a flight crosses a diagnostic sector (surface)
is compiled in the sector tallies. The
stimator is just $w g$, where
$w =$ is the weight when the track hits the sector.
Examples of the $g =$ functions include mass, incident angle, and energy.
This is effectively a collision estimator.

There is also a ``null'' or ``none''
estimator that can be used, say, to eliminate the
contribution of a particular reaction to a tally without having to 
remove that reaction from the problem.

For each tally, we need to specify
\begin{enumerate}
  \item A descriptive label,
  \item The type of tally (test, reaction, or sector),
  \item Its geometry, [volume (per zone), surface (diagnostic sector),
or detector (e.g., chord-integrated)],
  \item The dependent variable (e.g., energy, background momentum,
H$_{\alpha}$ emission rate). The requested variables 
are matched directly against the labels in the atomic physics data files.
If a new reaction and data file containing a new variable are added to
the code, that variable can be used as the basis for a corresponding tally
without having to modify the code.
  \item The rank (or dimensionality) of the tally,
  \item Its independent variables (1 $\rightarrow$ rank) [e.g.,
zone, test particle species, energy bin (diagnostic sector),
wavelength bin (detector).
  \item The estimator to be used (track-length, collision, post-processed,
or null). For reaction tallies, the estimator chosen can vary with
reaction. It may be advantageous to combine results 
from than one estimator, but no effort to implement such a capability
has been made yet.
  \item Any conversions that must be performed on the tally prior to
output.
Usually, these conversion will do some sort of scaling, e.g., by
mass or volume. An important example is 
converting $\vec{v}$ from Cartesian
to cylindrical coordinates when scoring.
\end{enumerate}

\section{Accumulation of Tallied Data}
\label{Accumulation}

The estimators described in Sec.~\ref{tallies} are calculated by
the routines in {\tt score.web}, which are invoked during tracking.
At the end of each flight, the resulting data are accumulated into the
global arrays described in the \href{classes.pdf#page.65}{statistics class} 
(package abbreviation {\tt sa}) via the subroutines in {\tt stat.web}.

The data accumulation process takes place over a range of scales
to optimize performance on massively parallel platforms.  The basic
breakdown is:
\begin{description}
  \item[flt] for ``flight''.  This always refers to a single flight.
  \item[frag] for ``fragment''.  To facilitate load balancing, 
only a fraction (or ``fragment'') 
of the total number of flights per processor
are sent at a time to each process.  
As each process completes, an additional fragment
is sent until all have been completed.
  \item[fin] for ``final''.  In the simplest case, this refers to
the global accumulation of data from all processes.  Again,
in the simplest case, the data are actually accumulated
directly in an array from the 
\href{classes.pdf#page.61}{output class} (see also
Sec.~\ref{outputformat}).
\end{description}

This process is complicated by two additional optimizations.
First, as is documented in the \href{classes.pdf#page.65}{statistics class},
these arrays may be compressed, especially at the {\tt flt} level,
to reduce the amount of data transmitted.  Second, in runs on
larger machines
(say, $> 100$ cores or processes) fragments are 
accumulated into intermediate ``final'' arrays over groups of
processors to reduce the number of accumulations that must
be performed by the master process, avoiding network 
transmission collisions. The machinery facilitating this is
contained in {\tt degasinit.web}.  

The sources of neutral particles to be simulated in a given
DEGAS 2 run are broken up into ``groups'' according to their
physical nature, e.g., ``plate recycling'', ``gas puff'', etc., 
as described in the \href{classes.pdf#page.26}{sources class}.  
Because the number of flights and total particle current
associated with each source group are specified independently
by the user, the data for each group are accumulated 
separately and are available by themselves within the
\href{classes.pdf#page.61}{output arrays}.

The fundamental expressions used in the data accumulation
process are, thus, for a specific source group $i$.  The 
other principal indexing parameters for the output arrays
are the geometric zones, $z$, and the tallies, Sec.~\ref{tallies}.
To minimize confusion, we do not introduce any explicit
index for the latter; the following expressions apply
independently for each tally.  The first accumulated moment,
to be associated with the mean in the end is:
\begin{equation}
S_{1}^{i}(z) = \sum_{j = 1}^{N_{i}} \rho_{j} \left (
\sum_{k \in z} \xi_{j,k} \right ) / \sum_{j = 1}^{N_{i}} \rho_{j},
\label{mean1}
\end{equation}
where: $i$ is the source group index, the ``1'' refers to the
first statistical moment, $j$ runs over all of the
$N_{i}$ flights in group $i$, $\rho_{j}$ is the relative
statistical weight of flight $j$, $k$ is the list of events 
scored
by flight $j$ (e.g., collisions)
in zone $z$, and $\xi_{j,k}$ is the estimator
value (as in Sec.~\ref{tallies}) 
for this tally recorded by flight $j$ during event $k$.
The statistical weight $\rho_{j} \equiv 1$ in most cases,
but can be set otherwise by the user to effect importance
sampling via the input arrays {\tt source\_segment\_rel\_wt}.
Note that this quantity is {\em fixed} for a given flight
(indeed, is the same for all flights launched from that
source segment) and is distinct from the statistical
weight, designated by $w$ in Secs.~\ref{SimpleEstimators},
\ref{generalest}, and \ref{tallies}, which may
vary during flight tracking, but {\em always} has 
an initial value of $w = 1$.

The equivalent expression for the second moment is:

\begin{equation}
S_{2}^{i}(z) = \sum_{j = 1}^{N_{i}} \rho_{j} \left (
\sum_{k \in z} \xi_{j,k} \right )^{2} / \sum_{j = 1}^{N_{i}} \rho_{j}
- \left [ \sum_{j = 1}^{N_{i}} \rho_{j} \left (
\sum_{k \in z} \xi_{j,k} \right ) / \sum_{j = 1}^{N_{i}} \rho_{j} \right ]^{2}.
\label{var1}
\end{equation}
While this expression is very close to that for the
variance appearing in statistics textbooks \cite{Bran2014}, it is not
suitable for the multi-level data accumulation used by
DEGAS 2.  
How one would break up the sum in the
numerator of Eq.~(\ref{mean1}) into smaller
sets (fragments) of
flights $j$ is clear.
To do likewise for the second moment, we need to instead
work with:
\begin{equation}
\sum_{j = 1}^{N_{i}} \rho_{j} \left (
\sum_{k \in z} \xi_{j,k} \right )^{2} / \sum_{j = 1}^{N_{i}} \rho_{j} 
= S_{2}^{i}(z) + 
\left [ \sum_{j = 1}^{N_{i}} \rho_{j} \left (
\sum_{k \in z} \xi_{j,k} \right ) / \sum_{j = 1}^{N_{i}} \rho_{j} \right ]^{2}.
\label{var2}
\end{equation}

We now motivate the specific expressions used to accumulate
data from an individual flight into a fragment, and a fragment
into a final total.
In the interest of clarity, 
we introduce some simplified notation:
$\sum_{j = 1}^{N_{i}} \rho_{j} \rightarrow w$,
$S_{1}^{i}(z) \rightarrow S_{1}$, and $S_{2}^{i}(z) \rightarrow S_{2}$.
The subscript $i$ in the following represents ``incremental'' data
(e.g., an individual flight), $b$ represents the ``base'' data
(e.g., the fragment containing data from earlier flights);
the result of the accumulation will be new base data, indicated
with a prime.
The familiar form for the first moment is:
\begin{equation}
S_{1b}^{\prime} = (w_{i} S_{1i} + w_{b} S_{1b}) / (w_{i} + w_{b}),
\label{mean2}
\end{equation}
with the new base weight being:
\begin{equation}
w_{b}^{\prime} = w_{i} + w_{b}.
\label{weight}
\end{equation}
But, the code actually uses an equivalent form to maximize
numerical accuracy:
\begin{equation}
S_{1b}^{\prime} = S_{1b} + (S_{1i} - S_{1b})\frac{w_{i}}{w_{i} + w_{b}}.
\end{equation}

The analogous expression following from Eq.~(\ref{var2}) is:
\begin{equation}
(S_{2b}^{\prime} + S_{1b}^{\prime 2} ) w_{b}^{\prime}
= (S_{2b} + S_{1b}^{2}) w_{b} 
+ (S_{2i} + S_{1i}^{2}) w_{i}.
\label{var3}
\end{equation}
Thus,
\begin{equation}
S_{2b}^{\prime} = [(S_{2b} + S_{1b}^{2}) w_{b} 
+ (S_{2i} + S_{1i}^{2}) w_{i} - S_{1b}^{\prime 2} w_{b}^{\prime} ] / w_{b}^{\prime}.
\label{var4}
\end{equation}
The equivalent form that is coded up in {\tt stat.web} is:
\begin{equation}
S_{2b}^{\prime} = [S_{2b} w_{b} + S_{2i} w_{i} 
+ (S_{1i} - S_{1b})(S_{1i} - S_{1b}^{\prime}) w_{i} ] / w_{b}^{\prime}.
\label{var5}
\end{equation}
The equivalence of Eqs.~(\ref{var4}) and (\ref{var5}) can be
demonstrated most readily by adding and subtracting 
$S_{1b}(S_{1i} - S_{1b}^{\prime}) w_{i}$ inside the square brackets of
Eq.~(\ref{var4}) and substituting in Eq.~(\ref{mean2}).  Again,
this form improves numerical accuracy, e.g., in the case
where $S_{1i} \simeq S_{1b}$.

The mean tally values for each source group are scaled by the
total current in that group, which is obtained by summing
over the group's $N_{{\rm seg},i}$ segments:
\begin{equation}
\Gamma_{i} = \sum_{\ell = 1}^{N_{\rm seg},i} \Gamma_{i,\ell}.
\end{equation}
In the steady state 
mode of operation, $\Gamma_{i}$ represents a number of 
particles per second; in the time dependent mode,
$\Gamma_{i}$ is a number of particles.  Obviously, these
quantities are used in combining the $S_{1}^{i}(z)$ and
$S_{2}^{i}(z)$ over source groups.  But, they are also
used to scale the user-specified relative statistical
weights to get the $\rho_{j}$ appearing in 
Eq.~(\ref{mean1}), etc. above and in setting the
probabilities for sampling flights from the source
segments. 

Namely, the user can specify
an arbitrary 
set of relative weights, $r_{i,\ell}$, for segments $\ell =
1 \rightarrow N_{{\rm seg}, i}$ in group $i$.  
A normalization factor is computed for each source group:
\begin{equation}
r_{{\rm norm}, i} = \left ( \sum_{\ell = 1}^{N_{{\rm seg}, i}} 
\Gamma_{i,\ell} / r_{i,\ell} \right ) / 
\left ( \sum_{\ell = 1}^{N_{{\rm seg}, i}} 
\Gamma_{i,\ell} \right ).
\label{wtnorm}
\end{equation}
The cumulative probability distribution used in sampling flights
in group $i$ is, e.g., for segment $\ell$:
\begin{equation}
p_{i, \ell} = \left ( \sum_{m=1}^{\ell -1} \Gamma_{i,m} / r_{i, \ell} \right )
/  \left ( \sum_{m=1}^{N_{{\rm seg},i}} \Gamma_{i,m} / r_{i, \ell} \right ).
\end{equation}
These probabilities are computed without the weight
normalization factor Eq.~(\ref{wtnorm}) since it
would appear in both numerator and denominator.
The actual sampling procedure is more complex
and offers additional flexibility; see
the \href{classes.pdf#page.26}{sources class} and
\href{sourcetest.pdf}{sourcetest}.

If flight $j$ in group $i$ is sampled from 
segment $i$, the relative weight
is set to:
\begin{equation}
\rho_{j} = r_{i, \ell} \times r_{{\rm norm}, i}.
\end{equation}

The final estimated value for a given tally in 
zone $z$ is, for each individual group:
\begin{equation}
E^{i}(z) = \Gamma_{i} S_{1}^{i}(z)
\end{equation}
In a time dependent run, this is divided by
the simulated time interval $\Delta t = t_{f} - t_{i}$
so that the dimensions are the same as in a steady
state run.  This is then summed over all 
groups to get the total:
\begin{equation}
E(z) = \sum_{i=1}^{N_{\rm groups}} E^{i}(z).
\end{equation}

The output quantity derived from the second
moment $S_{2}$ is actually the
relative standard deviation (see also
Sec.~\ref{statcompare}).  Namely, the expected
error in a statistical estimate of the mean $\mu$
via $N$ samples is $\sigma / \sqrt{N}$, where
$\sigma$ is the standard deviation and $\sigma^{2}$ 
is the variance.  
The relative standard deviation ``rsd'' is then
\begin{equation}
{\rm rsd} = \frac{\sigma}{\mu \sqrt{N}}.
\end{equation}

Note that Eq.~(\ref{var1}) and so on have 
$N \leftrightarrow \sum \rho_{j}$ in the denominator
of the variance even though $N-1$
appears in most textbook expressions (e.g., \cite{Bran2014}).
The ``-1'' accounts for the fact that the 
same sample is used to estimate the mean $\mu$ and the variance.
While $N \gg 1$ in most DEGAS 2 runs, rendering the difference
between $N$ and $N-1$ negligible, retaining $N-1$
yields a maximum relative standard deviation of exactly
unity when $\rho_{j} \equiv 1$ (a single score), facilitating interpretation of 
results.  In the process of computing the relative
standard deviation, we multiply the 
$S_{2}$ by $N/(N-1)$, and let the $N$ in this numerator
be cancelled the $N$ in the expression for the ``rsd''.

The ``rsd'' values required depend on the problem at hand
and on location within the problem space.  Some rough guidelines adapted
from the manual for the MCNP \cite{MCNP}  code can serve as a 
useful starting point:

\begin{center}
\begin{tabular}{cl}
rsd Range  &  Tally Quality  \\ \hline
$0.5 \rightarrow 1$ & Garbage  \\
$0.2 \rightarrow 0.5$ & Factor of a few  \\
$0.1 \rightarrow 0.2$ & Questionable \\
$< 0.1$ & generally reliable \\ 
\end{tabular}
\end{center}

Iterative problems, such as with neutral-neutral scattering
or coupling to a plasma code, demand 
precision, at least in high interaction regions,
greater than that desired in the overall solution.

\section{Tracking Procedure}
\label{TrackingProcedure}

The two previous subsections are tied together with an outline of the
DEGAS 2 subroutine for following flights in Fig.~\ref{flightflowchart}.

\begin{figure}
\begin{center}
\leavevmode
\includegraphics[width=6in]{flight_flowchart}
\vspace{-1.85in}
\caption{Simplified flow chart of DEGAS 2 subroutine {\tt follow}, showing the steps 
involved in tracking a flight from its birth to its termination. Note
the weight reduction, and collision and tracklength estimators 
(Sec.~\ref{Estimators}). The operators from the integral equations are
shown along with the pre- and post-collision densities. These labels 
tie the practical algorithmic operations to the equations in 
Sec.~\ref{IntegralEquationDetails}. Some details have been omitted
for clarity.}
\label{flightflowchart}
\end{center}
\end{figure}

The flight is initiated as it is sampled from the spatial and velocity 
distribution of a source. Conceptually, this step amounts to
specifying $\chi_{0}(\vec{r}^{\prime},\vec{v}^{\prime},i)$ using the
source term $Q(\vec{r}^{\prime},\vec{v}^{\prime})$. The subscript
$0$ denotes the initial track for the flight. 

The next stage of the algorithm involves computing the reaction
rates $\rho_{s}$, $\rho_{a}$, etc. for the flight. Continuing
flights which require a recomputation (due to crossing from one
plasma region to another or to a change in species or velocity)
repeat the main loop starting at this same point. The 
post-collision density $\chi_{n}$
in Fig.~\ref{flightflowchart} is intended to refer to them.

The time to the next collision is randomly sampled 
(Sec.~\ref{SamplingDistance}) from the sum of the non-absorbing
reaction rates using a random number $\zeta$.

The weight $w$ is compared with some minimum weight to guarantee
that the flight terminates at some point. The ``End flight''
process indicated in Fig.~\ref{flightflowchart} actually represents
a Russian Roulette procedure\cite{Span69} which may allow the flight
to live a little longer.

The transport kernel $T$ translates the flight from position
$\vec{r}^{\prime}$ to $\vec{r}$. At the same time, this step 
represents the conceptual transition from post-collision density
to pre-collision density $\psi$, Eq.~(\ref{post2pre}). The tracking 
algorithm stops when $t = t_{\rm max}$, the flight
crosses a zone boundary (denoting a change in plasma parameters), or
a ``sector'' (a material or diagnostic surface; these will be
discussed in Sec.~\ref{Geometry}) is reached. If a zone boundary
is crossed, the flight continues as is, although the rates must
be recomputed. Processing a sector may result in the continuation or the
termination of the flight, depending on the type of sector. Again,
the details have been omitted for simplicity.

If the flight reaches $t = t_{\rm max}$, the implication is that
it has had a collision. A particular collision is randomly chosen from 
amongst those possible for this particular flight.
At this point, precollision scores are recorded
and the collision itself is processed. Conceptually, this represents the
transition via the operator $C$ from $\psi_{n+1}$ to $\chi_{n+1}$.
If the reaction chosen results in no particles to track (e.g., 
in no neutral species; particles which DEGAS 2 has been instructed
to track are called ``test'' particles), the flight is terminated.
Test particles resulting from the collision are tracked in the same
manner, starting with the computation of the reaction rates.

\section{Random Numbers}
\label{Randomnumbers}

The concepts behind the random number generator used in DEGAS 2 are
described in the file {\tt random.web}. The 
\href{random.pdf#page.1}{documentation} at the top of the file
provides a nice introduction.

\section{Geometry}
\label{Geometry}

% Note: mostly copied this from the introduction to definegeometry2d.web;
% that was derived from readgeometry.web.
The hiearchy of geometry-related objects in
DEGAS 2 is, from the lowest level to the highest level, is:
\begin{enumerate}
  \item surface,
  \item cell,
  \item polygon,
  \item zone.
\end{enumerate}
Individual cells are composed of surfaces, as discussed in the documentation
for the internal geometry (e.g., see {\tt geometry.web}; follow
this \href{geometry.pdf#page.1}{link} to the PDF version of this file). 
The next level up
in two (and three, to a limited extent) dimensions
is a polygon. 
Finally, a zone may consist of one or more polygons; properties
are constant across a zone. For example, a single zone might be used to
represent the vacuum region around the plasma which is comprised of several
(possibly disconnected) polygons. Or, a plasma flux surface on which density
and temperature are constant might be a single zone. 
Note that polygons are used only for external interfaces to the
geometry such as that provided by {\tt definegeometry2d} or
{\tt readgeometry} and are not retained in the classes of the
main code.  Cells and surfaces
are essentially used only by the code itself; the user would have to deal
with them only in debugging. Zones are used primarily for the
specification of input and output data.  For more detailed
background on the DEGAS 2 geometry and some guidance on how to
set one up for your particular problem, see the introductions
to the external geomtry interface codes {\tt definegeometry}
(PDF file \href{definegeometry2d.pdf#page.1}{here}),
{\tt readgeometry} 
(PDF file \href{readgeometry.pdf#page.1}{here}), or
{\tt boxgen} (less well documented).

Another important component of the geometry are the ``sectors'';
they are defined at interfaces between adjacent zones of
different types to facilitate tracking or between zones of the same
type for diagnostic purposes.
Namely, flights are halted at sectors during tracking. If the next
zone along the track represents, say, a solid material, the
main code hands that flight off to the routines that
effect the interaction of the current particle with that material. 
While stopped at that sector,
scores to corresponding ``diagnostics'' (groups of sectors)
are also tabulated. The essential identifying information 
of a sector are its surface and zone numbers.
See the 
\href{classes.pdf#page.39}{sector class} documentation for a more
detailed description.  Default diagnostic sector groups defined
by all DEGAS 2 geometry setup programs are discussed in 
Sec.~\ref{Diagnostic}.

\section{Symmetry and Coordinate Systems}
\label{Symmetry}

The coordinate system in the main DEGAS 2 code is purely 3-D
Cartesian; that's essential to an efficient tracking algorithm.  There
are, however, two coordinate system related concepts.

\begin{enumerate}
  \item The geometry's symmetry. In most problems, the user's objectives
can be met with an approximate 2-D solution possessing an ignorable
coordinate. The two situations considered in DEGAS 2 are invariance
to translation along the $y$ axis (``plane'' symmetry) and
invariance to rotation about the $z$ axis (``cylindrical'' symmetry).
In both cases, the tracking is still done in 3-D, but the output
results are averaged over the third dimension.  As will
be discussed in more detail below, 1-D and 3-D ``symmetries''
are also permitted. Most geometry related variables refer to 
Cartesian coordinates.  However, in cylindrically symmetric or nearly symmetric
(3-D) cases, some geometry related variables are in cylindrical
coordinates.
  \item Coordinate system of input and output velocity fields.
The main need for this is in cylindrically symmetric or
nearly symmetric
problems, e.g., when DEGAS 2 is coupled to a 2-D
plasma code like UEDGE.  In these cases, the plasma flow velocities
are effectively specified at a toroidal angle $\phi = 0$, corresponding
to $y = 0$.
To use these velocities in a collision at some particular location
for which $\phi \neq 0$ (equivalently, $y \neq 0$), 
the velocities
need to be rotated through that angle $\phi$.
Likewise, a code like UEDGE expects to get from DEGAS 2 
toroidally symmetric ion momentum
sources specified at $\phi = 0$. So, the scores for those sources
need to be rotated from the local toroidal angle $\phi \neq 0$ 
to $\phi = 0$ as they are
being compiled.  These two transformations are handled by macros
in the \href{classes.pdf#page.23}{background class}.
In problems with
plane or no symmetry, these macros have no effect.
\end{enumerate}

The currently permissible symmetry types (see 
also {\tt geometry.hweb})
are:
\begin{description}
  \item[None] Currently indicates that the symmetry has not been defined,
{\em not} that the geometry has no symmetry.
  \item[Plane] A 2-D geometry in which the $y$ coordinate is ignored.
  \item[Cylindrical] A 2-D geometry in which the toroidal angle 
$\phi$ is ignored.
  \item[1-D] A 1-D (plane) geometry in which the $y$ and $z$ 
coordinates are ignored.
  \item[3-D Plane] This geometry is close to be planar symmetric;
all of the geometrical objects are defined using 2-D shapes
defined in a single plane. The problem is bounded in the $y$
direction, and the $y$ variation of the results can be 
computed and output. The nature (solid vs. plasma or vacuum)
of the objects can also vary in
the $y$ direction.  The problem is divided in the $y$ direction
by $y = $ constant surfaces.  This might be described as a
``bread slice'' approach to a 3-D geometry.
  \item[3-D Cylinder] This is the cylindrical equivalent of the
3-D Plane symmetry. The full range of toroidal angles (360 degrees
or $2 \pi$ radians) is included in this case. Obviously, the
problem is divided in the $\phi$ direction by $\phi =$ constant
surfaces. The analogous simplified characterization of this approach
is the ``pie slice'' technique.
  \item[3-D Cylindrical Section] This is also a cylindrical equivalent
of the 3-D plane symmetry.  The difference is that in this case
only a finite range of toroidal angles is considered.  The problem
is bounded by $\phi = $ constant surfaces at the extrema of this
range.  Extending the ``pie slice'' analogy, this case represents
some number of slices, but less than the whole pie!
\end{description}

\section{Time Dependence}
\label{TimeDep}
\index{time dependence}
\index{snapshot}
\index{defineback}

The default and most frequently used mode of running DEGAS 2 simulates a
steady state, time independent system.  The reason for this is physically
obvious: unconstrained by magnetic field lines, neutral particles can
travel through a plasma much more quickly than charged species, and
in many cases, the neutral profile equilibrates on a time scale
short compared with global transport times.
For a rapidly evolving system, however, a
time dependent calculation may be needed. The implementation
of time dependence in DEGAS 2 follows
that used in EIRENE \cite{Reit95}.

Only the {\tt defineback} input preprocessor is presently set up
to handle direct user specification of a time dependent run;
see the corresponding \href{defineback.pdf#page.7}{documentation} 
for details on how this is done.
However, the user can in principle also manually
set the controlling
parameters in the \href{classes.pdf#page.26}{source class}.

The user specified sources in a time dependent run are
either fixed in time during a time 
interval (uniform) or start at the beginning of the interval 
($\delta$ function).  More complex time variation of the sources
can be created by breaking up the source time dependence into a
set of discrete time intervals with the sources being uniform
over each.

To provide continuity between consecutive time intervals, an
additional ``snapshot'' source group is set up to sample
from the neutral particle distribution function (PDF) 
that was in the volume at
the end of the end of a previous interval.  
Again following \cite{Reit95}, we represent this PDF via the
particle positions,
velocities, and weights; this is the purpose of the
\href{classes.pdf#page.127}{snapshot class}.  These
data are written to the file indicated by the
\verb+snapshotfile+ in {\tt degas2.in} (Sec.~\ref{degas2in}).
With this
approach, the amount of kinetic detail in the PDF is
maximized.

If the user is truly solving an initial value problem
and wishes to track the evolution of the neutral density
in time from a null initial state, e.g., the
temporal evolution of a gas puff from the start, no
initial snapshot file is required.  A more typical 
application, however, is to follow the common evolution
of the neutral and plasma behavior at the same time
from some nominal initial state.  To do that, we
need a 
snapshot distribution representative of that
initial state.

To generate such a distribution, the code is run in a
special mode that is similar to a standard, steady
state simulation in that the plasma background
is fixed and the flights are tracked
until they are ionized or exit the simulation volume.
However, the path of each flight is broken up into
a set of uniform time steps of the same size as those
used in the actual time dependent calculation.  
The start time of a new flight coming
from the gas source (e.g., a gas puff) is uniformly sampled
from the initial time interval. The flight is tracked through the
remainder of that interval; at the end, its characteristics are
deposited in the snapshot distribution. The flight
is then tracked forward from that point through another time interval,
resulting in an additional entry to the snapshot distribution. The
process continues in this manner until the flight is effectively
ionized or leaves the volume. The end result is a snapshot
distribution that fills the volume and is statistically equivalent to
a steady state calculation.  The number of snapshot particles
generated per flight varies with the problem and is,
actually, a measure of the average neutral density in the
volume.  Consequently, the user should perform a few scoping runs
with a limited number of flights to estimate this ratio
and determine the number of flights required to yield
the desired snapshot population.

Sampling from a large snapshot distribution in a
parallel run of {\tt flighttest} can
slow the initial broadcast of the source information
to the slave processes and stress memory resources on
those cores.  The {\tt MASTER\_SAMPLE} switch in
the {\tt Makefile.local} file can be set to
``yes'' to allow sampling to be done instead on
the master process so that only the flight
data need be broadcast to the slaves.
For additional information, see
the documentation of the associated
\href{classes.pdf#page.133}{flight\_frag class}.


\section{Atomic Physics}
\label{AtomicPhysics}
\index{Atomic physics}
\index{collisional-radiative}
\index{ionization}
\index{recombination}
\index{rates, atomic physics}

The interaction of electrons with atoms and ions in a divertor
plasma is complicated by having comparable collision and radiative
decay times. The techniques for treating such systems were first 
described by Bates and McWhirter\cite{Bate61,McWh65}. The physical
processes included in such a collisional-radiative model are:
  \begin{enumerate}
    \item Electron collisional excitation $K_{e,mn}$ and de-excitation
$K_{d,mn}$,
    \item Spontaneous radiative transitions $A_{mn}$,
    \item Electron collisional ionization $K_{i,m}$,
    \item Three-body recombination $K_{r,mn}$,
    \item Radiative recombination $\beta_{rad,m}$,
    \item Dielectronic recombination (for multi-electron atoms) 
      $\beta_{dia,m}$.
\end{enumerate}
where $m$ and $n$ are principal quantum numbers. For the two state
terms, the initial state is the {\em second} number; the final is the
{\em first}. This apparently backward notation allegedly has a mathematical
origin.

A set of $m \geq 1$ equations describes the balance between these processes that
determines the density of the $m$-th excited state of the atom,
$N_{m}$,
\begin{eqnarray}
\frac{d N_{m}}{d t} & = & - N_{m} \left[ n_{e} \left( K_{i,m} + \sum_{n > m}
K_{e,nm} + \sum_{n < m} K_{d,nm} \right) \right. \nonumber \\
& & \qquad \quad \left. + \sum_{n < m} A_{nm} \right] \rule[-1ex]{0ex}{4ex} \nonumber \\
& & \quad + n_{e} \left( \sum_{n > m} N_{n} K_{d,mn} + \sum_{n < m} N_{n} K_{e,mn} 
    \right) \rule[-2ex]{0ex}{4ex} \nonumber \\ & & \quad +
  \sum_{n > m} N_{n} A_{mn} + n_{e} n_{i} ( \beta_{rad,m} + \beta_{dia,m} + n_{e}
                                                   K_{r,m})  \nonumber \\
  & & \quad + n_{e} n_{H_{2}} E_{H_{2},m} + n_{e} n_{H_{2}^{+}}(E_{H_{2}^+,m} + \alpha_{m}),
\label{crmodel}
\end{eqnarray}
where $n_{e}$ is the electron density and $n_{i}$ is the density of ions
corresponding to the atom under consideration ($m \rightarrow \infty$,
effectively).
The two terms at the end account for atoms coming from the dissociation
of H$_{2}$ and H$_{2}^{+}$ ($E_{H_{2},m}$ and $E_{H_{2}^+,m}$) and dissociative
recombination of H$_{2}^{+}$ ($\alpha_{m}$).

The first step in developing a collisional radiative model involves
breaking up Eq.~(\ref{crmodel}) into a single equation for the $m=1$ ground state
and a set of equations for all of the excited states, $m > 1$.
Writing the latter as a matrix
equation for $m, n \geq 2$: 
\begin{eqnarray}
\frac{d N_{m}}{dt} & = & n_{e} [ Q_{mn} N_{n} + N_{1} K_{e,m1}
+ n_{i} (\beta_{m} + n_{e} K_{r,m}) \nonumber \\
& & \quad + n_{H_{2}} E_{H_{2},m} +  n_{H_{2}^{+}}(E_{H_{2}^+,m} + \alpha_{m}) ],
\label{crmodel2}
\end{eqnarray}
where the $K_{e,mn}$, $K_{d,mn}$, $A_{mn}$, and $K_{i,m}$ contributions
have been lumped into the $Q_{mn}$ matrix.
Its diagonal terms represent
sinks of $N_{m}$ while the off-diagonal elements are sources
from the other excited states.  The remaining terms in Eq.~(\ref{crmodel2})
are sources coming from the ground state, $\propto N_{1}$; ions,
$\propto n_{i}$; molecules, $\propto n_{H_{2}}$; and molecular ions,
$\propto n_{H_{2}^{+}}$.

The collisional radiative approximation\cite{Bate61,McWh65} consists of assuming
that the excited states equilibrate amongst themselves much more
quickly than the ground state evolves so that
$\frac{d N_{m}}{dt} \rightarrow 0$; 
transport of the excited states is thus ignored.
In contrast, interactions of the ground state with the excited states
occur on time scales comparable to those of transport;
simulating these processes consistently to obtain the
$N_{1}$ in Eq.~(\ref{crmodel2})
is the job of the neutral transport
code.
The collisional radiative approximation allows the inversion of the
$m \geq 2$ matrix equation with these various source
terms, yielding a set of $N_{m}$ for each:  
\begin{equation}
N_{m} = \left( \frac{N^{(i)}_{m}}{N_{1}}\right) N_{1} + 
\left( \frac{N^{(ii)}_{m}}{n_{i}}\right) n_{i}
+ \left( \frac{N^{(iii)}_{m}}{n_{H_{2}}}\right) n_{H_{2}}
+ \left( \frac{N^{(iv)}_{m}}{n_{H_{2}^{+}}}\right) n_{H_{2}^{+}},
\label{crnm}
\end{equation}
where:
\begin{eqnarray}
  \frac{N^{(i)}_{m}}{N_{1}} & = &  - (Q^{-1})_{mn} K_{e,n1}   \\
  \frac{N^{(ii)}_{m}}{n_{i}} & = & - (Q^{-1})_{mn}  (\beta_{n} + n_{e} K_{r,n})    \\
  \frac{N^{(iii)}_{m}}{n_{H_{2}}} & = & - (Q^{-1})_{mn} E_{H_{2},n}  \\
  \frac{N^{(iv)}_{m}}{n_{H_{2}^{+}}} & = & - (Q^{-1})_{mn} E_{H_{2}^+,n}
                                           - (Q^{-1})_{mn} \alpha_{n}.
\end{eqnarray}
In practice, the solution of these equations involves first
obtaining the various reaction rates,  $K_{e,mn}$,
$K_{d,mn}$, $A_{mn}$, $K_{i,mn}$, $K_{r,mn}$, $\beta_{rad,m}$, $\beta_{dia,m}$,
as a function of $n_{e}$ and $T_{e}$.
At each $n_{e}$, $T_{e}$ pair, the inverse of the $Q_{mn}$ matrix is
computed and then used to generate each of the population
coefficients above.  Further discussion of the molecular terms
is deferred to the descriptions of the collisional radiative
versions of the molecular dissociation, ionization, and
dissociative recombination processes.

Instead focusing on only the ground state and ion
terms in Eq.~(\ref{crnm}), the $m = 1$ equation can becomes
\begin{equation}
\frac{d N_{1}}{d t} = - N_{1} n_{e} S_{eff} + n_{e} n_{i} R_{eff} + {\rm
transport, etc.},
\end{equation}
with
\begin{equation}
S_{eff} = \sum_{m \geq 1} \left( \frac{N^{(i)}_{m}}{N_{1}}\right) K_{i,m},
\end{equation}
and
\begin{equation}
R_{eff} = - \sum_{m \geq 2} \left( \frac{N^{(ii)}_{m}}{n_{i}}\right) K_{i,m}
+ \sum_{m \geq 1} \left( \beta_{rad,m} + \beta_{dia,m} + n_{e} K_{r,m} \right).
\end{equation}
$S_{eff}$ is thus the effective ionization rate of the ground state
neutral, incorporating the multi-step processes associated with all
of the excited states. Likewise, $R_{eff}$ is the effective recombination
rate. Both are functions of $n_{e}$ and $T_{e}$. Selected
values of
the population coefficients  $N^{(i)}_{m}/N_{1}$ and
$N^{(ii)}_{m}/n_{i}$ may also be tabulated or used to compute
photon emission rates for specific transitions of interest
(e.g., Balmer-$\alpha$).

The above assumes that just the ground state evolves on the transport
time scales relevant to DEGAS 2. If metastable excited states
are present in the system (e.g., helium's $2^{1}S$ and $2^{3}S$
states or hydrogen's $n = 2$ in an optically thick region),
they may need to be treated on this same footing. The need
for such a treatment and its formulation go beyond the
scope of this presentation. See the discussion below on
helium for some additional references.

Computing the electron energy sources and sinks associated with 
Eq.~(\ref{crmodel}) involves accounting for the energy exchanges
in each of the processes appearing in Eq.~(\ref{crmodel}):
\begin{eqnarray}
\frac{d E_{e}}{d t} & = & E_{1} \left( - N_{1} n_{e} S_{eff} + 
n_{e} n_{i} R_{eff}  \right) \nonumber \\
& & \qquad - \sum_{\stackrel{m,n}{\scriptstyle n < m}} N_{m} A_{nm}
\left( E_{n} - E_{m} \right) \nonumber \\
& & \qquad - n_{e} n_{i} 
\sum_{m \geq 1} \beta_{rad,m} \left( \langle E_{rad,m} \rangle + E_{m}
\right) \nonumber \\
& & \qquad - n_{e} n_{i} 
\sum_{m \geq 1} \beta_{dia,m} \left( \langle E_{dia,m} \rangle + E_{m}
\right),
% Note that in Goto's notation, N_{m} = n_e R_0 n_i + n_e R_1 n_(1)
\label{creloss}
\end{eqnarray}
where
\begin{itemize}
  \item $E_{1}$ is the ionization energy,
  \item the first term is the loss due to ionization,
  \item the second is the gain due to recombination. 
  \item The sums represent lost photons from radiative transitions and
recombination,
  \item $\langle E_{rad,m} \rangle$ and 
$\langle E_{dia,m} \rangle$ are the average energies of electrons in
radiative and dielectronic recombination, respectively.
\end{itemize}
The two average energies are computed using 
the expression given by Reiter\cite{Reit93},
\begin{equation}
\langle E_{x,m} \rangle  =  T_{e} \left[ \frac{3}{2} 
   + \frac{T_{e}}{\beta_{x,m}} \frac{ d \beta_{x,m}}{d T_{e}} \right],
\end{equation}
where $x$ denotes either $rad$ or $dia$.

Using Eq.~(\ref{crnm}), Eq.~(\ref{creloss}) becomes
\begin{equation}
\frac{d E_{e}}{d t} = - E_{1} \left( N_{1} n_{e} S_{eff} - 
n_{e} n_{i} R_{eff}  \right) - N_{1} E_{\rm loss}^{(i)} 
- n_{i} E_{\rm loss}^{(ii)},
\label{eloss}
\end{equation}
where the terms $E_{\rm loss}^{(i)}$, $E_{\rm loss}^{(ii)}$ represent the loss
rates due to coupling to the ground state and continuum, respectively.
A different arrangement of these terms which can be
generalized to the case involving multiple ``transported
states'' has also been used:
\begin{equation}
\frac{d E_{e}}{d t} = - N_{1} n_{e} E_{\rm loss}^{(i)\prime} +
n_{i} n_{e} E_{\rm source}^{(ii)\prime} - n_{i} n_{e} 
E_{\rm loss}^{(ii)\prime}.
\label{elossprime}
\end{equation}
The relationships between the terms are nearly obvious, but 
potentially prone to confusion.  For clarity, they are:
% In the interest of real clarity, the ``unprimed'' E_loss expressions
% are the ones appearing in ehr2.dat.  These are converted into the
% ``primed'' expressions by reactionwrite, via the
% following equations.  The helium data files
% start out with the ``primed'' expressions.  This should be
% clarified in the specific discussion below.
\begin{eqnarray}
E_{\rm loss}^{(i)\prime} & = & E_{1} S_{eff} + E_{\rm loss}^{(i)}
/ n_{e}, \\
E_{\rm source}^{(ii)\prime} & = & E_{1} R_{eff}, \\
E_{\rm loss}^{(ii)\prime} & = & E_{\rm loss}^{(ii)} / n_{e}.
\end{eqnarray}
Note that the ``unprimed'' energies have dimensions of
power, and the ``primed'' versions have dimensions of
power times volume.  All are stored in tabular form
as functions of $n_{e}$ and $T_{e}$ for
use in DEGAS 2.

\subsection{Hydrogen Collisional Radiative Model}
\label{HCR}

The code we use to solve these equations for
the hydrogen atom, {\tt collrad}, 
is based on the one described
by Weisheit\cite{Weis75}. The original
code assumed that ionization and recombination were in balance 
so that the neutral density could be expressed 
in terms of the ion density.  While this might be valid in
an astrophysical context, it is horribly wrong in fusion devices.
The first major change to the code is to have
it use Eq.~(\ref{crnm}) instead. The other
change is to replace the cross sections
with those provided in the Janev-Smith database\cite{Jane93}.
The resulting collisional-radiative data are believed to be
in good agreement with ADAS. Although the {\tt collrad} code is not
itself distributed with DEGAS 2, it can be obtained from the DEGAS 2
authors (see Sec.~\ref{help}).

The most recent data produced by the {\tt collrad} code
are contained in the file {\tt ehr5.dat}, which are then
incorporated in the {\tt hionize5.nc} and {\tt hrecombine5.nc}
files via the {\tt reactionwrite} code.  Older versions of this
file are included with DEGAS 2 for the purpose of verification,
including with the original DEGAS code.  The corresponding 
ionization netCDF
files in the {\tt degas2/data} directory are noted in parentheses;
the recombination files are numbered in the same way.
\begin{description}
  \item[{\tt eh.dat}] The file generated by Weisheit's original code.
  \item[{\tt ehr1.dat}] The result of modifying Weisheit's code to use
Eq.~(\ref{crnm}) ({\tt hionize1.nc}).
  \item[{\tt ehr2.dat}] Comprehensive update of the {\tt collrad}
code, primarily the switch to the Janev-Smith database\cite{Jane93}
({\tt hionize2.nc}).
  \item[{\tt ehr3.dat}] Minor modifications associated with the removal
of an erroneous factor of $0.75$ multiplying $A_{12}$ ({\tt hionize.nc}).
  \item[{\tt ehr4.dat}] Were {\em not} produced by {\tt collrad}, but via
ADAS.  Flaws in these calculations were subsequently identified; these
data should {\em not} be used.  Note that these data are distinct from
the hydrogen data one would obtain from either the ADAS or OpenADAS
libraries ({\tt hionize4.nc}; at this point the authors decided to
increment the number in the netCDF file name rather than continue
to modify {\tt hionize.nc}).
  \item[{\tt ehr5.dat}] Produced by the {\tt collrad} code following
``convergent close coupling'' (CCC) computed updates to the $n = 1$, 3, 4,
and 5 excitation rates; the fits to these cross sections were provided
by Janev and Reiter ({\tt hionize5.nc}).
\end{description}
See also Sec.~\ref{bestpractices}.

The {\tt ehr5.dat} file contains
\begin{enumerate}
  \item Ionization rate $S_{eff}$ in cm$^{3}$ s$^{-1}$,
  \item Recombination rate $R_{eff}$ in cm$^{3}$ s$^{-1}$,
  \item Neutral electron losses $E_{\rm loss}^{(i)}$ in erg s$^{-1}$,
  \item Continuum electron losses $E_{\rm loss}^{(ii)}$ in erg s$^{-1}$,
  \item Neutral ``n=3 / n=1'', $N_{3}^{(i)} / N_{1}$,
  \item Continuum ``n=3 / n=1'', $N_{3}^{(ii)} / n_{i}$,
  \item Neutral ``n=2 / n=1'', $N_{2}^{(i)} / N_{1}$,
  \item Continuum ``n=2 / n=1'', $N_{2}^{(ii)} / n_{i}$,
  \item Neutral ``n=4 / n=1'', $N_{4}^{(i)} / N_{1}$,
  \item Continuum ``n=4 / n=1'', $N_{4}^{(ii)} / n_{i}$,
  \item Neutral ``n=5 / n=1'', $N_{5}^{(i)} / N_{1}$,
  \item Continuum ``n=5 / n=1'', $N_{5}^{(ii)} / n_{i}$,
  \item Neutral ``n=6 / n=1'', $N_{6}^{(i)} / N_{1}$,
  \item Continuum ``n=6 / n=1'', $N_{6}^{(ii)} / n_{i}$,
  \item Neutral ``n=7 / n=1'', $N_{7}^{(i)} / N_{1}$,
  \item Continuum ``n=7 / n=1'', $N_{7}^{(ii)} / n_{i}$,
  \item Neutral ``n=8 / n=1'', $N_{8}^{(i)} / N_{1}$,
  \item Continuum ``n=8 / n=1'', $N_{8}^{(ii)} / n_{i}$,
  \item Neutral ``n=9 / n=1'', $N_{9}^{(i)} / N_{1}$,
  \item Continuum ``n=9 / n=1'', $N_{9}^{(ii)} / n_{i}$.
\end{enumerate}
Figure~\ref{ehr2} shows the
variation of the effective ionization and recombination rates
with electron density and temperature.

\begin{figure}
\begin{center}
\leavevmode
\includegraphics[width=5in]{Ionization_Recombination_Plot}
\caption{Effective hydrogen ionization and recombination rates as a function
of electron temperature for selected electron densities. {\em Note that 
these are actually data from} {\tt ehr2.dat}. {\em The most recent
data differ only quantitatively.} }
\label{ehr2}
\end{center}
\end{figure}

The portion of the file for each quantity consists of 15 separate sections 
(labeled by {\tt jn}), corresponding to the various electron density values,
\begin{equation}
n_{e}({\tt jn}) =  10^{10 + ({\tt jn} - 1)/2} \: {\rm cm}^{-3}.
\end{equation}
Each section consists of 60 data values, one for each $T_{e} ({\tt jt})$,
\begin{equation}
T_{e}({\tt jt}) =  10^{-1.2 + ({\tt jt} - 1)/10} \: {\rm eV}.
\end{equation}
The index {\tt jt} starts at 1 in the first column and row and proceeds 
row-by-row. There are 10 rows of 6 entries each.

The density of the $n = 3$ excited state can be found with
\begin{equation}
N_{3} = \left( \frac{N^{(i)}_{3}}{N_{1}}\right) N_{1} + 
\left( \frac{N^{(ii)}_{3}}{n_{i}}\right) n_{i}.
\end{equation}
The corresponding emission rate for 
the Balmer-$\alpha$ line (a.k.a. ``H$_{\alpha}$'', $n = 3 \rightarrow 2$)
is given by 
$N_{3} A_{23} = 4.41 \times
10^{7} N_{3}$ cm$^{-3}$ s$^{-1}$.  
Note that the energy of this transition is 
$E_{23} = 13.595 (2^{-2} - 3^{-2})$ eV.
Similar expressions give the
$n=2$ density and the Lyman-$\alpha$ ($n = 2 \rightarrow 1$) emission
rate. 

The data in {\tt ehr5.dat} are transformed into netCDF files 
containing the particular variables DEGAS 2 expects to find
(see Sec.~\ref{addreactions}) by the {\tt reactionwrite}
code.  Separate files for ionization and recombination
are written.  In these files, $S_{eff}$ and $R_{eff}$ become
the ``reaction rates''.  The ``emission rate'' variable 
is obtained by rewriting the above expression
for the Balmer-$\alpha$ emission rate as a power,
per atom (for ionization) or ion (for 
recombination) and per electron:
\begin{eqnarray}
P_{\alpha}^{(i)} & = & E_{23} A_{23} \left( \frac{N_{3}^{(i)}}{N_{1}} \right)
/ n_{e}, \label{pemiti} \\
P_{\alpha}^{(ii)} & = & E_{23} A_{23} \left( \frac{N_{3}^{(ii)}}{n_{i}} \right)
/ n_{e}. \label{pemitii}
\end{eqnarray}
These quantities have dimensions of power times volume.
The wavelengths of this transition for each of the
hydrogen isotopes are also inserted into the
file, $\lambda_{H_{\alpha}} = 6562.80 $ \AA, 
$\lambda_{D_{\alpha}} = 6561.04 $ \AA, and $\lambda_{T_{\alpha}}
= 6560.45 $ \AA.
The electron energy exchange terms in {\tt ehr5.dat} correspond
to those in Eq.~(\ref{eloss}); {\tt reactionwrite} transforms
them into those of Eq.~(\ref{elossprime}).

\subsection{Helium Collisional Radiative Model}
\label{HeCR}

The code used to solve the collisional radiative equations for
the helium atom was developed by Goto\cite{Goto2003},
based on Fujimoto's original\cite{Fuji79}.  This code actually
generates data for two collisional radiative models.  In the 
first, only the ground state ($1^{1}S$) atom is considered slowly
varying (i.e., must be treated explicitly as a transported
species in DEGAS 2); this model was designated as
``formulation II'' by Fujimoto.  In the second, the metastable
$2^{1}S$ and $2^{3}S$ states are also considered
slowly varying; Fujimoto called this ``formulation I''. 
The appropriate species, reactions, and
data needed to do simulations based on this model have been
incorporated into DEGAS 2.  However, investigations using it
have not demonstrated that the added detail yields any significant
differences in the final results.  Moreover, a separate analysis
of the helium collisional radiative equations [Eq.~(\ref{crmodel})]
suggests that this particular
model may not even be valid for typical fusion edge conditions.
For these reasons, we will not discuss this model any further and
will focus only on the simpler ``one state'' model.

The primary modification made to Goto's code in preparation
for using its output with DEGAS 2 were to add the 
computation of the electron energy exchange terms in
Eq.~(\ref{elossprime}). 
The resulting formulation II model and data are, not surprisingly, very
similar in structure to those of the hydrogen collisional 
radiative model.  
The file \verb+he_crII_alad.dat+ (in directory \verb+Aladdin/data+) is
an Aladdin formatted version of these data. These same data
written in a format more like that of {\tt ehr2.dat} are
contained in the file \verb+he_crII.txt+ 
% put in href??
that can be downloaded
from the ``Data'' section of the DEGAS 2 web site.  This file has
been provided as a convenience for those who would just like to
use the data, but not download the entire DEGAS 2 code.

The contents of this file are:
\begin{enumerate}
  \item Ionization rate $S_{eff}$ in cm$^{3}$ s$^{-1}$,
  \item Neutral electron losses $E_{\rm loss}^{(i)\prime}$ 
in eV cm$^{3}$ s$^{-1}$,
  \item Neutral 5877 \AA\ ($3^{3}D \rightarrow 2^{3}P$) emission
rate $P_{5877}^{(i)}$ in eV cm$^{3}$ s$^{-1}$,
  \item Neutral 6680 \AA\ ($3^{1}D \rightarrow 2^{1}P$) emission
rate $P_{6680}^{(i)}$ in eV cm$^{3}$ s$^{-1}$,
  \item Recombination rate $R_{eff}$ in cm$^{3}$ s$^{-1}$,
  \item Continuum electron losses $E_{\rm loss}^{(ii)\prime}$ 
in eV cm$^{3}$ s$^{-1}$,
  \item Continuum electron sources $E_{\rm source}^{(ii)\prime}$ 
in eV cm$^{3}$ s$^{-1}$,
  \item Continuum 5877 \AA\ ($3^{3}D \rightarrow 2^{3}P$) emission
rate $P_{5877}^{(ii)}$ in eV cm$^{3}$ s$^{-1}$,
  \item Continuum 6680 \AA\ ($3^{1}D \rightarrow 2^{1}P$) emission
rate $P_{5877}^{(i)}$ in eV cm$^{3}$ s$^{-1}$,
\end{enumerate}

The portion of the file for each quantity consists of 30 separate sections 
(labeled by {\tt jn}), corresponding to the various electron density values,
\begin{equation}
n_{e}({\tt jn}) =  10^{8 + 8*({\tt jn} - 1)/29} \: {\rm cm}^{-3}.
\end{equation}
Each section consists of 65 data values, one for each $T_{e} ({\tt jt})$,
\begin{equation}
T_{e}({\tt jt}) =  10^{({\tt jt} - 1)/16} \: {\rm eV}.
\end{equation}
The index {\tt jt} starts at 1 in the first column and row and proceeds 
row-by-row. There are 10 rows of 6 entries each, followed by a single
row of 5 entries.

The ``emission rates'' here correspond to the power in those lines 
and are analogous to Eqs.~(\ref{pemiti}) and (\ref{pemitii}) above.
The total emitted power (per unit volume) is obtained at run
time by combining them with the ground state and ion densities:
\begin{equation}
P_{5877} = P_{5877}^{(i)} N_{1^{1}S} n_{e} + 
P_{5877}^{(ii)} n_{i} n_{e}.
\end{equation}
Note that $E_{3^{3}D \rightarrow 2^{3}P} = 2.10955$ eV,
$A_{3^{3}D \rightarrow 2^{3}P} = 7.069 \times 10^{7}$ s$^{-1}$, and
$E_{3^{1}D \rightarrow 2^{1}P} = 1.85606$ eV,
$A_{3^{1}D \rightarrow 2^{1}P} = 6.369 \times 10^{7}$ s$^{-1}$.
The more precise (vacuum) wavelengths for these transitions
are 5877.1 \AA\ and 6679.7 \AA\ respectively. The first
ionization energy for helium is $E_{1}= 24.587$ eV.

The electron energy losses and sources can be used directly
in Eq.~(\ref{elossprime}).

The data in \verb+he_crII_alad.dat+ are transformed into netCDF files
by the {\tt ratecalc} code using the input files
\verb+heionII.input+ and \verb+herecII.input+ (in the {\tt data} directory).
See the \verb+DATA_HISTORY+ file for additional details.

Recent re-examination by D. Reiter, the code authors, 
and others of the cross section fits input to the Goto code have
identified some shortcomings that do quantitatively impact the 
collisional radiative rates.  We are awaiting updates to the
cross sections and the associated fits.  Users with rate
sensitive applications should contact the code authors 
(see Sec.~\ref{help}) regarding the status of this work.

\section{Elastic Scattering}
\label{ElasticScattering}
\index{Elastic scattering}

\subsection{Neutral-Ion Elastic Scattering}
\label{NeutralIon}
\index{Neutral-ion scattering}
\index{Charge exchange}
\index{Spin exchange}
\index{Diffusion cross section}
\index{Total cross section}
\index{Viscosity cross section}
\index{Scattering angle}
\index{Differential cross section}

{\em Note: this subsection was extracted from Ref.~\cite{Kanz00}. 
A few minor changes in notation have been made for consistency.}
Only limited experimental data are available for ion-neutral elastic
scattering at low energies.  The monograph by Janev\cite{Jane87} is one of the
most widely used data sources available, but only contains charge
exchange at moderate to high collision energies and no data for
elastic scattering.  As divertor technology has progressed, low
temperature operating environments have become a standard mode of
operation.  Under these conditions, a purely classical treatment
of charge exchange is no longer sufficient. 
In particular, below
approximately 2 eV for $D^{+} + D$ collisions, the processes
of elastic scattering and and resonant charge exchange can no longer
be
distinguished
due to quantum mechanical effects.  The need for a
comprehensive treatment of ion-neutral elastic collisions at low
energies motivated the Controlled Fusion Atomic Data Center at
ORNL to produce a collection of total differential elastic scattering
cross sections for 31 various isotopic combinations of hydrogen atoms
and molecules\cite{Krst98}.

DEGAS 2 represents an improvement over the
original DEGAS code in that the pseudo-collision algorithm originally
employed in scoring results\cite{Heif82} is replaced with a track-length type
estimator.  This allows efficient determination of momentum and energy
exchange with the background plasma over a wide range of background
densities.  The rate expressions used by the track-length estimator
involve the momentum transfer (or diffusion) 
cross section, $\sigma_{d}$\cite{Bach95}.

A general
expression for the integral cross sections is\cite{Bach95}
\begin{equation}
\sigma_{l} = 2 \pi \int_{0}^{\pi} \left( 1 - \cos \theta \right)^{l}
\sigma(\theta, E_{r}) \sin \theta \: d \theta,
\label{sigmal}
\end{equation}
where $l=0$
corresponds to the total elastic scattering cross section, $l=1$ to the
diffusion cross section and $l=2$ to the viscosity cross
section; $\theta$ is the center-of-mass scattering angle, $E_{r}$ 
is the relative
energy of the collision, and $\sigma(\theta,E_{r})$ is the differential elastic
scattering cross section.  The total elastic collision cross section,
$\sigma_{t}$, is used in Monte Carlo simulations to determine the time between
elastic collisions.  In this manner the size of the total elastic
collision cross section influences the execution time of numerical
simulations.

A simulation technique that reduces the total
collision cross section while keeping the momentum transfer cross
section invariant can be developed.  Using the expression for $\sigma_{d}$ and
dividing Eq.~(\ref{sigmal}) into small and large angle components,
\begin{equation}
\sigma_{d} (E_{r}) = 2 \pi \int_{0}^{\theta_{\rm min}} C(E_{r})
( 1 - \cos \theta) \delta(\theta - \theta_{\rm min}) \: d \theta 
+ 2 \pi \int^{\pi}_{\theta_{\rm min}} 
( 1 - \cos \theta) \sigma(\theta, E_{r}) \: d \theta.
\label{sigmad}
\end{equation}
Substituting $\sigma(\theta,E_{r})$ from the ORNL data set here and directly
computing $\sigma_{d}(E_{r})$  [via Eq.~(\ref{sigmal})], 
allows Eq.~(\ref{sigmad}) to be solved for
$C(E_{r})$ given a particular value of $\theta_{\rm min}$.  
The resulting value of $C(E_{r})$
is then used in a similar expression for the viscosity cross section.
This provides viscosity cross sections accurate to second order in
$\theta_{\rm min}$.  The value of $\theta_{\rm min}$ 
is chosen so as to reduce the total elastic
cross section as much as possible while retaining physical realism.

For low energy elastic collisions a significant portion of the
collisions involve small-angle deflections.  Due to this behavior, the
majority of the reduction in total elastic cross section is obtained
with a five-degree minimum scattering angle (Fig.~\ref{rjkfig1}).  Minimum
scattering angles greater than five degrees result in little
computational savings.  The resulting reduction in $\sigma_{t}$ 
is approximately
a factor of two at relative energies $E_{r} > 10^{-1}$ eV.  The values for 
$\sigma_{d}$
are unchanged during application of a small-angle cut-off. 

\begin{figure}
\begin{center}
\leavevmode
\includegraphics[width=5in]{rjk_fig1}
\caption{Comparison of total and diffusion cross sections for $D^{+} + D$
collisions with and without the small-angle cut-off approximation.}
\label{rjkfig1}
\end{center}
\end{figure}

The
use of this small-angle cut-off leads to an efficient algorithm for
simulating the dynamics of elastic collisions in Monte Carlo
simulations.  Previous models based on classical collision kinetics
for direct elastic scattering face difficulties in that the cross
section is a rapidly varying function of the impact parameter\cite{Bach95}.  
This
precludes the use of efficient table-look-up methods and forces a time
consuming evaluation of the collision integral for each individual
interaction.

Utilization of the differential elastic cross
section, $\sigma(\theta,E_{r})$, 
allows development of a table-look-up interpolation
for both direct elastic scattering and resonant charge exchange.
Following the work by Abou-Gabal and Emmert\cite{Abou91}, 
we divide the total
elastic cross section into $n$ pieces and interpret the result as a
cumulative probability distribution function,
\begin{equation}
P_{i} = 2 \pi \int_{-1}^{\cos \Theta_{i}} 
\frac{\sigma(\theta, E_{r})}{\sigma_{t}(E_{r})} \: d( \cos \theta),
\: i = 0, 1, \ldots n.
\label{scatprob}
\end{equation} 
Equation~\ref{scatprob}
is inverted to obtain the scattering angle as a function of relative
energy and $n$ equally spaced values of the cumulative probability
spacing, $\Theta_{i}(P_{i},E_{r})$, where $\Theta_{i}$ 
is the scattering angle in the
center-of-mass frame.  Examples of $\Theta_{i}(P_{i},E_{r})$
for the $D^{+} + D$ and $D^{+} + D_{2}$
collisions at a center-of-mass energy of 1 eV are presented in 
Fig.~\ref{rjkfig2}.
At this energy, inclusion of direct elastic scattering increases the
probability that colliding particles undergo small angle deflections.
Consideration of charge exchange alone would shift the probability
distribution to favor collisions resulting in $\Theta = \pi$.  As evident in
Fig.~\ref{rjkfig2}, $\Theta_{i}(P_{i},E_{r})$ for like-species 
collisions exhibits a higher
probability of undergoing large-angle collisions compared to the
ion-molecular collision case.  This is due to the absence of a charge
exchange contribution to the scattering function for ion-molecular
elastic collisions.  Monte Carlo collision processing begins with the
sampling of a suitable ion from a Maxwellian distribution at the local
ion temperature, yielding, $E_{r}$.  The $P_{i}$ 
are next sampled from a uniform
distribution.  The collision angle follows from a bi-linear
interpolation into this data table without evaluation of complex
integral functions.

\begin{figure}
\begin{center}
\leavevmode
\includegraphics[width=5in]{rjk_fig2}
\caption{Scattering function $\Theta_{i}(P_{i},E_{r})$ for
$D^{+} + D$ and $D^{+} + D_{2}$ ion-neutral elastic
interactions at an interaction energy of 1
eV.}
\label{rjkfig2}
\end{center}
\end{figure}

Previous treatments of charge exchange in
Monte Carlo algorithms employed a $180^{\circ}$ collision scattering model in
the center-of-mass frame.  The approach is simple to implement,
requiring just the exchange of the ion and neutral velocity vectors
and yields a momentum transfer cross section of about $2 \sigma_{cx}$, 
where $\sigma_{cx}$
is the charge exchange cross section.  This is, in fact, roughly equal
to the total diffusion cross section in ion-neutral
collisions, $\sigma_{d} \sim 2 \sigma_{cx}$ \cite{Dalg58,Hols52,Wade79}
The accuracy of this simple
expression is demonstrated in Fig.~\ref{rjkfig3}, which compares the 
diffusion cross section with results for the ``spin exchange'' cross
section from the ORNL data\cite{Krst98}.  
Spin exchange is defined as the exchange
of the intrinsic quantum mechanical property of particle spin; for
distinguishable particles this is equivalent to charge exchange.  For
like-species collisions, spin exchange and charge exchange are 
equivalent only within classical regimes.  The relation 
$\sigma_{d} \sim 2 \sigma_{cx}$
holds within the classical regime above approximately 2 eV for $D^{+} + D$.
At lower energies, the particles must be treated as indistinguishable
and separation of direct elastic scattering and charge exchange is no
longer possible.  At these low energies, spin exchange is no longer
equivalent to charge exchange and the relation 
$\sigma_{d} \sim 2 \sigma_{cx}$ no longer
holds.  This illustrates the difficulties in treating ion-neutral
elastic collisions by charge exchange alone under low-temperature
detached conditions.

\begin{figure}
\begin{center}
\leavevmode
\includegraphics[width=5in]{rjk_fig3}
\caption{Comparison of diffusion and charge exchange cross sections for 
$D^{+} + D$ elastic collisions.}
\label{rjkfig3}
\end{center}
\end{figure}

The addition of these data present the user with yet another option for 
handling hydrogen charge exchange in DEGAS 2. Figure~\ref{cxsummary} 
compares the new CFADC data (labeled ``Krstic'') with previous 
versions and with experimental data. The cross sections have been
plotted as a function of the relative velocity so that both the
hydrogen and deuterium data can be included. The Krstic cross sections
are the ``spin exchange'' cross sections described in Ref.~\cite{Krst98}.
They are thus {\em smaller} than the total elastic scattering
cross sections that are to be used in DEGAS 2. This underscores the point
that these new elastic scattering cross sections are intended to
{\em replace} the existing charge exchange data, not used in addition to
them.

The default data used heretofore
in DEGAS 2 have been computed from the ``Janev-Smith'' fit\cite{Jane93}.
This fit was derived from older data published by Barnett\cite{Barn90}.
Barnett's data were a compilation of various experimental results. 
Included was the work by Newman\cite{Newm82}. Why the
Barnett compilation diverges slightly from the 
Newman data at the lowest energies is
not clear. The Reviere\cite{Revi71} curve in Fig.~\ref{cxsummary} represents the
charge exchange cross section employed in the original DEGAS code. Not
surprisingly, they are used in DEGAS 2 in the {\tt Degas\_box\_bench}
example. Note that Reviere's intended application was neutral beam
penetration. Hence, the lack of agreement at divertor-relevant 
velocities is not surprising.

\begin{figure}
\begin{center}
\leavevmode
\includegraphics[width=5in]{CX_Summary_Data}
\caption{Comparison of various charge exchange cross sections. The curve 
labeled ``Janev-Smith'' represents the default data used in DEGAS 2
prior to the addition of the newer CFADC elastic scattering data,
labeled here by ``Krstic'' (one curve for H, one for D).}
\label{cxsummary}
\end{center}
\end{figure}

\subsection{Neutral-Neutral Elastic Scattering}
\label{NeutralNeutral}

% Do these Include: statement regarding cross species rates, 
% pointer to balance checks.
% Convergence criterion.

Low temperature (a few eV or less) plasma conditions may result in
neutral densities that are comparable to the plasma density. In such
cases, collisions amongst the neutral species can be important in 
determining the exchange of momentum and energy in the system. For example, hydrogen
molecules generated at a surface typically have energies comparable to 
room temperature, say, $0.03$ eV. Yet, these molecules may be surrounded
by much more energetic atoms, say 3 to 5 eV,  resulting from molecular
dissociation or charge exchange. Elastic collisions between the two
species can significantly increase the molecules' energy, enabling them
to penetrate much further into the plasma before becoming dissociated
or ionized.

When the neutral-neutral mean free path is very short, 
$K_{n} = \lambda_{\rm mfp} / d \ll 1$ ($d$ is a measure of the 
system size or of gradient scale lengths; $K_{n}$ is the 
Knudsen number) cases, 
the most efficient approach would be to use a fluid model
for the neutral species. This regime is also known as the ``viscous
flow'' regime\cite{OHan89}. However, $K_{n}$ is usually close to one in most
fusion problems, even in vacuum regions. In that case, a kinetic
treatment is required. Fortunately, such ``transition regime'' cases
can be handled effectively with a Monte Carlo code. When $K_{n} \gg 1$,
the ``molecular flow'' regime, neutral-neutral collisions can be
neglected. For a practical discussion
of these three regimes as they apply to the tokamak divertor, see
Ref.~\cite{Niem95}.

The great difficulty with treating neutral-neutral collisions is that they
make the problem nonlinear.  Namely, the collision kernel is then a 
quadratic function of the neutral distribution function rather than a 
linear one. The exact problem is very difficult to solve, and one is
forced to resort to approximations. Reiter et al. addressed this issue
first with the EIRENE code and concluded that an iterative ``BGK'' approach
(referring to the original algorithm developed by Bhatnagar,
Gross, and Krook\cite{Bhat54}; proposed independently by 
Welander\cite{Wela54}) was sufficient\cite{Mayxx,Reit97}. The BGK
approximation replaces the collision integral with a much simpler 
expression in which the distribution function relaxes to a Maxwellian
distribution
with a velocity independent collision time. A physically relevant result
is ensured by determining the parameters of the Maxwellian distribution so as
to enforce conservation of mass, momentum, and energy in the collisions.
The collision time remains arbitrary and is chosen so as to reproduce
measured viscosity and diffusion rates; this procedure is discussed in
greater detail in Sec.~\ref{BGKRates}.

Neutral-neutral collisions are included in DEGAS 2 in exactly the same
way as in EIRENE. While more details on nature of the BGK algorithm
are provided in Kanzleiter's thesis\cite{Kanz99}, the essential elements
and references are covered in the more readily available Ref.~\cite{Reit97}.
The BGK algorithm approximates the binary collision integral for species $i$ as
\begin{equation}
\frac{\partial f_{i}}{\partial t} \approx \frac{M_{i} - f_{i}}{\tau_{i}}
+ \sum_{j \neq i} \frac{M_{ij} - f_{i}}{\tau_{ij}},
\label{bgk}
\end{equation}
where the $M$ represent Maxwellian distributions,
\begin{equation}
M_{\alpha} \equiv n_{\alpha} \left( \frac{m_{\alpha}}{2 \pi k T_{\alpha}}
\right)^{3/2} \exp \left[ - \frac{m_{\alpha} 
(\vec{v}_{\alpha} - \vec{U}_{\alpha})^{2}}{2 k T_{\alpha}} \right].
\label{bgkMax}
\end{equation}
Here, the individual particle velocity is $\vec{v}_{\alpha}$; the mass
$m_{\alpha}$ is always that of the species represented on the
left-hand side of Eq.~(\ref{bgk}).
In Eq.~(\ref{bgk}), a single subscript refers to a physical species. 
Dual subscripts below refer to fictitious, ``mixed'' species introduced
for the purpose of handling binary collisions between unlike neutral
species. 

Conservation of mass, momentum, and energy in the self-collision case
yields
\begin{eqnarray}
n_{i} & = & \int_{-\infty}^{+\infty} d^{3}v f_{i}(\vec{v}), \\
\label{nmoment}
n_{i} m_{i} \vec{U}_{i} & = & \int_{-\infty}^{+\infty} d^{3}v m_{i} 
\vec{v} f_{i}(\vec{v}), \\
\label{vmoment}
\frac{3}{2} n_{i} k T_{i} & = & \int_{-\infty}^{+\infty} d^{3}v 
\frac{1}{2} m_{i} v^{2} f_{i}(\vec{v}).
\label{tmoment}
\end{eqnarray}

Conservation of mass, momentum, and energy are also enforced in the
the mixed-species case to determine the $M_{ij}$. Additional constraints
are obtained by requiring that the ratio of the momentum and energy
relaxation rates is the same as for the full collision integral.
The results are\cite{Mayxx,Reit97,Kanz99} are
\begin{eqnarray}
\vec{U}_{ij} & = & \vec{U}_{ji} = \frac{m_{i} \vec{U}_{i} + m_{j} \vec{U}_{j}}
{m_{i} + m_{j}}, \\
k T_{ij} & = & k T_{i}  - \frac{2 m_{i} m_{j}}{(m_{i} + m_{j})^{2}}
\left[ (k T_{i} - k T_{j} ) - \frac{m_{i}}{6} (\vec{U}_{i} - \vec{u}_{j})^{2}
\right].
\end{eqnarray}
One additional important constraint on the collision times is obtained,
\begin{equation}
n_{j} \tau_{ij} = n_{i} \tau_{ji}.
\end{equation}
With $\tau_{ij} \equiv (n_{j} \langle \sigma v \rangle_{ij})^{-1}$
representing the typical time between physical species $i$ and 
the fictitious species represented by $M_{ij}$,
this becomes just
\begin{equation}
\langle \sigma v \rangle_{ij} = \langle \sigma v \rangle_{ji}.
\label{sigvconstraint}
\end{equation}
May\cite{Mayxx,Reit97} sets the values of the reaction rates for
mixed collisions using empirical diffusion coefficients; 
the rates for self-collisions
come from measured viscosity coefficients. The end result is that
in both cases the reaction rate $\langle \sigma v \rangle_{i}
\propto T_{i}^{1/4}$ 
or $\langle \sigma v \rangle_{ij}
\propto T_{ij}^{1/4}$.  The validity of these expressions is examined in the 
next subsection.

The practical implementation of the algorithm is iterative. By default, the
initial iteration of the code proceeds without neutral-neutral collisions.
The reason for this is that the neutral ``background'' with which the 
test species collide initially has zero density. A nonzero density
could be specified during the setup of the background, if desired.
The fluid moments of Eqs.~(\ref{nmoment}), (\ref{vmoment}), and (\ref{tmoment})
are computed at the end of the iteration and inserted into the background
density, velocity, and temperature arrays for the neutral background 
species. The calculation is then restarted; this time the test neutral
species scatter against the neutral background. At the end of the run,
the fluid moments are again updated. The iterations continue either a
specified number of times or until some convergence criterion is
satisfied. 

Accurately simulating neutral-neutral scattering requires greater care
in setting up the geometry and deciding on the number of flights than
is exercised in linear runs. First,
the entire problem space must be divided up finely enough to resolve all
anticipated spatial variations of the neutral density
and the other fluid moments. In comparison, 
without neutral-neutral
scattering, the vacuum regions can be treated as large monolithic zones.
The variation of the neutral parameters between iterations is asymptotically
limited by Monte Carlo noise. In other words, if too few flights are used,
the background densities are not accurately determined, causing the
elastic scattering of the test species to be doubly inaccurate.
The variations in neutral density, etc. in a given simulation will be more
effectively reduced by increasing the number of flights than by 
carrying out more iterations.

\subsection{Verification of BGK Model and Rates}
\label{BGKRates}
\index{BGK}
\index{Viscosity}
\index{Diffusion}
\index{Neutral-neutral scattering}
\index{Maxwell molecules}

As described by May\cite{Mayxx}, the reaction rates for the BGK algorithm
can be determined by matching physical quantities, the viscosity and
diffusion cofficients, against measured values. However, since May's thesis
work and the publication of Reiter\cite{Reit97}, the ORNL 
CFADC has
computed differential scattering cross sections \cite{Krst98}, 
analogous to those
described in Sec.~\ref{NeutralIon}, for these neutral-neutral systems.
In this subsection we will show that the two approaches are largely
consistent, at least near room temperature.

We proceed in a manner analogous to May. Namely, we will obtain from the
CFADC differential cross sections values for the viscosity and diffusion
coefficients. These will be transformed
into reaction rates that can be compared with those described in 
Sec.~\ref{NeutralNeutral}.

For self-collisions, consider the viscosity. The general expression 
derived in Chapman and Cowling\cite{Chap70} [their Eq.~(10.1,4)] is
\begin{equation}
\mu_{1} = \frac{5}{8} \frac{k T}{\Omega_{1}^{(2)}(2)},
\label{visc1}
\end{equation}
where $\Omega_{1}^{(l)}{(n)}$ corresponds directly to the $\Omega^{l,n}$
integrals described by Bachmann\cite{Bach95}.  The subscript ``1'' here 
is to be associated with the cross section and refers to collisions of species 
``1'' with itself (as opposed to collisions between species ``1'' and 
``2'' which will be labeled with ``12'' below).

Likewise, we have for the diffusion coefficient [Eq.~(9.81,1) in
Ref.~\cite{Chap70}],
\begin{equation}
D_{12} = \frac{3}{16} \frac{k T}{(n_{1} + n_{2})m_{r} \Omega_{12}^{(1)}(1)},
\label{diff1}
\end{equation}
where $n_{1}$ and $n_{2}$ are the species densities and $m_{r}$ is the
reduced mass for the system, $m_{r} = m_{1} m_{2} / (m_{1} + m_{2})$.

The $\Omega^{(l)}(n)$ integrals represent averaging of the collision 
cross sections over Maxwellian distributions for both species. The
{\tt ratecalc} code is able to integrate these cross sections over
a single Maxwellian distribution,
namely, the $I_{l,n}$ described in Reiter\cite{Reit93}
and the \href{ratecalc.pdf#page.3}{documentation} for {\tt ratecalc}.
Fortunately, the double integral we need here can be found as a limiting case of 
the single integral:
\begin{equation}
\langle I_{l,n} \rangle(u_{p},u_{n}) = \lim_{u^{\prime} \rightarrow 0} 
I_{l,n}(u_{p}^{\prime},u^{\prime}),
\label{Iln1}
\end{equation}
where $u_{p}^{\prime 2} \equiv u_{p}^{2} + u_{n}^{2}$, and 
$u_{p} = \sqrt{2 T_{p} / m_{p}}$ is the thermal velocity of the
second species. Likewise, $u_{n} = \sqrt{2 T_{n} / m_{n}}$ is the
thermal velocity associated with the Maxwellian distribution used in the
averaging over $u$.  The $\langle I_{l,n}\rangle$ is then related to
$\Omega^{(l)}(n)$,
\begin{equation}
\langle I_{l,n} \rangle = 8 \left( \frac{u_{p}^{\prime}}{u_{p}} \right)^{n}
\Omega^{(l)}(n/2).
\label{Iln2}
\end{equation}
The same result is quoted in Ref.~\cite{Bach95} and in \cite{Reit93},
although in the latter case the variable equivalent to $u_{p}^{\prime}$
is incorrectly defined.

With these expressions, we can compute the single species viscosity and
two species diffusion coefficients from the CFADC differential scattering
cross sections. The corresponding expressions for the ``Maxwell molecule''
interaction [intermolecular force given by $\kappa / r^{5}$; e.g., see
Eq.~(12.1,1) in Ref.~\cite{Chap70})] are
\begin{equation}
\mu_{1} = \frac{k T}{3 \pi A_{2}(5)} \sqrt{\frac{2 m}{\kappa}},
\label{visc2}
\end{equation}
where $A_{2}(5) = 0.436$ is the result of a dimensionless integral.
And from Eq.~(14.2,2) of Ref.~\cite{Chap70},
\begin{equation}
D_{12} = \frac{k T}{2 \pi A_{1}(5)} \frac{1}{n \sqrt{\kappa_{12} m_{r}}},
\label{diff2}
\end{equation}
where $A_{1}(5) = 0.422$ is the result of another dimensionless integration.

The reaction rates needed for the BGK algorithm are related to the constant
in the Maxwell molecule force law by (see Ref.~\cite{Mors63} or \cite{Mayxx}),
\begin{equation}
\langle \sigma v \rangle_{11} = 2 \sqrt{2} \pi A_{1}(5)
\sqrt{\frac{\kappa}{m}},
\label{sv11}
\end{equation}
and 
\begin{equation}
\langle \sigma v \rangle_{12} = 2 \pi A_{1}(5)
\sqrt{\frac{\kappa_{12}}{m_{r}}}.
\label{sv12}
\end{equation}

Now, we can use Eq.~(\ref{visc2}) to eliminate $\kappa$ in Eq.~(\ref{sv11})
in favor of $\mu_{1}$. Then, we can insert  $\mu_{1}$ as given by
Eq.~(\ref{visc1}) to obtain an effective BGK reaction rate that has the
same viscosity as that computed directly from the CFADC differential
cross sections:
\begin{equation}
\langle \sigma v \rangle_{11} = 2.065 \Omega_{1}^{(2)}(2).
\end{equation}
The analogous procedure for the diffusion coefficient yields
\begin{equation}
\langle \sigma v \rangle_{12} = \frac{16}{3} \Omega_{12}^{(1)}(1).
\end{equation}

Finally, we can use Eqs.~(\ref{Iln1}) and (\ref{Iln2}) to
replace the $\Omega$ integrals with the equivalent integrals obtainable
from {\tt ratecalc},
\begin{equation}
\langle \sigma v \rangle_{11} = \frac{2.065}{32} 
\lim_{u^{\prime} \rightarrow 0} I_{2,4} (u_{p}^{\prime} = \sqrt{2} u_{p},
u^{\prime}),
\end{equation}
and
\begin{equation}
\langle \sigma v \rangle_{12} = \frac{2}{3} 
\frac{u_{p}^{2}}{u_{p}^{2} + u_{n}^{2}} 
\lim_{u^{\prime} \rightarrow 0} I_{1,2} 
(u_{p}^{\prime} = \sqrt{u_{p}^{2} + u_{n}^{2}}, u^{\prime}).
\label{sv122}
\end{equation}

In evaluating the latter expression, we will assume that one of the
species (e.g., molecules at the wall temperature) are much cooler
than the other, i.e., $u_{n} \ll u_{p}$. Equation~(\ref{sv122}) then
becomes a function of only $u_{p}$.

The CFADC data exist only for the $D + D_{2}$ and $D + D$ systems 
(and isotopic equivalents; molecules are computationally intensive, hence,
the absence of molecule-molecule data). These are compared with the 
May data in Fig.~\ref{scattering}. 

One concern with May's treatment of the $D_{2} + D_{2}$ case is that the 
empirical viscosity formula quoted in Ref.~\cite{Mayxx} does not
appear in the literature, nor is its origin clearly stated. The best
guess for its origin is that it was obtained by rescaling the 
``Fuller'' expression for the diffusion coefficient\cite{Reid87}
into a viscosity using the theoretical expressions to relate the two.
In Fig.~\ref{viscosities}, we compare May's expression with two that do
appear directly in Ref.~\cite{Reid87} (``Chung'' and ``Lewis'') and with
individual data points listed in the CRC Handbook\cite{CRC}.  The 
agreement of these values near room temperature ($\sim 0.03$ eV) is
expected. However, May's formula has a stronger temperature scaling than
the others and exceeds them by about a factor of 15 at 10 eV. The
relevant temperature range for this process is expected to be between
1 and 10 eV.

% NOTE: have not yet made these changes to the code!

Figure~\ref{scattering} shows that at least near room temperature,  
May's expressions are also consistent with the CFADC data. However,
in the 1 to 10 eV range, they are considerably larger than the
CFADC values,
suggesting that the values used in DEGAS 2 (and EIRENE) may be too
large. However, we will 
demonstrate in Sec.~\ref{cmod1d} that the temperature dependence
ascribed to these rates by May is inconsistent with Eq.~(\ref{sigvconstraint}).
Consequently, DEGAS 2 now uses constant neutral-neutral reaction rates,
i.e., independent of the ``background''
temperature.  Since agreement of the May expressions with the CFADC
and empirical viscosities is good at room temperature, 
and since room temperatures are physically relevant (at least for
desorbed D$_{2}$) we fix the reaction rates at its
room temperature (300 K = 0.02585 eV) value according to May's formulas.

\begin{figure}
\begin{center}
\leavevmode
\includegraphics[width=5in]{scattering_plot}
\caption{Comparison of reaction rates for the $D + D_{2}$ and $D + D$ 
systems.  The original rates suggested by May\cite{Mayxx} and effective
rates computed from the CFADC data are shown.}
\label{scattering}
\end{center}
\end{figure}

\begin{figure}
\begin{center}
\leavevmode
\includegraphics[width=5in]{viscosity_plot}
\caption{Comparison of viscosity formulas from May\cite{Mayxx}, 
Reid\cite{Reid87} (the ``Chung'' and ``Lewis'' expressions), and the CRC
Handbook\cite{CRC}.}
\label{viscosities}
\end{center}
\end{figure}

\section{Atomic Physics: Best Practices}
\label{bestpractices}

The rates used for atomic physics processes in DEGAS 2 have evolved over
the years.  Yet, the older rates are retained in the interest of
backward compatibility and verifiability of results.  Details on
the provenance of each set of rates are maintained in the
{\tt DATA\_HISTORY} file in the {\tt degas2/data} directory. 
The following subsections go into more detail on this and
other points users should consider
in setting up their problems.

\subsection{Atomic Hydrogen Ionization and Recombination}

{\tt hionize5} should be used for essentially all applications involving 
atomic hydrogen in a plasma; see Sec.~\ref{HCR} for a detailed description
of earlier versions.  The {\tt hrecombine5} data come
from the same calculation as {\tt hionize5} and represent the best available
data for recombination.  Note, however, that recombination is significant
only in high density ($n_{e} > 10^{20}$ m$^{-3}$), low temperature ($T_{e} < 5$ eV)
plasmas.  If in doubt, add the process to the problem input file and 
generate the background netCDF file, e.g., via {\tt defineback}.
The code will automatically add a recombination source to the
end of the list of sources.
Compare its {\tt source\_total\_current} 
with the others in the problem.  If the recombination source is much
smaller, it need not be included.

\subsection{Helium Ionization and Recombination}

As noted in Sec.~\ref{HeCR}, the simpler ``formulation II'' model
embodied in {\tt heionizeII} and {\tt herecombineII} should suffice for virtually all
applications.  The considerations noted above for hydrogen recombination
also apply to helium. 

\subsection{Hydrogen Ion - Atom Charge Exchange and Elastic Scattering}

A key point made in Sec.~\ref{NeutralIon} that bears repeating is that
{\em all} of the hydrogen ion-atom elastic cross section data provided with DEGAS 2
effectively include classical charge exchange as well.  Thus, only one
of the two processes should be included in a problem; using both 
amounts to doubling the charge exchange reaction rate.  If the 
atom-ion interaction energies in a problem are 
well above 1 eV, a charge exchange process
based on the ORNL data, e.g., {\tt dd\_chargex}
can be chosen; if lower energies are expected, use {\em only}
the corresponding elastic process, e.g., {\tt dd\_05\_elastic.nc}.

A second important point is that the ORNL elastic / charge exchange
data produced by ORNL\cite{Krst98} is considered definitive
and unlikely to be superseded; the other charge exchange data
files are retained only for backward compatibility and verification
exercises.  All of the elastic scattering data, regardless of species,
are made available with both 0 and 5 degree minimum scattering angles.
Users generally should use the more efficient 5 degree versions; the
0 degree files can be used for a verification test in the event
of questions.

The ORNL elastic scattering computations were performed for
specific species; for interaction energies below a few eV, the cross 
sections do indeed differ (see Fig.~\ref{cxsummary}).  Above this
range, where classical charge exchange is the dominant process,
the cross sections are essentially identical when expressed
as a function of the relative projectile velocity.  Should one
need to model a mix of isotopes, one can modify the reaction 
specification, e.g., for {\tt hh\_chargex}, in a local copy 
of the {\tt reactions.input} file to change the reaction
specification from ``species specific'' (with \verb+=>+ 
appearing between reagents and products) to ``generic''
(using \verb+->+ instead), just as was done for older 
charge exchange reactions, like {\tt hchex}.  As obvious
as this equivalence may seem, the detailed considerations
are more complex; the interested user should consult
Ref.~\cite{Krst99} for additional details.

An additional consideration is that great care has been taken with
the files based on the ORNL data to ensure that the energy and
temperature ranges used for the integrals are large enough that
the user can be confident that the track length estimator 
will be accurate.  With the older data files, e.g., {\tt hchex}, 
the limits of the data can be exceeded, causing track length
and collision estimator results to differ.

\subsection{Hydrogen Ion - Molecule Elastic Scattering}

This is a critical process in low temperature, high density
divertor plasmas since it transfers momentum and energy
from the plasma to the (cold) molecules, greatly increasing their
mean free paths, and, hence, their ability to penetrate
into the plasma.  The comments made above regarding ion-atom
elastic scattering largely apply, except that one 
typically ignores charge exchange (unless the molecules are
vibrationally excited).  The only available data are those
from ORNL\cite{Krst98}; again, one should use the files with
the 5 degree scattering cutoff, e.g., {\tt dd2\_05\_elastic.nc}.
The scaling of the cross sections with isotope is less
straightforward than in the ion-atom case; users needing
to model mixed isotope situations should contact
the code authors (see Sec.~\ref{help}).

\subsection{Examples}

Template problem input files for
the indicated situations can be found in the following
examples; they have been updated to reflect the above
guidance.
\begin{description}
  \item[NCSX, UEDGE\_CMod] Basic deuterium with Balmer-$\alpha$ light emission, 
  \item[CMod\_1D] Deuterium with nonlinear elastic scattering; does
not include all Balmer-$\alpha$ processes, 
  \item[He\_GPI] Helium puffed into deuterium plasma, 
  \item[He\_metastables] Single state helium model
in {\tt pr\_boxgen\_both.input} and three state helium model
in {\tt pr\_boxgen\_mboth.input}.
\end{description}
All of the these helium examples include light emission for 
the lines noted in Sec.~\ref{HeCR} as part of the collisional
radiative model.

Note that the other examples' problem input
files have {\em not} been modified in this way
so as to
maintain consistency with previous results documented either in this manual
or in the published literature.

\section{Impurity Atomic Physics}

The ability to incorporate impurity atom (and ion) generalized
collisional radiative rate data from ADAS \cite{ADAS} was developed
in preparation for the research described in Ref.~\cite{Stot2017}.
However, apart from the {\tt adaswrite} (Sec.~\ref{othersetup})
utility that reads
the ADAS {\tt adf11} files and reformats the data into DEGAS 2 netCDF
files, none of that capability is included in this distribution.
Interested users should contact the code authors (Sec.~\ref{help}).

\section{Recycling}

Recycling of both ions and neutral species in DEGAS 2 is controlled by
the set of PMI (``Plasma Material Interactions'', e.g., see 
the PMI \href{pmisetup.pdf#page.1}{setup} file) and one additional
parameter generically referred to as the ``recycling coefficient''.
The value of the recycling coefficient is a constant associated with
each wall and target sector (see the \href{classes.pdf#page.39}{sector class}).
These are prescribed during geometry definition, e.g., via 
the {\tt recyc\_coef} keyword in {\tt definegeometry2d}.

The best way to understand the recycling coefficient, $\cal{R}$,
is to note that $P_{a} = 1 - \cal{R}$ is the probability for
a flight (the ``projectile'') striking that sector to be absorbed.
Adsorption, is effectively 
one of the plasma-material interactions that the
code will choose from in deciding the fate of the projectile.

The second PMI that is always present is referred to as the
``default PMI'' since the code defaults to that PMI should
none of the other prescribed PMI (including adsorption) be
selected.  Since all PMI must have an
associated probability or ``yield'', a ``default PMI'' is
assigned a negative (i.e., unphysical) yield.
The most common default PMI is desorption.

The third PMI we consider here is ``reflection'' or ``backscattering''
in which the incident particle returns immediately to the 
plasma with most of its incident energy.
The data associated with the reflection PMI
are used to compute a probability or yield, $P_{r}$, for
the process.
This probability is, in general, a function of the
projectile energy and angle with respect to the surface normal.

Consider now the common case of an incident hydrogen atom for which
both reflection and desorption (default) PMI's have been specified,
as well as a recycling coefficient $\cal{R}$.  The probability
assigned to the desorption PMI is 
$P_{d} = 1 - P_{r} - P_{a} = {\cal R} - P_{r}$.  In the case
of unity recycling (the surface is in steady state), the physical 
interpretation is that atoms not reflected are adsorbed, 
eventually recombine with another atom, and are desorbed as
molecules.  In a steady state simulation, the time interval
associated with ``eventually'' is meaningless, so the process might
as well be instantaneous.

In the case of less than unity recycling, i.e. net adsorption, 
the value of $P_{d}$ is reduced accordingly.  If ${\cal R} < P_{r}$,
$P_{d}$ is negative and there is no desorption and the
reflection probability is effectively $\cal{R}$.

If ${\cal R} = 0$, the sector is equivalent to an ``exit''.
One operational difference is that flights striking an
exit are terminated immediately.  Flights striking a 
zero recycling sector will still undergo the process of
selecting a PMI, even though the result will always
be adsorption.

\chapter{Running DEGAS 2}

\section{Getting DEGAS 2}

DEGAS 2 is now govern by a 
\href{http://theorycodes.pppl.wikispaces.net/Theory+Department+Codes}{licensing
procedure} developed and managed by the PPPL Theory and Computation
Department. Prospective users of the code should visit this site and
follow the instructions spelled out there.

The DEGAS 2 source code is available in three formats:
\begin{enumerate}
\item {\tt degas2.tar} An uncompressed tar file; extract with:
\begin{verbatim}
tar xf degas2.tar
\end{verbatim}
\item {\tt degas2.tar.Z} The tar file compressed with the {\tt compress}
utility. To extract the code,
\begin{verbatim}
uncompress degas2.tar.Z
tar xf degas2.tar
\end{verbatim}
\item {\tt degas2.tar.gz} The tar file compressed with the {\tt gzip}
utility. To extract the code,
\begin{verbatim}
gunzip degas2.tar.gz
tar xf degas2.tar
\end{verbatim}
\end{enumerate}

There are a number of older versions of the code in this directory. The most
recent is always linked to the names listed above. Please consult with the
authors if you think you need an earlier version.

\section{Getting Supporting Tools}

DEGAS 2 relies on several pieces of utility software.  In some cases (netCDF),
these utilities are absolutely required. These utilities 
can be obtained from the indicated sites. You
will probably want to have your system administrator download, compile,
and install these packages system-wide for your convenience.


\subsection{GNU Software}
\label{gnu}
\index{gmake}
\index{emacs}
\index{Tags}
\index{make}
\index{Makefile}

The DEGAS 2 Makefile automates several chores, including generating and compiling
FORTRAN code, as well as weaving \TeX\ files from from the FWEB source,
The Makefile also controls the application of platform-specific lines of code
and compilation options.  Although not completely necessary, using DEGAS 2
without using GNU Make to run this Makefile would be extraordinarily inconvenient.
The {\tt gmake} (the GNU Make utility) can be found at
\href{https://www.gnu.org/software/make/}{https://www.gnu.org/software/make/}

Likewise, the Emacs editor is not essential, but
its use simplifies a number of tasks.
For instance, Emacs provides a facility which allows the user to search for variable
and routine names across many files. For example, say you're
editing {\tt randomtest.web} and you'd like to see the source
code for {\tt subroutine next\_seed} called there. All you
have to do is execute the Emacs command \verb+M-. next_seed+ (it
will ask for the name of the tags file; the default value of {\tt TAGS}
is correct and you just need to hit {\tt Enter} here). All of DEGAS 2
source and header files are included in this list of tags. Additional
facilities such as a tag-query-replace  (e.g., to change {\em all} 
occurences of a variable name) are available. See the Emacs Info
file under the node ``Tags'' for more information.

% Insert URL for Emacs

\subsection{netCDF}
\label{netcdf}
\index{netCDF}
\index{Emacs}
\index{ncgen}
\index{ncdum}
\index{ncgen-filter}
\index{ncdump-filter}

netCDF provides a device-independent binary file
format. Furthermore, these files are self-describing. Most of the effort
required to read and write these files is handled by macros. These are
documented in the file {\tt netcdf.hweb}. By adding the commands 
in {\tt src/.emacs} to
your .emacs file, Emacs will be able to convert automatically these files
to text and back to binary for perusal and editing.
netCDF is fully documented in the Info facility of Emacs. The netCDF
libraries are required to run DEGAS 2 since virtually all of the input
data is stored in netCDF files. The associated netCDF utilities {\tt ncdump}
(for translating a netCDF file into human-readable text) and {\tt ncgen}
(for generating a netCDF binary file from a text file) are useful for
reading and making minor changes to DEGAS 2 input files.
Two short scripts, {\tt ncdump-filter} and {\tt ncgen-filter}, allow
Emacs to invoke these utilities via the commands in the .emacs file.
These are not part of the netCDF distribution and are included amongst
the DEGAS 2 source files for your convenience. These should be copied
to a directory which is covered by your shell's {\tt PATH} variable.

An additional note: if you use {\tt ncgen}, be sure to either specify
the \verb+-b+ or \verb+-o netcdf_filename+ option to let the utility know
that you want it to generate a netCDF file (see: {\tt man ncgen}).

netCDF can be obtained from  
\href{http://www.unidata.ucar.edu/software/netcdf/}{http://www.unidata.ucar.edu/software/netcdf/}.  
DEGAS 2 has been tested recently with versions 4.2.0, 4.4.1.1,
and 4.4.4.2.  
Note, that the {\tt ncgen}
that comes with V. 4.1.1 and V. 4.2.0 has a bug that causes problems
when creating large files; the work around is to use instead
{\tt ncgen3}.
If you are using {\tt ncgen-filter}, you would need
use {\tt ncgen3} there as well.

\subsection{FWEB}
\label{fweb}
\index{FWEB}
\index{ftangle}
\index{fweave}
\index{TeX}
\index{LaTeX}

The FWEB package provides a facility for code documentation as well
as a powerful preprocessor. The typesetting abilities of FWEB allow authors
to insert \TeX\ (or \LaTeX) comments into their code. In DEGAS 2, the
preprocessor is utilized extensively for writing macros which allow the
code seen by the user to be significantly more powerful than pure FORTRAN.
In fact, the code can be processed into FORTRAN 77 or FORTRAN 90 as
needed. However, FWEB is {\em not optional}. The user must have access
to at least its {\tt ftangle} executable on some system in order to
generate compilable FORTRAN files (once generated, these can be copied
to other systems for compilation).
For examples of the overall structure of FWEB programs, see the *.web
files; most of the macros are located in the header *.hweb files.
FWEB is fully documented in the Info facility of Emacs. 

FWEB has been compiled and run on a wide variety of platforms,
including, for example, the Edison machine at NERSC.  Users
should first try the most recent version:
\href{ftp://ftp.pppl.gov/pub/degas2/fweb_550bc1c.tar.gz}{ftp://ftp.pppl.gov/pub/degas2/fweb\_550bc1c.tar.gz}
Two slightly older versions are also available in the
DEGAS 2 FTP area as alternatives should problems
arise.  The last version distributed by John
Krommes (ca. 1998) can be found at:
\href{ftp://ftp.pppl.gov/pub/fweb/fweb-1.62.tar.gz}{ftp://ftp.pppl.gov/pub/fweb/fweb-1.62.tar.gz}.

The installation instructions distributed with FWEB are
antiquated and needlessly complex for most systems.  A
simple approach should suffice:
\begin{enumerate}
  \item Unpack the tarball in a convenient location.
  \item In the {\tt fweb/Web} directory, run: {\tt ./configure}
  \item The {\tt configure} script will choose a C compiler
for you, likely {\tt gcc} (recommended).  You can override
its choice by setting the {\tt CC} environment variable,
  \begin{enumerate}
    \item If you have already run {\tt configure}, delete the
{\tt config.cache} file before running it again.
    \item Alternatively, one can set the {\tt CC} variable
in {\tt defaults.mk} after running {\tt configure}.
    \item As of this writing, a bug in the ``2017''
version of the {\tt icc} (Intel) compiler prevents it
from working properly with FWEB's default optimization.  The
user can simply delete {\tt -O2} from the {\tt CFLAGS} variable
in {\tt defaults.mk}.
  \end{enumerate}
  \item Run: {\tt make} to compile both {\tt ftangle} and 
{\tt fweave}.
  \item These executables can then simply be copied to a location
in your {\tt PATH}.
\end{enumerate}

To typeset DEGAS 2 source code in ``human-readable'' formats such as DVI,
PostScript, or PDF, the user will need FWEB's {\tt fweave} utility and
a version of \TeX\ or \LaTeX.  E.g., with \verb+gmake tallysetup.pdf+
Note that for most of the {\tt .web} files, {\tt fweave} will 
sound an error beep associated with the ``memory allocation
interface'' macro; this is harmless.

\LaTeX\  is available in lots of places. To begin with, see the
introductory material and pointers on the
official CTAN \TeX\ network page at  
\href{http://www.tex.ac.uk/ctan}{http://www.tex.ac.uk/ctan}. The 
recommended distribution for Unix (including Linux) systems
is \TeX\ Live (\href{http://www.tug.org/texlive/}{http://www.tug.org/texlive/}.
\TeX\ Live is also available for Windows.

\subsection{Triangle}
\label{triangle}
\index{Triangle}
\index{definegeometry2d}

Jonathan Shewchuck's 
{\tt Triangle} program is used by the {\tt definegeometry2d}
code to break polygons up into triangles. The home for {\tt Triangle}
is \\
\href{http://www.cs.cmu.edu/~quake/triangle.html}{http://www.cs.cmu.edu/~quake/triangle.html} 
The {\tt Triangle} V. 1.6 was used in the creation of
the {\tt definegeometry2d} examples described in
Sec.~\ref{defexamples}.
Anyone downloading a more recent version should notify the {\tt DEGAS 2}
authors. The {\tt DEGAS 2} {\tt Makefile} will expect to find
the {\tt Triangle} files, \verb+triangle.c+, \verb+triangle.h+, etc., 
in \verb+$HOME/Triangle+. If your copy is installed elsewhere, you can
either create an appropriate link or modify {\tt Makefile}.

\subsection{Silo and HDF}
\label{silohdf}
\index{Silo}
\index{HDF}
\index{HDF5}

DEGAS 2 produces no direct graphical output.  Instead, the {\tt geomtesta}
utility (see Secs.~\ref{maincode} and the
\href{geomtesta.pdf#page.1}{source code} documentation.) generates
device-independent binary files in either the HDF (more specifically,
HDF Version 4) or Silo formats that can be processed by
external graphical packages, such as VisIt or IDL.  The VisIt
package in particular is freely available \\
(see \href{https://wci.llnl.gov/codes/visit/home.html}{https://wci.llnl.gov/codes/visit/home.html}), can read Silo files directly, and can generate
both 2-D and 3-D images.  A similar alternative to
{\tt geomtesta}, {\tt ucd\_plot}, is now available.  It's
more restrictive in that it requires the Silo library
and works only with 2-D geometries created by 
{\tt definegeometry2d}.  On the other hand, the data
contained in the Silo file are effectively
represented on the actual DEGAS 2 mesh instead of
a pixelated representation.  Visualization of the
{\tt ucd\_plot} Silo has only been performed with VisIt.


The Silo library can be found at \href{https://wci.llnl.gov/codes/silo/index.html}{https://wci.llnl.gov/codes/silo/index.html}.  It can use either its own
PDB I/O driver or HDF5.  In the latter case, the HDF5 library
must also be installed.  For HDF5 and / or HDF4, see
\href{http://www.hdfgroup.org/}{http://www.hdfgroup.org/}.
% HAVE NO EVIDENCE OF THIS BEING A PROBLEM WITH RECENT VERSIONS:
% Note that as of HDF version 4, some of the netCDF functionality is 
% implemented in HDF. To avoid conflicts with the official netCDF libraries,
% do not install the netCDF portion of the HDF distribution.

Note that {\tt definegeometry2d} also produces an HDF or Silo
file, {\tt geomtestc.hdf} or {\tt geomtestc.silo}, as part of
its test of the geometry.  

\subsection{Git}
\label{Git}

The central DEGAS 2 source code is maintained in a Git repository.
In principle, users planning to coordinate with the code authors
can obtain access to the repository, although in practice none have
done so.  Instead, suggested modifications have been passed to the
code authors via e-mail (See Sec.~\ref{help}).

\section{Structure of the DEGAS 2 File System}
\label{structure}
\index{directories}
\index{HOME}
\index{Makefile.local}

Here is an example structure of the directories under the top-level 
{\tt degas2} directory. Those labeled with ``(D)'' represent
source and data directories maintained with CVS. The others are created
by the user as described in the compiling section. Note that the Makefile 
assumes that this {\tt degas2} directory sits in the directory pointed to
be the \verb+$HOME$+ 
environment variable on your workstation. If you have installed
the code elsewhere, see Sec.~\ref{adjustments}. 

\begin{description}
\item[data] Input atomic and surface physics data and information (D)
\item[src] All source code (D)
\item[dg\_carre] Source and executable files for DG and Carre (D)
\item[Doc] Location of the various formats of this document and ancilliary
files (D)
\item[examples] Input and output files from example and benchmark runs (D)
\item[Aladdin] Source code and data for use with IAEA Aladdin package (D)
\item[ALPHA] Object and executable files for Dec Alpha machines
\item[LINUX] Object and executable files for Linux machines
\item[LINUX64] Object and executable files for 64-bit Linux machines
\item[SUN] Object and executable files for Sun Solaris and SunOS machines
\item[CRAY] Pre-processed source code for export to Cray computers
\item[tex] Intermediate files and final DVI documentation files generated 
from source code
\end{description}

If you need to have multiple run directories for a given system (e.g., focussed
on different problems or using different compilers), you can create additional
directories with names like \verb+SUN-foo+ where ``foo'' is any string you like.

\section{Components of the Code}
\label{components}

DEGAS 2 is not a single code, but a complex package of 
setup, test, simulation, and post-processing tools. Which ones are
needed depend on the user's objectives and familiarity with DEGAS 2.
A few additional ``targets'' are listed in the Makefile. These are
either sufficiently out of date or infrequently used to warrant
mention here.

\subsection{Main Code}
\label{maincode}
\index{problemsetup}
\index{readgeometry}
\index{readbackground}
\index{tallysetup}
\index{flighttest}
\index{geomtesta}
\index{ucd\_plot}
\index{definegeometry2d}
\index{defineback}

Virtually every user will be running these executables,
usually in this order.

\begin{description}
\item[problemsetup] Reads a text file describing the species, 
materials, atomic, and surface physics required for the
current problem. The code
accesses the required "reference" data, and compiles them
into a netCDF file (the {\tt problem.nc} file).
\item[definegeometry2d] Generates a DEGAS 2 netCDF geometry file 
from from a text
input file. Input files for simple geometries may be
generated manually; more complex geometries can be specified via other computer
generated text files
files referenced by the input file. For example, {\tt definegeometry2d} is
able to read data produced by the {\tt DG} package 
(extracted from the {\tt SOLPS} distribution
and included with DEGAS 2 for convenience)
or by UEDGE. A newer option suitable for
``main chamber'' geometries of limited spatial
extent is {\tt efit2dg2d}, the output of which can
be pasted directly into a {\tt definegeometry2d}
input file.
Of the three tools for setting up
DEGAS 2 geometries {\tt definegeometry2d} is the most flexible. 
\item[readgeometry] Generates a DEGAS 2 netCDF geometry file from an
existing DEGAS, UEDGE, or SONNET geometry description. This package has
largely been superseded by {\tt definegeometry2d}.  However, users with legacy
input files for the old DEGAS code will still want to use {\tt readgeometry}.
A short text input file helps readgeometry
through the transformation process.
\item[defineback] Uses a short text input file to control the mapping of plasma
data in external text files onto the DEGAS 2 geometry. Designed to be used in
conjunction with {\tt definegeometry2d}, {\tt defineback} can utilize data
produced by other codes, including UEDGE.  In particular, 
{\tt defineback} can be run in an iterative mode with UEDGE 
analogous to that provided by {\tt updatebackground} (see below
and in the source \href{defineback.pdf#page.13}{file}
{\tt defineback.web}).
\item[readbackground] Plasma data in an external file is mapped onto
the zones of the DEGAS 2 geometry. Presently, only DEGAS and UEDGE files
can be used; readgeometry and readbackground read the same data
file. 
\item[updatebackground] Is similar to {\tt readbackground}, but is intended
for use in iterations with the UEDGE code {\em only}. The source
information generated by the initial run of {\tt readbackground} is 
stored in a separate file ({\tt oldsourcefile},
Sec.~\ref{degas2in}). On subsequent
iterations, small changes in the source terms are accounted for
by allowing non-unit weighting factors at each source segment.
Once the accumulated changes exceed a certain size (specified by
parameters in the sources class), the DEGAS 2 source sampling
arrays are re-initialized (i.e., to once again use unit weights)
and the {\tt oldsourcefile} is rewritten. The effectiveness of this
procedure is currently being evaluated and should be considered
as experimental.
\item[tallysetup] This code specifies the types of ``scores'' (the
essential output of the code), along with their
dependencies, which will be computed in the DEGAS 2 run. Although the input
files(s) should not be modified by the novice user, this code needs to be
run to factor the values of physics and geometry parameters for the
problem at hand into the specification of these scores.
\item[flighttest] The core of DEGAS 2. This program launches, follows,
and scores the Monte Carlo trajectories which make up a DEGAS 2 simulation.
(The name is a holdover from early versions in which this did tracking 
with a minimum of hardwired physics).
\item[outputbrowser] Reads the output netCDF file generated by
{\tt flighttest} (in addition to all other netCDF files for the
problem) and allows the user to interactively browse its contents.
A script facility is provided for batch-like operation. 
\item[geomtesta] Another inappropriately named routine which serves as
the principal post-processing code for DEGAS 2. Although originally
designed as a diagnostic for the geometry, it has proven too easy to 
continue extending this code to warrant development of an honest
post-processing tool. The DEGAS 2 input and output files are read in by
the code. Using the geometry information, the zone-based plasma and neutral
data are transcribed onto a high density uniform rectangular mesh and
dumped into Silo or HDF files suitable for further manipulation by 
external graphics packages such as VisIt or IDL.
\item[ucd\_plot] Alternative post-processing code that uses DEGAS 2's
own mesh for representing output data.  Creates a Silo format file
for visualization with VisIt.
\item[degas2\_xgc.a] This library contains the core routines of DEGAS 2
and an interface for a subroutine based coupling to a plasma code.
The primary quantities exchanged between the codes are moments of the
plasma and neutral distribution functions.  An interface into the
DEGAS 2 atomic physics routines is also provided.  This library was
designed for and is used with the XGC0 guiding center neoclassical
particle transport code as part of the Center for Plasma Edge
Simulation \cite{Chan2008}.
\end{description}

\subsection{Other Setup Routines}
\label{othersetup}
\index{datasetup}
\index{ratecalc}
\index{reactionwrite}
\index{pmiwrite}
\index{boxgen}
\index{efit2dg2d}
\index{adaswrite}

These are occasionally used to add new atomic
or surface physics reactions, or to generate new atomic and surface physics
data files.

\begin{description}
\item[datasetup] This routine possibly belongs in the first list. It
reads text files describing the complete lists of elements, species, reactions,
materials, and plasma-material interactions available to DEGAS 2 and generates
the corresponding DEGAS 2 netCDF files. These lists are occasionally referred 
to in the code as the "reference" lists; the input file for problemsetup
specifies subsets of these lists. A new reaction is added to DEGAS 2 by 
inserting the requisite information into the ``reactionfile'' 
and rerunning datasetup. The new reaction can
then be added to the problem input file.
\item[ratecalc] A routine for computing averages of atomic physics
cross sections. Since it is designed to read data from the Aladdin 
database, ratecalc is presently limited in its flexibility. The principal
application thus far has been the generation of reaction rates and higher
moments for charge exchange and other interactions between atoms, molecules,
and ions.
\item[reactionwrite] A simple routine designed to translate collisional
radiative data for the electron impact ionization of hydrogen from the
format used by the old DEGAS code into netCDF files (ionization and
recombination) for DEGAS 2. This code has been run a handful of times
as needed to keep up with changes in the content and format of the atomic
physics data files.
\item[pmiwrite] Is analogous to reactionwrite, but handles almost all
of the plasma-material interaction data. Presently, two data files from
the old DEGAS code and one from EIRENE are used as input. A few processes
describable without external data are also included. This code can be used,
in a somewhat clumsy manner, to add new plasma-material interaction data
to DEGAS 2 using either one of these three existing data bases or via
a new datafile (with the rest of the code serving as an example).
\item[adaswrite] Reads ADAS {\tt adf11} \cite{ADAS} 
formatted files containing generalized
collisional radiative rates, arranges the data to conform to DEGAS 2's 
{\tt xsection} format, and writes them out as netCDF files.
\item[boxgen] Is the third of geometry / background generation tool provided 
with DEGAS 2.  It is intended to serve as an example of how to generate
geometry and background files from scratch. As the name implies, {\tt boxgen}
sets up a simple box geometry with a linearly varying plasma. However, 
no input file for {\tt boxgen} has been developed; all modifications have to
be effected through direct source code modification. The typical user
in search of a simple linear geometry to study can be overwhelmed by the
complexity of the code encountered there. For this reason, the 
{\tt definegeometry2d} / {\tt defineback} pair may be more suitable for
setting up simple problems.
\item[efit2dg2d] Is effectively a pre-processor for {\tt definegeometry2d}
used for problems such as simulating gas puff imaging in the main chamber
of a tokamak in which the spatial extent in the $R$-$Z$ plane can be
confined to a rectangular region with no intervening material surfaces
or X-points.  An EQDSK ``g'' file is read in and then used to trace out
flux surface contours in the rectangle.  These are transformed into 
``polygons'', and a {\tt definegeometry2d}-compatible representation
of them is produced on output.  Note, however, that additional effort
is required by the user to assemble the {\tt definegeometry2d} input
file.
\item[tri\_to\_sonnet] Transforms an unstructured mesh specified in
the {\tt Triangle} format into a Sonnet-format mesh file suitable for
use in {\tt definegeometry2d}.
\item[poly\_to\_tri] Reads the ``plasma'' and ``vacuum''
triangles and rectangles from a polygon netCDF file (as produced by
{\tt definegeometry2d}), splits the rectangles into two triangles,
and writes them all out as {\tt Triangle} element and node files.
Zone index and number data are transferred as well.
\end{description}

\subsection{Test Routines}
\label{testroutines}
\index{reactiontest}
\index{pmitest}
\index{sourcetest}
\index{dataexam}
\index{sysdeptest}
\index{randomtest}

\begin{description}
\item[reactiontest] Can be run after {\tt problemsetup} 
to check the computation
of the reaction rate, and collision products (including velocities and
scoring data) for a given reaction.
\item[pmitest] The analog of {\tt reactiontest} for plasma-material
interactions. Again, this code can be executed for a given problem once
{\tt problemsetup} has been run. Since the velocity distributions of some
plasma-material interactions are described statistically, pmitest
is also set up to average the outcome of a given interaction over a
specified number of trials and print out a list of the product velocities
for external manipulation.
\item[sourcetest] Tests the distribution in phase space of the source
terms for a particular problem. Since the source is described in the 
``background'' netCDF file, {\tt defineback}, 
{\tt readbackground}, {\tt boxgen}, or something equivalent
must be run prior to using this routine. Like {\tt pmitest}, the code is able
to run many trials, compute averages, and printout data from each 
individual trial.
\item[dataexam] Provides the user with a convenient interface to the
atomic physics data files. Although netCDF files can be converted to
a human-readable text form, interpreting those data is since DEGAS 2
collapses the multi-dimensional structure components of the atomic physics data
into a single one-dimensional array. This routine reads the ``raw'' (at
the ``reference'' level) data files prior to their insertion into the
problem netCDF file and, thus, can be run at any time. An analogous
routine for the plasma material interactions is needed.
\item[sysdeptest] Tests several system-dependent features of DEGAS 2.
This would be run only if the code was being ported to a new operating
system or if the operating system on a currently supported architecture
underwent a signficant upgrade.
\item[randomtest] Is used to verify that DEGAS 2 will give the same 
results on 
different architectures (provided the random number seeds are set the
same!). 
\end{description}

\subsection{Miscellaneous Routines}
\label{miscroutines}
\index{matchout}
\index{matcheir}
\index{datamatch}

\begin{description}
\item[matchout] Is used to compare the current DEGAS 2 output netCDF
file with another from an equivalent run (i.e., all of the arrays must
be exactly the same size) on the command line. The most frequent use
of {\tt matchout} is to verify that two separate runs have yielded 
results which are the same to within roundoff error. Because of the size of 
the ouptut netCDF file, verifying this by visually comparing the files 
is impractical.  
\item[matcheir] Is specifically designed for comparing with output from
the EIRENE code. The name of the EIRENE file is specified on the command
line. See Sec.~\ref{outofthebox}.
\item[datamatch] Provides a means for comparing the contents of the two
netCDF data files specified on the command line.
\end{description} 

\subsection{DG and Carre}
\label{dgcarre}
\index{DG}
\index{Carre}

The graphical interface DG code is used here
in conjunction with {\tt definegeometry2d}, as
in B2-Eirene / SOLPS (its original application), to simplify specification of 
2-D divertor tokamak geometries .  Since DG is also available
separately from DEGAS 2, its manual, {\tt dg.pdf} is contained in the
\verb+dg_carre/DG/DOC+ directory.  Carre\cite{Marc96} generates 2-D
quasi-orthogonal meshes of the sort used by B2 and UEDGE. This
version of Carre includes pre- and post-processing scripts that
simplify its use with DG.  A number of utility programs for 
manipulating magnetic equilibrium files also come with this
DG distribution; see the DG manual for details. 
Note, however, that these versions of DG and Carre have not
been kept up to date with those distributed with SOLPS-ITER.
Users requiring additional capabilities not available in
the DEGAS 2 versions may wish to contact the SOLPS-ITER
development team.

\section{Input Files}

The essential inputs to a Monte Carlo neutral transport code are:

\begin{itemize}
\item The identity of the ``test'' species (usually neutral in plasma problems)
that are of interest,
\item The identity and phase space distributions of the ``background''
(typically electrically charged in plasma problems)
species with which
the test species will interact, 
\item The interaction processes between the test and background species. 
The relative probability of these processes as a function of the 
test and local background properties are required as well as a 
prescription for specifying the outcome of each interaction.
\item The simulation geometry, including a framework for
holding the background species data, as are a means for tracking
the test species through the background.
\item The boundaries of the geometry, including a prescription
for specifying the outcome of interactions between the test species and
these boundaries.
\end{itemize}

DEGAS 2 utilizes several text based input files for controlling the actions
of its component executables (see Sec.~\ref{components}). 
In this section we will discuss a few of these input files. The more complex
input files are described in their respective codes; links to the documentation
for those files will be given instead. 
More
detailed information on running DEGAS 2 will follow in later sections.

These text input files are, for the most part, 
read with the same set of text processing utilities that adhere to the following
guidelines. In some cases, particularly
those involving files generated by other codes, the files are read directly
with formatted FORTRAN read statements.

\begin{enumerate}
  \item Spacing and blank lines are ignored. The user is free to utilize white
space in whatever way facilitates reading and maintaining the input file.
  \item Comments, either a complete line or at the end of a line, are begun
with a \verb+#+ sign.
  \item File names can be a relative or full path name.
\end{enumerate}

\subsection{degas2.in}
\label{degas2in}
\index{symbolic name}

{\tt degas2.in} is the input file which controls (most of) the other
input files. Each line consists of DEGAS 2's symbolic name for the
file (this can be changed only by modifying {\tt readfilenames.hweb}
and {\tt readfilenames.web}) followed by the path name for the file
to be used in the current problem. The order of the lines in this file
is unimportant. 

\begin{verbatim}
elementsfile  ../data/elements.nc 
backgroundfile bk_uers.nc 
geometryfile ge_uers.nc 
problemfile pr_uers.nc
reactionfile reactions.nc
speciesfile ../data/species.nc 
aladinfile ../data/aladinput.nc
aladoutfile ../data/aladoutput.nc 
elements_infile ../data/elements.input
problem_infile  pr_uers.input
reaction_infile reactions.input
species_infile ../data/species.input
materials_infile ../data/materials.input
materialsfile ../data/materials.nc
pmi_infile ../data/pmi.input
pmifile ../data/pmi.nc
cramdproblemfile cramdprob.nc
tallyfile tally_uers.nc
outputfile degas2_uers_out.nc
oldsourcefile os_uers.nc
snapshotfile sn_uers.nc
\end{verbatim}

Note that some of these are 
outdated or rarely used (unused entries may be deleted from the file, if
desired);
see the \href{classes.pdf#page.129}{readfilenames class} for
more specific information. Most of the files come in pairs with the name
\verb+XXX_infile+ corresponding to a text input file and
\verb+XXXfile+ being a netCDF file generated by a program
like {\tt datasetup}, {\tt problemsetup}, etc.

\subsection{elements\_infile}
\label{elementsinfile}
\index{elements\_infile}
\index{elementsfile}

{\tt elements\_infile} lists all of the elements available to DEGAS 2 for
use in constructing the species (see Sec.~\ref{speciesinfile}). 
Additional detail is given at the beginning of the 
\href{elementsetup.pdf#page.1}{file} {\tt elementsetup.web}.

\subsection{species\_infile}
\label{speciesinfile}
\index{species\_infile}
\index{species}
\index{generic}
\index{speciesfile}

Inputs for species are contained in the file with the symbolic
name {\tt species\_infile}. Additional detail is given at the beginning
of the \href{speciesetup.pdf#page.1}{file} {\tt speciesetup.web}.

\subsection{reaction\_infile}
\label{reactioninfile}
\index{reaction\_infile}
\index{reaction type}
\index{chargex}
\index{elastic}
\index{dissoc}
\index{dissoc\_rec}
\index{ionize}
\index{ionize\_suppress}
\index{recombination}
\index{reactionfile}
\index{generic reactions}
\index{non-generic reactions}

The file {\tt reaction\_infile} describes all of the reactions available in
DEGAS 2. Additional detail is given at the beginning of the
\href{reactionsetup.pdf#page.1}{file} {\tt reactionsetup.web}.

Note that adding a new reaction to this file involves two tasks beyond 
inserting the appropriate lines in {\tt reaction\_infile}. First, the 
netCDF file for the atomic physics data must be generated (see
Sec.~\ref{addreactions}). The
second task would be to write subroutines for setting up the products and
handling collisions (see Sec.~\ref{reactionprocessing}). 
This would be necessary only if a new reaction type
were being added.

\subsection{ratecalc Input}
\label{ratecalcinput}

The {\tt ratecalc} code is used to compute atomic physics reaction
rates from cross section data, as well as a few other tasks. 
The code is documented internally. Follow this
% To get this to work, had to add some stuff to fweb.sty (the hyperref
% package). Can then do: pdflatex ratecalc.tex. 
% Note that this page
% number is the one TeX assigns; not the one acroread uses.
% pdflatex will complain occasionally about ``non-pdf special'', but
% I don't see any deleterious effects so far (appears to have something
% to do with section headers).
\href{ratecalc.pdf#page.1}{link} to
the introductory section to learn more.


\subsection{materials\_infile}
\label{materialsinfile}
\index{materials\_infile}
\index{materialsfile}

The materials described in {\tt materials\_infile} are essentially just labels
which will be used in conjunction with the plasma material interactions 
(see Sec.~\ref{pmiinfile}) to specify how test 
species interact with non-transparent
surfaces. Additional detail is given at the beginning of the 
\href{materialsetup.pdf#page.1}{file} {\tt materialsetup.web}.

\subsection{pmi\_infile}
\label{pmiinfile}
\index{pmi\_infile}
\index{PMI}
\index{generic PMI}
\index{non-generic PMI}
\index{reflection}
\index{adsorption}
\index{desorption}

The file {\tt pmi\_infile} describes all of the plasma-material interactions
(PMI) available in DEGAS 2. Additional detail is given at the beginning
of the \href{pmisetup.pdf#page.1}{file} {\tt pmisetup.web}.

\subsection{problem\_infile}
\label{probleminfile}
\index{reference level}
\index{subset data}
\index{test species}
\index{background species}
\index{problem\_infile}
\index{problemfile}

As noted already in the section describing the DEGAS 2 components
(see Sec.~\ref{components}), there are two levels of physics input to DEGAS 2.
The input files noted thus far comprise the reference level: in
practice, the sum total of the data available to the code (although
in principle it could be smaller). The second level is the problem level
of data. This prescribes the species, reactions, materials, and
PMI which will serve as the physical
model to be used in carrying the simulation at hand. In some
parts of the (internal) code, these are also referred to as the
subset data since they represent a subset of the reference data
(see also Sec.~\ref{addreactions} and Fig.~\ref{classeschart}).

Additional concepts alluded to above in connection with the
reaction input file (see Sec.~\ref{reactioninfile}) are those of 
test species
and background species. Most simply, test species are the ones
DEGAS 2 will track as they collide off of background species.
The use of the word species here is important: both of the test
and background lists are subsets of the ``species'' list 
(see Sec.~\ref{speciesinfile}). One more precise distinction between the
two species types is that we assume that we know the distribution
function (in space and velocity) of the background species; in
fact, such information is required input to DEGAS 2. On the other
hand, we are attempting to {\em compute} the test species
distribution function, moments of which serve as the primary output
of DEGAS 2.

More information about the input file can be found at the beginning of
the \href{problemsetup.pdf#page.1}{file} {\tt problemsetup.web}.

\subsection{readgeometry Input}
\index{readgeometry input}

% Again, generated this using FWEB to get the TeX file and
% pdflatex to get the PDF file.

The documentation for {\tt readgeometry} input is maintained in the
source code, {\tt readgeometry.web}. For readers of the PDF version
of this manual, here's a 
\href{readgeometry.pdf#page.1}{link} to the
corresponding PDF file.

\subsection{definegeometry2d Input}
\index{definegeometry2d input}

The documentation for {\tt definegeometry2d} input is maintained in the
source code, {\tt definegeometry2d.web}. For readers of the PDF version
of this manual, here's a 
\href{definegeometry2d.pdf#page.1}{link} to the
corresponding PDF file.

\subsection{defineback Input}
\index{defineback input}

The documentation for {\tt defineback} input is maintained in the
source code, {\tt defineback.web}. For readers of the PDF version
of this manual, here's a 
\href{defineback.pdf#page.1}{link} to the
corresponding PDF file.

\subsection{readbackground Input}
\index{readbackground input}

Currently, the only possible inputs to {\tt readbackground} are the
{\tt UEDGE} (a specifically formatted text file generated during
{\tt UEDGE} post-processing) and {\tt DEGAS} formats (an input file for
the old DEGAS code). This routine makes some specific assumptions about
the contents of these files. Since this situation is unsatisfactory from
a number of viewpoints, a better long-term approach is being contemplated.
For that reason, no additional documentation is provided at this point.

\subsection{tally\_infile}
\label{tallyinfile}
\index{tally\_infile}
\index{tallyfile}
\index{tallysetup}
\index{tally}

The documentation on the input file for {\tt tallysetup} is maintained
in the source code, {\tt tallysetup.web}. This 
\href{tallysetup.pdf#page.1}{link} will take the reader to the corresponding
PDF file. The \href{classes.pdf#page.52}{tally class} contains more 
extensive and detailed information.

\subsection{outputbrowser Input}
\index{outputscript}

This post-processing utility can be run interactively or with an
input file. When run interactively, the session is logged into a
script file called {\tt outputscript} (in the run directory) which
can be subsequently used as input to {\tt outputbrowser}. Details on
the use of {\tt outputbrowser} and on the format of the script
file can be found in the introduction to the 
\href{outputbrowser.pdf#page.1}{file} {\tt outputbrowser.web}.

\subsection{geomtesta Input}
\index{geometry.inp}

The documentation for {\tt geomtesta} is maintained in
the source code, {\tt geomtesta.web}.  For readers of the
PDF version of the manual, here's a 
\href{geomtesta.pdf#page.1}{link} to the corresponding PDF file.

\subsection{ucd\_plot Input}

The documentation for {\tt ucd\_plot} is maintained in
the source code, {\tt ucd\_plot.web}.  For readers of the
PDF version of the manual, here's a 
\href{ucd_plot.pdf#page.1}{link} to the corresponding PDF file.

\section{Compiling DEGAS 2}

DEGAS 2 is compiled with \href{http://www.cmake.org}{CMake}, a widely used standard for building portable software and managing dependencies.
To configure: 
\begin{verbatim}
cd ~/[your-path-to]/degas2
mkdir build
cd build
cmake ..
\end{verbatim}
Components are built as before from the build directory:
\begin{verbatim}
make randomtest
../bin/randomtest
\end{verbatim}
The resulting executables are output to the bin directory. 
For flexible use of these tools, it is recommended to add the bin directory 
to your shell's \$PATH environment variable.

At CMake configure time, several things happen beyond generating the Makefile. 
Most importantly, FWEB is compiled from the source distributed with DEGAS 2. 
That way, {\tt ftangle} is available to CMake to build the fortran source files.
Silo and Triangle are likewise also distributed with DEGAS 2, but CMake will instead use an existing implementation of Silo on the system if it is found.
Template input files for the data directory are also generated with the correct global path to the degas2/data directory.

%By default, executables are compiled in ``Debug'' mode. To gain more performance, they can be compiled as ``Release''s as follows:
%\begin{verbatim}
%cd [your-path-to]/degas2/build
%cmake -DCMAKE_BUILD_TYPE=Release ..
%make flighttest
%\end{verbatim}
MPI is enabled by default if CMake can find the MPI library.

Some cautionary remarks:
\begin{itemize}
   \item If the use-defined source files \verb+usr2ddetector.web+ or \verb+usr2dplasma.web+ are being used, the CMake configure step needs to be rerun so that these files get recognized as part of the compilation.
   \item Currently, CMake does not track changes to .hweb files. If these are edited and changes pulled from the repository, CMake will not recognize this. It is therefore recommended practice to {\tt make clean} when pulling updates.
   \item The file CMakeLists.txt contains a section (clearly labelled) which is automatically generated to track dependencies. If you make changes to the code such that dependencies change, be sure to include these dependencies in {\tt src/Makefile} as before, run {\tt make depend} and copy the contents of {\tt cmakecommands.txt} to the relevant section of CMakeLists.txt. This extra step is temporary until the legacy build system is no longer supported.
   \item The CMake implementation is intended to reproduce all useful functionality of the legacy build system. Please reach out the developers \ref{help} if it does not suit your needs.
\end{itemize}

\subsection{Legacy build instructions}

This subsection describes how to compile DEGAS 2 with the included Makefile.
This method is still supported, but will likely be phased out. 
If compiling this way, replace the {\tt build} subdirectory with your {\tt SYSTEM} name in the examples \ref{examples}.

The Makefile has been designed to compile on  a variety of architectures
and operating systems. It uses the shell command {\tt uname} to
determine the operating system upon which {\tt gmake} is being run.  The root
name of the current directory (see Sec.~\ref{basics}) tells it the
operating system for which the code is being compiled.  In most cases,
these are the same; a mechanism for cross compiling (i.e., target is a
different operating system) is provided (Sec.~\ref{cross}).
The systems used thus far are: 

\begin{center}
\begin{tabular}{lp{10 em}l}
Generic Name & {\tt uname} & Run Directory Root \\ \hline
Linux 32-bit & {\tt Linux} & {\tt LINUX} \\
Linux 64-bit & {\tt Linux}, {\tt uname -m} contains {\tt 64} & {\tt
  LINUX64} \\
Sun Solaris & {\tt SunOS} & {\tt SUN} \\
IBM AIX & {\tt AIX} & {\tt IBM} \\
Digital OSF1 & {\tt OSF1} & {\tt ALPHA} \\
Silicon Graphics IRIX & {\tt IRIX} or {\tt IRIX64} & {\tt SGI} \\
Mac OS X & {\tt Darwin} & {\tt MACOSX} \\
\end{tabular}
\end{center}

Note that due to the proliferation of relatively cheap and fast
Linux clusters, usage of DEGAS 2 on other systems has essentially
ceased and support for those systems cannot be guaranteed.
The most frequently used compiler on Linux systems
is Portland Group F90 compiler ({\tt PGROUP} in
the {\tt Makefile}), albeit with F77 format code 
({\tt FORTRAN90=no}).  The Pathscale ({\tt PATHSCALE} in 
the {\tt Makefile}) and {\tt gfortran} (version 4.3) have been
successfully tested with the F90 formatted code (in
{\tt DEBUG} mode only). The use of
the F90 option eliminates the need for system dependent code
so that other compilers can be used without modifications to
the {\tt sysdep.hweb} or {\tt sysdep.web} files; only the
compiler
related variables in the {\tt Makefile} or {\tt Makefile.local}
need be defined. On the other hand, some F9x compilers (e.g.,
Lahey-Fujitsu F95) may result in code that runs very slowly.

By default, the Makefile assumes that the DEGAS 2 main directory
exists at \verb+$HOME/degas2+. If you unpacked the {\tt tar} file
in some other location, you'll need to tell the Makefile where to look
(see Sec.~\ref{adjustments}). 
Blindly change directory names is not recommended.
Some flexibility is again provided; see below for details.

\subsubsection{Basics}
\label{basics}

Let's start by compiling the {\tt randomtest} utility (see 
Sec.~\ref{testroutines}) on a Sun system. To use other systems,
replace {\tt SUN} with the appropriate directory root name (more than
one run directory can be used; see Sec.~\ref{adjustments}) from the above list.

First, change to the main DEGAS 2 directory and create the SUN 
subdirectory:
\begin{verbatim}
cd ~/degas2
mkdir SUN
cd SUN
cp ../src/Makefile .
\end{verbatim}

Once the Makefile is present here, it will ``update itself'' with respect
to the copy in the {\tt src} directory when needed. To get an idea of
how {\tt ftangle} works, just ``make'' the main FORTRAN file for 
{\tt randomtest}:
\begin{verbatim}
gmake randomtest.f
\end{verbatim}
At some point, you may want to visually compare this file with the
original source code {\tt randomtest.web} (in the {\tt src} directory, with
the other source files) to gain an appreciation
for the amount of work the FWEB macros do in DEGAS 2. For more information
on these macros, see the documentation in {\tt array.hweb} and the
other header files.

Now, finish making {\tt randomtest}:
\begin{verbatim}
gmake randomtest
\end{verbatim}

To run {\tt randomtest}, just type
\begin{verbatim}
./randomtest
\end{verbatim}
The output to the right of the ``='' sign should match the numbers in
parentheses.

\subsection{Making Documents}

Generating ``woven'' (i.e., using FWEB's {\tt fweave} utility)  
documentation is similar:
\begin{verbatim}
cd ~/degas2
mkdir tex
cd tex
cp ../src/Makefile .
gmake randomtest.dvi
\end{verbatim}
The DVI file can then be printed or viewed (using {\tt xdvi}), as desired.
In fact, the Makefile also provides additional ``targets'', e.g.,
{\tt randomtest.print} (uses {\tt lpr} to print to your default printer),
{\tt randomtest.view} (launches {\tt xdvi}) for you, and
{\tt randomtest.ps} (a PostScript file generated via {\tt dvips}). You
can generate woven documents for all of the ``{\tt .web}'' source files.
For some of the more useful ones, see Sec.~\ref{otherdoc}.

You could also just generate the \TeX\ file (again, from the \verb+degas2\tex+
directory),
\begin{verbatim}
gmake randomtest.tex
\end{verbatim}
You could then use {\tt latex}, {\tt pdflatex}, or whatever other \TeX\ application
you wished. 

\subsection{Adjustments to Makefile}
\label{adjustments}

A number of the default settings in the Makefile can be overridden by
placing the desired values in a file called {\tt Makefile.local} in the
working (e.g., {\tt SUN}) directory. Some of the more frequently
changed variables are:
\begin{itemize}
  \item {\tt DEBUG=no} turns on optimization. The default {\tt DEBUG=yes}
is required for debugging.
  \item {\tt DEGASROOT=somedir} will tell Makefile to find the main
DEGAS 2 directory at {\tt somedir/degas2}. {\tt somedir} should be an
absolute path name. E.g.,
\begin{verbatim}
DEGASROOT = /u/somedir
\end{verbatim}
  \item {\tt FORTRAN90=yes} switches from using the default FORTRAN 77 
compiler to FORTRAN 90. Be aware that only some of the FORTRAN 90
compilers available work satisfactorily. Some will compile the code,
but run slowly.
  \item {\tt FCF77}, {\tt FCF90} tell {\tt Makefile} which ``normal'' (non-MPI)
compiler to use.
  \item {\tt MPI=yes} will enable the MPI commands in the source
code and direct the {\tt Makefile} to use the MPI compiler flags.
  \item {\tt FCMPI77}, {\tt FCMPI90} tell Makefile which MPI 
compliant compiler to use.
  \item {\tt NETCDFV2=yes} tells DEGAS 2 that your system has only the
older version 2 of the netCDF library (by default, DEGAS 2 expects version 3).
  \item {\tt GRAPH\_FILE} can be set to either {\tt HDF4},
{\tt SILO}, or {\tt SILO\_HDF5} to set the type of graphics
output produced by {\tt geomtesta} and {\tt definegeometry2d}, 
according to the availability and capabilities of the Silo library.
Note that this flag is also used in compiling {\tt ucd\_plot},
but that the code requires one of the two Silo options.  The
{\tt postdetector} code does not use this flag since it produces
only HDF4 files.
  \item The Makefile will accomodate more than one working directory on
a given machine provided the additional directory has a name like
{\tt SUN-junk}, where ``junk'' is some string meaningful to you.
  \item If you would like to have more than one source directory, we
recommend creating a new {\tt degas2} directory heirarchy somewhere else
in your file system and use the {\tt DEGASROOT} variable to get 
the Makefile working there (you can set up symbolic links to common
directories such as {\tt data}).
\end{itemize}


\subsubsection{Cross-Compiling}
\label{cross}

{\bf The mechanism for cross-compilation has been updated to use
{\tt ssh} and {\tt scp} and no longer relies on the presence
of AFS.  However, this section of the
manual has not been correspondingly updated.
Please contact the code authors if you need this capability.}

As an example of how to cross-compile, here are the procedures for 
making {\tt randomtest} on the NERSC Crays. This assumes that you
have access to AFS from your local workstation (and at NERSC, of
course):

\begin{verbatim}
mkdir <some-afs-directory>/degas2/CRAY
ln -s <same-afs-directory>/degas2/CRAY ~/degas2/CRAY
cd ~/degas2/CRAY
cp ~/degas2/src/Makefile .
gmake randomtest
\end{verbatim}

This will use {\tt ftangle} on the local workstation to generate
the FORTRAN source code files from the FWEB files in the \verb+src+
directory. It is advisable to set up a {\tt Makefile.local} 
(see Sec.~\ref{adjustments}) in this
directory that contains the same flags that will be used on the
remote machine, particularly {\tt DEBUG}, {\tt FORTRAN90}
and {\tt MPI}.

Then, log on to the remote machine and do:
\begin{verbatim}
ln -s <same-afs-directory>/degas2 ~/degas2
cd <some-work-directory>/CRAY
cp ~/degas2/CRAY/Makefile .
make randomtest
\end{verbatim}

In addition to the usual variables in {\tt Makefile.local}, one
needs to add {\tt STANDALONE=no}. This will tell {\tt make} that 
the FORTRAN (e.g., {\tt .f}) files have already been generated on
your local workstation, and that it can find them in
\verb+~/degas2/CRAY+ on the remote machine, which should be
linked back to your local workstation via AFS.

The \verb+degas2/data+ directory will likely also be needed on the 
remote machine. You can manually copy the directory and its contents
to the remote machine (e.g., via ftp) 
or keep a copy in the AFS directory and use 
a link to allow the code on the remote machine to find the files.
E.g., for the second option, starting on the local workstation:
\begin{verbatim}
cd ~/degas2
cp -R data <same-afs-directory>/degas2
\end{verbatim}
And on the remote machine
\begin{verbatim}
cd <same-work-directory>
ln -s <same-afs-directory>/degas2/data data
\end{verbatim}
Be aware that in either case, the \verb+data+ directory could eventually
get out of sync with the one on your local workstation. If this
is a likely possibility, one could consider also linking the 
\verb+data+ directory on the local workstation to the one in the
AFS directory, making the latter the only real copy.

\subsubsection{Miscellaneous Targets}

The Makefile provides some other occasionally useful targets. Please
see the Makefile for more details.

\begin{itemize}
  \item {\tt clean} removes source, object, and other unessential
files. This is pretty thorough and indescriminate; use with caution.
  \item {\tt TAGS} updates the Emacs tags table (see Sec.~\ref{gnu}).
You may need to do this if you subtstantially altered any of the source
files or have other reason to believe that the tags table is out-of-date.
You would run \verb+gmake TAGS+ from the {\tt src} directory only.
  \item {\tt depend} updates the Makefile dependencies 
(in {\tt Makefile.depends}). Do this if the header file dependencies
of one or more executables have been changed or if you believe that the
dependencies are out of date. Note that source code (i.e., object files)
dependencies are
explicitly stated in the Makefile and must be updated by hand.
  \item {\tt foof} dumps out values of some of the internal Makefile
variables. This is useful for debugging Makefile problems.
\end{itemize}

% \subsection{CVS Access to DEGAS 2}

% Properly privileged (in the UNIX sense) authors of DEGAS 2 will be given
% access to the
% CVS (see Sec.~\ref{cvs}) repository for the code. With this access,
% one can obtain the latest version of the source with:

% \begin{verbatim}
% setenv CVSROOT /afs/pppl.gov/common/cvs
% cd
% cvs co degas2
% \end{verbatim}

% (this will place the code in \verb+~/degas2+; see also Sec.~\ref{adjustments}).
% Of course, the command setting the {\tt CVSROOT} variable can be more
% conveniently executed from a login (e.g., \verb+~/.login+) or shell startup
% script (e.g., \verb+~/.cshrc+). The user will also have access to
% all older versions of the source and data files through the appropriate
% CVS commands. More importantly, the user will be able to {\em update}
% the CVS repository. For this reason, such access is controlled by 
% the primary authors of DEGAS 2. Contact them if you feel that you can
% contribute to the DEGAS 2 source code
% (see Sec.~\ref{help}).

\section{Examples}
\label{examples}

The {\tt examples} directory contains several documented DEGAS 2 examples,
with input output, and auxiliary files. They are included not only
for instructive purposes, but also for verifying that successive versions
of the code give the same results. Namely, the user should be able to
download the code and generate exactly (to within roundoff error) the
results contained within each of the example subdirectories. Any
discrepancies should be reported to the primary code authors (see
Sec.~\ref{help}).

\subsection{Analytic\_fluid\_bench}
\label{analyticfluidbench}

This simple run demonstrates that DEGAS 2 matches the results of an analytic
model in the fluid limit. More information on this comparison can be
found in the IAEA proceedings\cite{Stot96} or on the Web as a 
\href{http://bp.pppl.gov/pub_report/1997/PPPL-3221.ps}{PostScript}
or 
\href{http://bp.pppl.gov/pub_report/1997/PPPL-3221.pdf}{PDF} file.

To run this example, follow these steps:
\begin{enumerate}
  \item {\bf Set source code switch.} The source code file {\tt boxgen.web}
has within it a few different settings available which can be selected
by changing the value of the flag {\tt BOXRUN}. With this file open in
an editor, change the line
\begin{verbatim}
@m BOXRUN 0
\end{verbatim}
so that it now reads
\begin{verbatim}
@m BOXRUN 32
\end{verbatim}
Save the file and exit the editor. 
  \item {\bf Switch to the example directory.} If the main DEGAS 2 
directory sits within your home directory:
\begin{verbatim}
cd ~/degas2/examples/Analytic_fluid_bench
\end{verbatim}
  \item {\bf Copy the} {\tt degas2.in} {\bf file into your working
directory.}
\begin{verbatim}
mkdir ../../workdir
cp degas2_boxgen.in ../../workdir/degas2.in
cd ../../build
\end{verbatim}
(overwriting any {\tt degas2.in} file you may have had there already!).
  \item {\bf Prepare reference data.} This example uses the standard
reference data files in the {\tt data} directory. Hence, this step can
be skipped if
those files have not been altered since you downloaded the code.
\begin{verbatim}
make datasetup
cd ../workdir
../bin/datasetup
\end{verbatim}
  \item {\bf Prepare problem data.} This and subsequent steps are 
{\em not} optional! Note that the {\tt degas2.in} file points back to
the {\tt examples} directory for the problem input file (see
Sec.~\ref{probleminfile}).
\begin{verbatim}
cd ../build
make problemsetup
cd ../workdir
../bin/problemsetup
\end{verbatim}
  \item {\bf Generate geometry and background files.} The {\tt boxgen}
program takes care of both of these:
\begin{verbatim}
cd ../build
make boxgen
cd ../workdir
../bin/boxgen
\end{verbatim}
  \item {\bf Set the number of flights.} The background file {\tt bk\_boxgen.nc}
contains the specification of the number of flights to be run. Open this
file in an editor (Emacs will be the most convenient) and find the line
\begin{verbatim}
 source_num_flights = 100 ;
\end{verbatim}
Change the ``100'' to either ``1600'' or ``6400''. The example contains the
output for both cases. Save the file and exit the editor.
  \item {\bf Set up the tallies.} 
\begin{verbatim}
cd ../build
make tallysetup
cd ../workdir
../bin/tallysetup
\end{verbatim}
  \item {\bf Run the code.}
\begin{verbatim}
cd ../build
make flighttest
cd ../workdir
../bin/flighttest
\end{verbatim}
  \item {\bf Check the text output.} The file {\tt density.out} can be directly
compared with either {\tt density\_1600.out} or {\tt density\_6400.out},
depending on the number of flights. The numbers in here have only five
digits of precision. If an exact match is not obtained, something is
amiss. Please contact the code authors (see Sec.~\ref{help}) if you feel
that a problem exists with the code as you downloaded it. 
  \item {\bf Check the binary output.} You can use the {\tt matchout}
utility (see Sec.~\ref{miscroutines}) to do this. 
  \item {\bf Compare with the analytic solution.} The file {\tt soln\_32}
contains the columns:
\begin{description}
  \item[x] Distance along the problem space in meters.
  \item[N1(x)/N(0)] Relative density (to density at $x=0$) predicted by
one analytic model (see the above references) solution.
  \item[N2(x)/N(0)] Relative density predicted by a second analytic model
solution.
  \item[Ti(x)] Ion temperature profile.
  \item[Ni(x)] Ion density profile.
\end{description}
The columns of the {\tt density.out} file contain
\begin{enumerate}
  \item First zone number. These zones correspond directly to the $x$
values in the {\tt soln\_32} file.
  \item Second zone number. This is always 0 since this is a 1-D problem.
  \item Neutral hydrogen density in m$^{-3}$. To normalize this column,
divide the whole column by the value in the first row. Since this corresponds
to the first value of $x$ in {\tt soln\_32} and {\em not} to $x = 0$, 
multiply the whole column by the relative density from the analytic
solution at the first value of $x$, normalizing the DEGAS 2 result
to the analytic solution at one point.
  \item Relative standard deviation of the neutral density (see
Sec.~\ref{statcompare}).
  \item Neutral hydrogen pressure in pascals.
  \item Relative standard deviation of the neutral pressure.
  \item The last two columns are no longer used and should be filled with
zeroes.
\end{enumerate}
\end{enumerate}

\subsection{Neutral-Neutral Scattering Examples}

The initial implementation of the BGK algorithm for handling neutral-neutral
collisions (see Sec.~\ref{NeutralNeutral}) in EIRENE 
has been tested by comparison
with analytic expressions for the time dependence of the relaxation of 
a distribution function (both in self-collision and
mixed-species collision systems) and against models describing the
Couette flow problem.\cite{Mayxx} The relaxation benchmarks have been
repeated with DEGAS 2, but have not been formally made into an example
run and will not be discussed here. 
Further details can be provided upon request. However, the Couette flow
problem can be set up with only a few minor modifications to the code and
does make a suitable example.

\subsubsection{Couette Flow}
\label{couetteflow}

The Couette flow problem of fluid mechanics involves the flow of fluid between
two parallel, sliding plates. The fluid is assumed to have ``no slip''
boundary condition at the plates. The viscosity of the fluid drags
along adjacent fluid elements, resulting in a velocity gradient
between the boundary velocities represented by the two plates.

In this example, two semi-infinite plates are a distance $d$ apart 
in the $x$ direction. 
Periodic boundary conditions are enforced in the $y$ and $z$ directions so that
the problem is effectively infinite in those directions.  Particles
are initialized at the $x = 0$ plate with a thermally distributed
velocity to which a constant velocity in the $y$ direction, 
$v_{y0}$,
is added; this represents the velocity of the plate. The second plate
at $x = d$ is treated as stationary. Particles striking it are assumed to be
re-emitted with a thermal distribution pointed in the $-x$ direction.
Particles striking the $x = 0$ plate are absorbed. Hence, in the no
collision limit, particles make only one round trip across the box.

The choice of the distributions used at the plates is crucial
to reproducing the analytic models. In particular, the ``Maxwell flux''
distribution \cite{Reit92} must be used. This distribution describes a 
recycling thermal flux
of particles moving in a particular direction (in this case, perpendicularly
away from the wall). This is the distribution used in some of the
PMI data files, such as in \verb+h2_des_maxw_mo+. In this example, however,
the distributions of the particles coming off of the two plates 
are enforced through hardwired code,
not through the data files. 

The initial velocity is chosen from the distribution
\begin{equation}
Q(\vec{v}) \propto v_{x}  
\exp \left\{ - \frac{m}{2 k T} 
[v_{x}^{2} + (v_{y} - v_{y0})^{2} + v_{z}^{2}] \right\}.
\label{maxflux}
\end{equation}
Note the plate velocity $v_{y0}$ in the exponent. For particles 
``reflected'' at the $x = d$ surface, $v_{y0} \equiv 0$. This is just a
thermal (Maxwellian) distribution multiplied by $v_{x}$, the velocity in the
direction normal to the surface and is required for describing sources that
simulate a recycling process. 

A brief digression will clarify the relationship between this {\em source} 
distribution and the thermal distribution that characterizes the 
particles in the volume of the fluid. Consider the number of particles
of arbitrary volume distribution $f$
that recycle at, say, the $x = 0$ surface in a time interval $\Delta t$
(a source distribution specifies a density in phase space per unit time).
For particles in the velocity interval $v_{x} \rightarrow v_{x} + dv_{x}$,
this is $f dv_{x} A \Delta x$, where $A$ is the area of the recycling surface
and $\Delta x = v_{x} \Delta t$ is the distance over which particles with
velocity $v_{x}$ can reach $x = 0$ in $\Delta t$. So, the phase space
density of the recycling particles per unit time (dividing this number
by $\Delta t$ and $dv_{x}$) is $\propto v_{x} f$. Another way of saying this
same thing is that the source needs to emphasize the faster particles
in the distribution since more of them will reach the recycling surface in
a given time interval. One upshot is that the average energy of the
source particles is $\langle E \rangle_{\rm source} = 2 T$ while the
average energy in the volume (for $v_{y0} = 0$) is the familiar $3/2 T$.

For clarity, we note that the volume distribution function in the
free-molecular limit is
\begin{eqnarray}
f(\vec{v}) & = & f_{+}(\vec{v}) + f_{-}(\vec{v}), \\
\label{fmax}
f_{+}(\vec{v}) & = & \frac{n}{2} 2 \left( \frac{m}{2 \pi T} \right)^{3/2}
\exp \left \{ - \frac{m}{2T} [v_{x}^{2} + (v_{y} - v_{y0})^{2} + v_{z}^{2}]
\right \}, \; v_{x} > 0 \\
f_{-}(\vec{v}) & = & \frac{n}{2} 2 \left( \frac{m}{2 \pi T} \right)^{3/2}
\exp \left [ - \frac{m}{2T} (v_{x}^{2} + v_{y}^{2} + v_{z}^{2})
\right ], \; v_{x} < 0 
\end{eqnarray}

With this distribution, one can show that the fluid pressure is
\begin{equation}
P = \frac{2}{3} \langle n E \rangle = nT + \frac{1}{6}mv_{y0}^{2}.
\end{equation}
The fluid velocity is
\begin{eqnarray}
\langle \vec{v} \rangle & = & \langle \vec{v} \rangle_{+} 
+ \langle \vec{v} \rangle_{-} \nonumber \\
& = & \frac{1}{2} v_{y0} \hat{y} + 0,
\end{eqnarray} 
so that $| \langle \vec{v} \rangle | = v_{y0} / 2$.
To compute the fluid temperature used in the BGK distributions, we need
to subtract the drift energy,
\begin{eqnarray} 
T & = & \frac{P}{n} - \frac{1}{3} m \langle \vec{v} \rangle^{2}  \\
\label{tdrift}
  & = & T + \frac{1}{12} v_{y0}^{2}.
\label{tcouette}
\end{eqnarray}

The primary physical quantity appearing in the analytic 
models\cite{Will62,Cerc69,Cerc75} of Couette
flow is the x-y component of the stress tensor,
\begin{equation}
\Pi_{xy} \equiv \int d^{3} v \: m v_{x} v_{y} f(\vec{v})
- n m U_{x} U_{y},
\end{equation}
where $\vec{U}$ is the first (velocity) moment of $f$ (the flow 
velocity). Note that while we should have $U_{y} \neq 0$, $U_{x}$ should
be identically 0 since there is no net flow in the $x$ direction.
In the molecular-flow regime, the distribution function is not changed by
collisions as the particles move across the box. Furthermore, particles
coming from the $x = d$ side of the box make no net contribution to
$\Pi_{xy}$ (since the integrand is an odd function of $v_{y}$). 
Inserting Eq.~(\ref{fmax}), we find
\begin{equation}
\Pi_{xy}^{\rm fm} = \frac{n m}{4} \left( \frac{8 T}{\pi m} \right)^{1/2} v_{y0}.
\label{fmstress}
\end{equation}
As collisions become more important, 
the plate velocity information is
dissipated into random motions of the fluid to a greater and greater
extent, and $\Pi_{xy}$ falls below $\Pi_{xy}^{\rm fm}$.

The variational theory of Cercignani\cite{Cerc69,Cerc75} yields a 
relatively compact formula for $\Pi_{xy}$ that matches the 
numerical results of Willis\cite{Will62} to within $0.5 \%$.
With $\delta = 1 / K_{n}$,
\begin{equation}
\frac{\Pi_{xy}}{\Pi_{xy}^{\rm fm}} = \frac{a + \sqrt{\pi} \delta}
{a + b \delta + c \delta^{2}},
\end{equation}
where for the case of a BGK collision operator,
\begin{eqnarray}
a & = & \frac{4 - \pi}{\pi - 2} \nonumber \\
b & = & \frac{\pi^{3/2}}{2(\pi - 2)} \nonumber \\
c & = & 1.
\label{cercig}
\end{eqnarray}

The $x$-$y$ component of the stress tensor is computed along with other
test particle data during tracking. However, final processing of that
information into a usable form is only done when the \verb+COUETTE+
macro is enabled in {\tt flighttest.web}.

The Couette flow problem can be simulated in DEGAS 2 with only a few
code modifications to set up the required boundary conditions (at the moment,
the required particle distributions cannot be specified completely
by input or data files). Because the problem is relatively simple, it
represents a good starting point for learning how to use 
neutral-neutral collisions in DEGAS 2.

Most of these changes are invoked with  a \verb+COUETTE+ 
FWEB macro inserted into the files

To run this example, follow these steps:
\begin{enumerate}
  \item {\bf Set source code switches.} The source code file {\tt boxgen.web}
has within it a few different settings available which can be selected
by changing the value of the flag \verb+BOXRUN+. With this file open in
an editor, change the lines
\begin{verbatim}
@m DIM 2

@m BOXRUN 0
@m SOLN 1
\end{verbatim}
so that they now read
\begin{verbatim}
@m DIM 1

@m BOXRUN 41
@m SOLN 0
\end{verbatim}
Save the file and exit the editor. The \verb+DIM+ flag controls the
dimensionality of the geometry.  By setting it to 1, we get a purely
one-dimensional geometry, with periodic boundary conditions in the $y$
and $z$ directions.  The \verb+SOLN+ flag controls the analytic 
solution used with the ``Analytic\_fluid\_bench'' example 
(Sec.~\ref{analyticfluidbench}); since that's not needed here, we can
turn it off.
The files {\tt plate.web}, {\tt sources.web}, and {\tt flighttest.web} each
have the line
\begin{verbatim}
@m COUETTE 0
\end{verbatim}
near the top. Edit these files, changing the ``0'' to a ``1'' in each case
(be sure to revert to the default setting of ``0'' before running other 
problems).
  \item {\bf Switch to the example directory.} If the main DEGAS 2 
directory sits within your home directory:
\begin{verbatim}
cd ~/degas2/examples/Couette_flow
\end{verbatim}
  \item {\bf Copy the} {\tt degas2.in} {\bf file into your working
directory.}
\begin{verbatim}
cp degas2.in ../../workdir/degas2.in
cd ../../build
\end{verbatim}
(overwriting any {\tt degas2.in} file you may have had there already!).
  \item {\bf Prepare reference data.} This example uses some nonstandard
data files so, unlike the previous example, {\tt datasetup}
must be run. Note that the {\tt degas2.in} file points back to
the {\tt examples} directory for the input files needed here and in 
subsequent steps.
\begin{verbatim}
make datasetup
cd ../workdir
../bin/datasetup
\end{verbatim}
  \item {\bf Prepare problem data.}
\begin{verbatim}
cd ../build
make problemsetup
cd ../workdir
../bin/problemsetup
\end{verbatim}
If you look at the problem input file, you will see that there is only
one reaction, that for the neutral-neutral collisions. Note also that
the test species \verb+D2+ also appears in the list of background species
(the usual background species \verb+e+ and \verb-D+- appear here just because
{\tt boxgen} expects them to be there; they are not used in this simulation).
  \item {\bf Generate geometry and background files.} The {\tt boxgen}
program takes care of both of these (be sure that the \verb+BOXRUN+ macro
has been set to 41):
\begin{verbatim}
cd ../build
make boxgen
cd ../workdir
../bin/boxgen
\end{verbatim}
  \item {\bf Set the number of flights.} The background file {\tt bk\_bgktest.nc}
contains the specification of the number of flights to be run. Open this
file in an editor (Emacs will be the most convenient) and find the line
\begin{verbatim}
 source_num_flights = 100 ;
\end{verbatim}
The results in this directory were obtained with 10 BGK iterations of
4,000,000 flights each. Such a run could easily require
more than a day if run on a single processor (this run was executed
in 15 minutes using 16 Dual-core AMD Opteron 1 GHz processors).
You can probably run a smaller
number of flights (say, 100,000) and still get reasonable results.
Save the file and exit the editor.
  \item {\bf Set up the tallies.} 
\begin{verbatim}
cd ../build
make tallysetup
cd ../workdir
../bin/tallysetup
\end{verbatim}
  \item {\bf Run the code.}
\begin{verbatim}
cd ../build
make flighttest
cd ../workdir
../bin/flighttest
\end{verbatim}
The number of BGK iterations is controlled by the value of the 
\verb+bgk_max+ macro in flighttest.web. The value used here is 10. If you just
want to get an idea of how the calculation proceeds, you can choose a
smaller number. Note that the macro parameters controlling the convergence
tests, \verb+bgk_cvg_dens+ and \verb+bgk_cvg_pres+, 
are set to ridiculously small
values ($10^{-8}$) to ensure that all 10 iterations are completed.
  \item {\bf Check the text output.} At each iteration, the code writes out the 
test species density and other data in a file with a name of the form
{\tt densityxx.out} where {\tt xx} is replaced by the iteration number. 
The contents of these files are analogous to those of the other examples
(the density and pressure for each zone are in the third and fifth
columns, respectively; 
see the description in Sec.~\ref{analyticfluidbench}), except for the
last two columns. These are normally not used; here, they contain the
values and standard deviations of the $x$-$y$ stress tensor, $\Pi_{xy}$.
The comparison of those values with theory will be discussed below.
 
At the end of
each iteration, the background netCDF file ({\tt bk\_bgktest.nc}) is 
also updated (overwritten) with
the density, velocity, and temperature of the neutral background species.
A third file, {\tt cvg\_global.txt}, will provide information on the
global progress of the BGK iterations towards convergence. The four columns are:
\begin{enumerate}
  \item Iteration number,
  \item Test species number,
  \item Global fractional change in density (dimensionless); this is compared
with \verb+bgk_cvg_dens+ to decide whether or not to stop the iterations.
  \item Global fractional change in pressure. (dimensionless); this is compared
with \verb+bgk_cvg_pres+ to decide whether or not to stop the iterations.
\end{enumerate}
Note that on most machines this file 
does not appear until the unit is closed at the end of the run.

These files can be directly 
compared with the corresponding files in the {\tt examples} directory.
If you have repeated the run at the above-described length, you should
be able to match the text files exactly. The level of agreement in the
netCDF file should be several digits (perhaps as many as 10).
If a satisfactory match is not obtained, something is
amiss. Please contact the code authors (see Sec.~\ref{help}) if you feel
that a problem exists with the code as you downloaded it. 
  \item {\bf Check the binary output.} You can use the {\tt matchout}
utility (see Sec.~\ref{miscroutines}) to do this. Again, you should expect
the largest differences to be $\sim 10^{-8}$ or better.
  \item {\bf Compare with the analytic solution.} The rightmost set of data
in the {\tt densityxx.out} file are the stress tensor values, $\Pi_{xy}$. 
Ideally, these are constant across the problem space. Some deviation
exists here due to Monte Carlo noise. You
can compute the free molecular value from Eq.~(\ref{fmstress}). With both of
these numbers in hand, you can compare with Eq.~(\ref{cercig}). To do that, you
need to know $\delta$. The explicit formula used is
\begin{equation}
\delta = \frac{n \langle \sigma v \rangle d}{\sqrt{2 T / m}},
\end{equation}
where $n = 10^{20}$ m$^{-3}$ (enforced in {\tt flighttest.web} by code 
enabled with the \verb+COUETTE+ macro), $\langle \sigma v \rangle = 1.1
\times 10^{-16}$ m$^{3}/$s (from the {\tt d2d2\_bgktest.nc} file), $d = 0.1$ m
(set in {\tt boxgen.web}), and $m = m_{{\rm D}_{2}} = 6.689 \times 10^{-27}$ kg
(also from {\tt boxgen.web}). One would expect the temperature $T$ of the
fluid to be equal to that of the wall by virtue of the initial conditions. 
However, in cases where the energy associated with the plate velocity 
$v_{y0}$ is significant compared with $T_{\rm wall}$, the actual temperature
of the fluid is slightly higher than that of the wall as is indicated
by Eq.~(\ref{tcouette}). 
In fact, 
this example with $v_{y0} = 1000$ m/s 
shows better agreement with the theoretical result Eq.~(\ref{cercig})
if $T$ is set to the simulated fluid temperature (e.g., from the  
the {\tt bk\_bgktest.nc} file). 

With $T = T_{\rm wall} = 0.0258 {\rm eV} = 4.1336 \times 10^{-21}$ J, we find 
$\delta_{\rm wall} = 0.9894$.  Plugging this into Eq.~(\ref{cercig}),
\begin{equation}
\left( \frac{\Pi_{xy}}{\Pi_{xy}^{\rm fm}} \right)_{\rm c}(\delta_{\rm wall}) = 
0.6047.
\end{equation}
The free-molecular stress, Eq.~(\ref{fmstress}), is 
$\Pi_{xy}^{\rm fm}(T_{\rm wall}) = 0.20977$ Pa.
If we average the 10 values for $\Pi_{xy}$ in the {\tt density10.out} file, we
get $\Pi_{xy, {\rm sim}} = 0.13045$ Pa and 
\begin{equation}
\left( \frac{\Pi_{xy}}{\Pi_{xy}^{\rm fm}} \right)_{\rm sim} = 0.6219.
\end{equation}

This would qualify as pretty good agreement. However, if we average the fluid 
temperature values in the {\tt bk\_bgktest.nc} file, we get $T_{\rm sim} =
4.5866 \times 10^{-21}$ J. Then, $\delta_{\rm sim} = 0.9394$, and
\begin{equation}
\left( \frac{\Pi_{xy}}{\Pi_{xy}^{\rm fm}} \right)_{\rm c}(\delta_{\rm sim}) = 
0.6157,
\end{equation}
about a factor of two better agreement. Now, one might argue that we should 
also use
$T_{\rm sim}$ in computing $\Pi_{xy}^{\rm fm}$. However, the 
free-molecular stress, Eq.~(\ref{fmstress}) was computed directly
from the general distribution function for arbitrary $v_{y0}$ so that
$T = T_{\rm wall}$ in that formula, without any ambiguity. 
Qualitatively speaking, for small
$\delta$ the character of the source distribution dominates. As $\delta$
approaches or exceeds $1$, the temperature of the fluid in the volume,
$T_{\rm sim}$ 
computed from Eq.~(\ref{tdrift}), better describes the distribution.
The better agreement is shown over a range of inverse Knudsen 
numbers and two values of $v_{y0}$ in Fig.~\ref{couetteplot}.
\end{enumerate}

\begin{figure}
\begin{center}
\leavevmode
\includegraphics[width=5in]{Couette_flow}
\caption{Comparison of 8 DEGAS 2 simulations (symbols) 
of the Couette flow problem
with the analytic expression for the normalized shear stress 
Eq.~(\ref{cercig}) obtained by Cercignani\cite{Cerc69,Cerc75}. Two
different plate velocities $V = v_{y0}$ were used in the simulations. 
For the higher velocity, better agreement with the analytic expression
is obtained if the inverse Knudsen parameter is computed from the
actual fluid temperature from the simulation, $T_{\rm sim}$, rather
than the wall temperature, $T_{\rm wall}$.}
\label{couetteplot}
\end{center}
\end{figure}



\subsubsection{C-Mod 1-D: Conservation Checks}
\label{cmod1d}

This example demonstrates a run with complete hydrogen physics in a simple
geometry. In particular, this run includes both neutral-ion and 
neutral-neutral 
elastic scattering processes. The complexity of the physics also
makes for a good demonstration of DEGAS 2's conservation checks.

Experiments by Pitcher et al. \cite{Pitc00,Pitc01} on the Alcator C-Mod tokamak
indicated that the flow of neutral gas through material structures around
the plasma was being limited by its diffusion rate through the divertor plasma.
Pitcher \cite{Pitc01} developed a semi-analytic model that reproduced this
behavior. An attempt was made to precisely reproduce the
results of that model with DEGAS 2. The overall effort consisted of several
simulations in which the model assumptions were gradually relaxed towards those
consistent with normal DEGAS 2 operation. This example is from a simulation in
having atomic and surface physics assumptions more like a normal DEGAS 2 run.
The remaining difference from a ``full physics'' DEGAS 2 run (apart from the
greatly simplified geometry) is that the neutral source is a gas puff rather
than a recycling source. The full physics and geometry simulations are described 
in Ref.~\cite{Stot01}.

To run this example, follow these steps:
\begin{enumerate}
  \item {\bf Set source code switch.} The source code file {\tt boxgen.web}
has within it a few different settings available which can be selected
by changing the value of the flag {\tt BOXRUN}. With this file open in
an editor, change the lines
\begin{verbatim}
@m BOXRUN 0
@m SOLN 1
\end{verbatim}
so that they now read
\begin{verbatim}
@m BOXRUN 40
@m SOLN 0
\end{verbatim}
Save the file and exit the editor. The second parameter disables the
analytic solutions that are used in the original boxgen
example (Sec.~\ref{analyticfluidbench}).  Note that the 
\verb+DIM+ parameter should remain at its default value of
2.
  \item {\bf Switch to the example directory.} If the main DEGAS 2 
directory sits within your home directory:
\begin{verbatim}
cd ~/degas2/examples/CMod-1D
\end{verbatim}
  \item {\bf Copy the} {\tt degas2.in} {\bf file into your working
directory.} E.g., if you are working on a Sun, 
\begin{verbatim}
cp degas2.in ../../SUN/degas2.in
cd ../../SUN
\end{verbatim}
(overwriting any {\tt degas2.in} file you may have had there already!).
  \item {\bf Prepare reference data.} This example uses some nonstandard
data files so, unlike the previous example, {\tt datasetup}
must be run. Note that the {\tt degas2.in} file points back to
the {\tt examples} directory for the input files needed here and in 
subsequent steps.
\begin{verbatim}
gmake datasetup
./datasetup
\end{verbatim}
  \item {\bf Prepare problem data.}
\begin{verbatim}
gmake problemsetup
./problemsetup
\end{verbatim}
  \item {\bf Generate geometry and background files.} The {\tt boxgen}
program takes care of both of these. Be sure the \verb+BOXRUN+ macro was
set to 40:
\begin{verbatim}
gmake boxgen
./boxgen
\end{verbatim}
  \item {\bf Set the number of flights.} The background file {\tt bk\_oned.nc}
contains the specification of the number of flights to be run. Open this
file in an editor (Emacs will be the most convenient) and find the line
\begin{verbatim}
 source_num_flights = 40000 ;
\end{verbatim}
Save the file and exit the editor.
  \item {\bf Set up the tallies.} 
\begin{verbatim}
gmake tallysetup
./tallysetup
\end{verbatim}
  \item {\bf Run the code.}
\begin{verbatim}
gmake flighttest
./flighttest
\end{verbatim}
The code will execute 5 iterations automatically, each with 40,000 flights.
The number of BGK iterations is controlled
by the parameter \verb+bgk_max+ in flighttest.web; the default setting is 5. The
BGK iterations can alternatively be controlled by the convergence parameters
\verb+bgk_cvg_dens+ and \verb+bgk_cvg_pres+. 
The corresponding tests monitor global measures
of the changes in the density and pressure of the BGK species. Note that
these can be satisfied only if there are sufficiently many flights
to achieve a correspondingly high precision in the local density and
pressure values, no matter how many BGK iterations are run. In this example,
the parameters chosen are such that
\verb+bgk_max+ will be governing the number of iterations.

  \item {\bf Check the text output.} At each iteration, the code writes out the 
test species density and other data in a file with a name of the form
{\tt densityxx.out} where {\tt xx}
is replaced by the iteration number. 
The contents of these files are analogous to those of the other examples
(the density and pressure for each zone are in the third and fifth
columns, respectively; 
see the description in Sec.~\ref{analyticfluidbench}).

At the end of
each iteration, the background netCDF file ({\tt bk\_oned.nc}) is 
also updated with
the density, velocity, and temperature of the neutral background species.
This file can be used to restart the BGK iterations, e.g., with an increased 
number of flights. 
A third file, {\tt cvg\_global.txt}, providea information on the
global progress of the BGK iterations towards convergence. The four columns are:
\begin{enumerate}
  \item Iteration number,
  \item Test species number,
  \item Global fractional change in density (dimensionless); this is compared
with \verb+bgk_cvg_dens+ to decide whether or not to stop the iterations.
  \item Global fractional change in pressure. (dimensionless); this is compared
with \verb+bgk_cvg_pres+ to decide whether or not to stop the iterations.
\end{enumerate}
Note that on most 
machines this file does not appear until the unit is closed at the end of 
the run. The output netCDF file from this run is included as well.

These files can be directly
compared with the corresponding files in the {\tt examples} directory.
The numbers in the text files have only five
digits of precision and should be matched exactly. The level of agreement
in the netCDF file should be several digits (perhaps as many as 10).
If a satisfactory match is not obtained, something is
amiss. Please contact the code authors (see Sec.~\ref{help}) if you feel
that a problem exists with the code as you downloaded it. 
  \item {\bf Check the binary output.} You can use the {\tt matchout}
utility (see Sec.~\ref{miscroutines}) to do this. Again, you should expect
the largest differences to be $\sim 10^{-8}$ or better.
  \item {\bf Examine test particle balances} To simplify examination 
of the conservation 
checks, an input script for the {\tt outputbrowser} code is included,
{\tt bgkbalances}. The results are in {\tt bgkbalances.out}. The checks can be 
performed for the density, momentum, and energy of each of the test
species. Since \verb-D2+- does not move, the checks on it are satisfied
trivially and are not considered further.\footnote{In fact, the 
values of the ``total'' tallies
for D2+ are all determined by roundoff errors. For this reason,
you will not likely match these values exactly. You should convince
yourself that they are numerically much smaller than the other totals. If you
use {\tt matchout} to compare with the reference output netCDF file, these
D2+ totals will also stand out as apparently significant discrepancies.}
For each quantity, the
procedure consists of totaling the negative of the quantity lost to the
walls (the "current into wall" tally), the quantity coming in from the
walls (the "current out of wall" tally), and the "source rate" tally.  The
last item is further 
broken down by reaction (or source) to provide an indication of the
relative importance of each process.  For this example, we consider only
the $X$-component of the momentum (dimension \#1) since the problem does not vary
in the other two directions.  

For this case, the following totals are obtained:

\begin{description}
  \item[particles] \ \\ 
  \begin{tabbing}
       D$_{2}$ \qquad \= $3 \times 10^{15}$ \kill
       D \>  $6 \times 10^{16}$  \\
       D$_{2}$ \> $5 \times 10^{15}$ \\
  \end{tabbing}
  \item[momentum 1] \ \\
  \begin{tabbing}
       D$_{2}$ \qquad \= $< 3 \times 10^{15}$ \kill
       D  \> $< 1 \times 10^{-6}$ \\
       D$_{2}$ \> $3 \times 10^{-8}$ \\
  \end{tabbing}
  \item[energy] \ \\
  \begin{tabbing}
       D$_{2}$ \qquad \= $3 \times 10^{15}$ \kill
       D \> $2 \times 10^{-2}$ \\
       D$_{2}$ \> $-2 \times 10^{-3}$ \\
  \end{tabbing}
\end{description}

In addition, a global check involving the background particle number can 
be performed:

\begin{eqnarray*}
  &+& \mbox{Total source of deuterium atoms (the puff source),} \\
  & &    +4.0 \times 10^{21} \\
  & &  \\
  &-& \mbox{Total number of deuterium atoms lost to the walls,} \\
  & &   -(3.97174 \times 10^{22} + 2. \times 4.88694 \times 10^{21}) \\
  & & + (3.68843 \times 10^{22} + 2. \times 5.61804 \times 10^{21}) \\
  & &  \\
  &-& \mbox{Total number of deuterium atoms lost to the plasma (D+ source rate),} \\
  & & - 2.62903 \times 10^{21}  \\
  & &  \\
  & & = 7. \times 10^{16} \\
\end{eqnarray*}

There will usually be a small remainder in these balances; 
its magnitude in this case roughly coincides with the accuracy provided by
{\tt outputbrowser}. This remainder results from the enforcement of a minimum
weight via Russian roulette in subroutine {\tt follow}. By default, that 
minimum weight is $10^{-3}$ of the initial weight. 
% The Russian roulette 
% procedure is set up so as to not lose mass. (This is also true in
% general only on average
Consequently, particles, 
momentum, and energy are conserved 
only on the average (i.e., after many flights).
You can demonstrate this effect by reducing the minimum weight (\verb+WMIN+ in 
{\tt flight.web}) to, say, $10^{-7}$. 
Doing so will make the run substantially longer.
To see the difference in the conservation checks, you may
need to examine the output netCDF file directly (hint: run {\tt outputbrowser}
first to get ``rough'' values for these quantities and then use the 
search facility in Emacs to locate the corresponding data in the output
netCDF file).

  \item {\bf Transfers between BGK species.} In looking at the energy and 
momentum
transfers between BGK species, you should notice that no sources 
appear for the background partners. The reason for this is that the transfers
to and from these species do not have a clear physical meaning; all 
relevant information can be obtained from the corresponding test species.

However, the transfers between BGK test species, say between \verb+D+ and 
\verb+D2+
here, are significant. The process \verb+d2d_neut+ (see the ``momentum
source vector by reaction'' and ``energy source by reaction'' sections
of {\tt bgkbalances.out} under reaction 11) transfers
energy and momentum between the background \verb+D2+ and the test \verb+D+
(problem species, ``\verb+problem_sp+'', number 6); likewise, 
\verb+dd2_neut+ (reaction number 12) connects the
background \verb+D+ and test \verb+D2+ (\verb+problem_sp+ 7). 
One would expect that in a
converged state these two transfer rates would be equal (this actually
goes into the derivation of the algorithm). In this example, they are
(to within the error bars),
\begin{tabbing}
  mom. 1 to \verb+D+: \= xx $-1.31488 \times 10^{-2}$ xx \=  
mom. 1 to \verb+D2+: \= xx $+1.23601 \times 10^{-2}$ \kill
  mom. 1 to \verb+D+: \> $-1.31488 \times 10^{-2}$ \>  
mom. 1 to \verb+D2+: \> $+1.23601 \times 10^{-2}$ \\
  energy to \verb+D+: \> $-5.62387 \times 10^{2}$ \>
   energy to \verb+D2+: \> $+5.49343 \times 10^{2}$ \\
\end{tabbing}
(The molecules come off the wall at low energy and are being heated by the
warmer atoms that arise from elastic scattering with the plasma ions and
from dissociation.)

For this to occur, the reaction rate for these two processes must
be numerically identical. The original implementation of the BGK
algorithm in EIRENE had the same temperature dependent expression for
the two rates. However, because the temperatures of \verb+D2+ and 
\verb+D+ species are different
in general, the resulting rates had different numerical values. The
initial implementation of the BGK algorithm in DEGAS 2 
repeated this same mistake.  Due to a dearth of applications for
these reactions, this problem has not yet been remedied.

To obtain the nearly equal transfer rates seen above,
the {\tt degas2.in} file in this example actually points to a 
dedicated {\tt reactions.input} file in the same directory 
that uses modified (``hacked'') data files for these two elastic 
scattering processes.  To see the effect of unequal reaction 
rates, modify the {\tt degas2.in} file to instead point
to the {\tt reactions.input} in the {\tt data} directory.
Rerun both {\tt datasetup} and {\tt problemsetup}, then
run the code again.  Note that you should also rerun {\tt boxgen}
to start the code off in the same initial state.  The
result after five iterations can
be compared with 
{\tt bgkbalances\_default.out} that comes with the examples:
\begin{tabbing}
  mom. 1 to \verb+D+: \= xx $-9.39835 \times 10^{-3}$ xx \=  
mom. 1 to \verb+D2+: \= xx $+9.19450 \times 10^{-3}$ \kill
  mom. 1 to \verb+D+: \> $-6.15071 \times 10^{-3}$ \>  
mom. 1 to \verb+D2+: \> $+9.65537 \times 10^{-3}$ \\
  energy to D: \> $-2.68765 \times 10^{2}$ \>  
  energy to D2: \> $+4.44565 \times 10^{2}$ \\
\end{tabbing}

  \item {\bf Other results.} For reference purposes, 
this directory contains output files from a 
single iteration (on a Linux machine running
Portland Group Fortran with \verb+DEBUG=yes+) with 1000 flights. These
output files contain the string \verb+1K+. 
The primary purpose of these files
is to serve as a known, reproducible point of contact for testing
subsequent code modifications. It also demonstrates that these
conservation checks hold even if the code is run with relatively few
flights. To duplicate these files, you will need to stop the BGK iteration
process after the first (zeroth, actually) iteration. The easiest way
to do this is to change the number of iterations, near the
top of {\tt flighttest.web}, from:
\begin{verbatim}
@m bgk_max 5
\end{verbatim}
to
\begin{verbatim}
@m bgk_max 0
\end{verbatim}

You may notice that the tallies used for the conservation checks 
predominantly utilize collision estimators while the standard (zone-resolved)
tallies are based on track length estimators.  The reason is that the
former will explicitly demonstrate conservation of mass, momentum, and
energy since they are computed using the instantaneous test particle
attributes. The track length estimators will generally provide more
accurate values for the code output, but will exhibit conservation only
in a statistical sense.  You can demonstrate this by replacing the
estimator lists in each of the ``total'' tallies with the lists from the
corresponding zone-resolved tallies.
\end{enumerate}

\subsection{definegeometry2d and defineback Examples}
\label{defexamples}
\index{definegeometry2d}
\index{defineback}

The {\tt definegeometry2d} (for setting up geometries, obviously)
and {\tt defineback} (for specifying the corresponding sources
and plasma background) are sufficiently flexible and general
that specific examples are useful in illustrating the ways that
they can be used.

\subsubsection{Gas Puff Imaging Example}
\label{HeGPI}
\index{DG}
\index{Carre}

This example illustrates the use of the DG and Carre codes
with {\tt definegeometry2d}, and also demonstrates the 3-D
capabilities of {\tt definegeometry2d}.  
It is based on the DEGAS 2 simulation of the NSTX
Gas Puff Imaging (GPI) GPI shot 112811, documented in 
\cite{Stot2007}.

To run the 2-D version of this example, follow
these steps:
\begin{enumerate}   
  \item {\bf Switch to the example directory.}  If the main 
DEGAS 2 directory sits
within your home directory:
\begin{verbatim}
cd ~/degas2/examples/He_GPI
\end{verbatim}
  \item {\bf Copy the {\tt degas2.in} file into your working directory.}  
E.g., if you are
working on a 64-bit Linux machine,
\begin{verbatim}
cp degas2.in ../../LINUX64
cd ../../LINUX64
\end{verbatim}
  \item {\bf Prepare reference data.}
This example uses the standard reference data
files in the data directory. Hence, this step can be skipped if those files
have not been altered since you downloaded the code.
\begin{verbatim}
gmake datasetup
./datasetup
\end{verbatim}
  \item {\bf Prepare problem data.}
\begin{verbatim}
gmake problemsetup
./problemsetup
\end{verbatim}
  \item {\bf Generate geometry file.} Here, we outline 
briefly the preliminary steps 
leading up to the creation
of the objects referred to in the {\tt definegeometry2d} input
file. Users interested in learning
more about that process should consult the DG User's Manual (in 
the \verb+dg_carre/DG/DOC+ directory; 
especially the DEGAS 2 specific sections under
``Customizing DG'').

The starting points for the DG code are a specification of the
vacuum vessel shape [$(R,Z)$ coordinates; \verb+nstx_walls+ in this
case] and an equilibrium (values of
the poloidal flux on a rectangular $R,Z$ grid,
\verb+g112811.00250.equ+, already converted from the EFIT format to DG's 
preferred file format).  Once these are loaded into DG,
the user proceeds to graphically generate the input for
the Carre code.  Carre is run separately outside of DG to 
create a nearly orthogonal 2-D mesh with one
coordinate that follows the flux surfaces.  Although DEGAS 2 doesn't
really need such a mesh, having one facilitates specifying plasma parameters
(typically a function of a flux variable).  That mesh is then loaded into DG 
and the user proceeds to connect it to the vacuum vessel 
points forming the ``DG polygons'' 
(see the entry in the \href{definegeometry2d.pdf#page.17}{definegeometry2d
manual}).
This
process results in three files: the DG file (\verb+nstx_25.dg+),
the mesh (\verb+nstx.carre.021+) and
DG's objects (\verb+nstx_25.dgo+).  Only the latter two 
are read in by {\tt definegeometry2d}.

There are two {\tt definegeometry2d} input 
files here, one for the 2-D (axisymmetric: 
\verb+nstx_25_dg2d.in+) case and one for 3-D (\verb+nstx_25_dg3d.in+).  
The results
contained in the \verb+examples/He_GPI+ directory 
are for the much quicker running 2-D
case.  Note that 
many of the 3-D lines appear in the 2-D file, but are commented
out.  Both runs utilize just the portion of the Carre mesh directly
in front of the gas manifold (``puffer'' in the {\tt definegeometry2d}
input file); this is done using the additional arguments
to the \verb+sonnet_mesh+ keyword.  One would establish the values of these
indices during the creation of the DG polygons.  The 3-D run then only 
simulates a portion of the torus (arguments to the \verb+bounds+ and 
\verb+y_values+
keywords).  These steps keep the problem size and run time down.

The emulation of the GPI camera view is done with the code in 
\verb+gpicamera.web+.
Prior to compiling {\tt definegeometry2d}, the user should copy this file
to \verb+usr2ddetector.web+ (see the brief explanation at the top of the
default version of the file \verb+def2ddetector.web+).  This is actually
just a placeholder $5 \times 5$ 
pixel version of the camera.  The full resolution
image is created during post-processing with the {\tt postdetector} routine.
\begin{verbatim}
cd ../src
cp gpicamera.web usr2ddetector.web
cd ../LINUX64
gmake definegeometry2d 
./definegeometry2d ../examples/He_GPI/nstx_25_dg2d.in
\end{verbatim}
  \item {\bf Generate background file.} The detailed
specification of the plasma parameters is actually done with
the code in \verb+nstxgpi2.web+.  Like \verb+gpicamera.web+, 
this is a user-defined
subroutine (\verb+get_n_t+) that needs to be compiled into the driver code.
This routine will read the plasma density and temperature profiles 
obtained with the NSTX 
Thomson scattering system, \verb+ts_112811_base.TXT+ and map them on to 
the entire geometry assuming that the plasma density
and temperature are constant on flux surfaces.  On the Carre-generated
mesh, this is straightforward.  Outside of it, where there are only
triangles, an approximate procedure is used.

The \verb+plasma_file+ keyword in the {\tt defineback} input file has two
parameters used by the {\tt nstxgpi2.web} \verb+get_n_t+ routine.
The second is the name of the Thomson scattering data file mentioned above.
The first argument
provides the subroutine with the location 
of the separatrix in the Carre mesh
subset used by {\tt definegeometry2d}.
Specifically, the separatrix is assumed to be between the cells \verb+iz_sep+
and \verb!iz_sep+1!.
If one knows the number of ``cells'' used in DG's
``Create Surface(s)...'' dialog for the core (\verb+n_cells+) 
and the minimum \verb+iz+
used in the mesh subset (specified via the \verb+sonnet_mesh+ keyword in 
{\tt definegeometry2d}), \verb+iz_min+, 
the value of this is: \verb!iz_sep = n_cells - iz_min + 2.!

The second piece of input to {\tt defineback} is the specification of the
neutral sources.  In the GPI simulations, this is a gas puff from
the polygon representing the gas manifold, a vertical line
in the 2-D case.
The physical locations of the 2-D sources are
provided by the ``stratum'' and ``segment'' indices
in the file \verb+source_nstx_26+.  The ``stratum'' number
corresponds to the ``stratum'' keyword used in {\tt definegeometry2d}.
In general, the simplest way to determine the segment numbers
is to use the ``\verb+*+'' option (see the 
\href{defineback.pdf#page.9}{defineback documentation}) to
select all of the segments associated with a stratum and then
manually identify those desired for the problem at hand.
In this
problem, a constant flux source is used since the sizes of
the segments vary.  The overall magnitude of the source is arbitrary
in this problem, but physically realistic.  The source
file for the 3-D case (\verb+source_nstx_26_3d+) is similar, but also provides
the toroidal location of the source segments (\verb+Segment_iy+).  These
were chosen so that the source points fall along a nearly straight
line at the same angle with respect to horizontal as the actual gas
manifold. 
\begin{verbatim}
cd ../src
cp nstxgpi2.web usr2dplasma.web
cd ../LINUX64
gmake defineback
./defineback ../examples/He_GPI/nstx_30_db.in
\end{verbatim}
  \item {\bf Set the number of flights.}
The background file \verb+bk_nstx_30.nc+ contains the specification of the 
number of flights to be run.  Open this file in an editor (Emacs
will be the most convenient) and find the line: 
\begin{verbatim}
 source_num_flights = 100 ;
\end{verbatim}
The results in this directory were obtained with 100,000 flights.
However, the 3-D production run used  2,000,000.
  \item {\bf Set up the tallies.}
\begin{verbatim}
gmake tallysetup
./tallysetup
\end{verbatim}
  \item {\bf Run the code.}
\begin{verbatim}
gmake flighttest
./flighttest
\end{verbatim}
  \item {\bf Check the text output.}
The file {\tt density.out} can be directly compared with the one in the 
\verb+examples/He_GPI+ directory.  The numbers in here have only five 
digits of precision.  If an exact match is not obtained, something
is amiss. These results were obtained with {\tt Triangle} V.\ 1.6.
  \item {\bf Check the binary output.}
You can use the {\tt matchout} utility to do this.  The largest differences
should be $\sim 10^{-10}$ or smaller.
  \item {\bf Generate graphical output.}
For the 2-D simulation,
\begin{verbatim}
cp ../examples/He_GPI/geometry.inp .
gmake geomtesta
./geomtesta
\end{verbatim}
The currently recommended approach is to compile the code with the
\verb+GRAPH_FILE=SILO+ or \verb+SILO_HDF5+ 
option and use VisIt (built with Silo
support) to view the Silo format file generated by {\tt geomtesta}.
The primary output variables in this problem are the emission rate
of the helium 5877 \AA\  (\verb+he_5877+) viewed by the GPI camera
and the neutral helium density, \verb+spHeden+.

For 3-D simulations, two example input files for {\tt geomtesta} are
provided.  Keep in mind that these need to be copied to the
run directory and renamed {\tt geometry.inp} in order for {\tt geomtesta} to
find them.
\item {\bf Generate camera image and associated data.}
The camera simulated by the main code uses only a $5 \times 5$ mesh in
this problem so as to not lengthen the run.  The
generation of actual images suitable for comparison with
experiment is deferred to postprocessing with the {\tt postdetector}
code.  In principle, {\tt postdetector} can be parallelized (although recent
changes to the code have yet to be propagated into the MPI-specific
sections), further mitigating the time expense for producing
the simulated camera image.

While {\tt postdetector} has been generalized with regard to species
(e.g., He or D$_{2}$), there are still application specific settings
in the code.  These are primarily the camera resolution and
specification of the target plane; note that these same
quantities appear in the preamble of {\tt gpicamera.web}.
The orientation of the simulated image can also vary with
application; see subroutine \verb+wrdatap+ in {\tt postdetector.web}.

Open {\tt postdetector.web} in an editor and find the section of
code just above the start of subroutine |fill\_views\_master|, and
set the application to ``NSTX'':
\begin{verbatim}
@m APP NSTX
\end{verbatim}
As was the case in compiling {\tt definegeometry2d}, the file 
{\tt usr2ddetector.web}
should be a copy of {\tt gpicamera.web}.  Note that the camera resolution
parameters in {\tt usr2ddetector.web} 
are overridden by those in {\tt postdetector.web}
and are not needed here.
\begin{verbatim}
gmake postdetector
./postdetector
\end{verbatim}
The camera image and other data are written to HDF format files (this
code has not yet been adapted for use with the Silo format).  Here
are the contents of the most useful files:
\begin{description}
  \item[emtt] Total emission rate of the 5877 \AA\ line 
[in W/(m$^{2}$ ster)].  This is  the simulated camera image.
  \item[tane] Electron density (in $m^{-3}$) at the 
target plane (intersection of the camera
view and the idealized gas puff sheet).
  \item[tate] Electron temperature (in eV) at the target plane.
  \item[tarr] Major radius (in m) at the target plane.
  \item[tarz] Vertical coordinate (in m) at the target plane.
  \item[tns1] Neutral helium density (in m$^{-3}$) at the target plane.
\end{description}

Note, however, that the identification of the ``target plane zones'' is
an automated optimization process that assumes a relatively fine toroidal
discretization.  As such, its results are rather difficult to
interpret in the 2-D case.

The quantities in the other files require much more space to explain.  
The
interested reader can learn more in \cite{Stot2004} 
  \item {\bf Three-dimensional simulation.}
The process for producing the analogous three-dimensional simulation is
the same as above.  The only modifications are:
\begin{enumerate}
  \item The input file for {\tt definegeometry2d} is 
\verb+nstx_25_dg3d.in+.  Note
that {\tt definegeometry2d} will require much more time (on the order of an
hour) to run.
  \item The pointer to the source file in the {\tt defineback} input
file \verb+nstx_30_db.in+ must be changed to point to 
\verb+source_nstx_26_3d+.
  \item Many more flights are required for the main code, 
e.g., 2,000,000, and
the run time is correspondingly longer.
  \item One of the two 3-D input files should be used for
{\tt geomtesta}.
\end{enumerate}
The resulting camera image corresponds to Fig. 4(a) in 
\cite{Stot2007}.  However,
the latter also incorporated the effects of camera vignetting, which
are not accounted for by the {\tt postdetector} routine.
\end{enumerate}

\subsubsection{NCSX Example}

This example illustrates the use of {\tt definegeometry2d} with 
input files generated with a relatively simple
stand-alone code.  The same concepts could be used
to manually create {\tt definegeometry2d} input files.
This particular simulation provided
neutral penetraton estimates used 
in the design of plasma facing components for the
National Compact Stellarator Experiment \cite{Zarn2001}.  
Since the bulk of the recycled neutrals leave
the surface with a cosine distribution (i.e., peaked about the normal 
to the material surface), the
greatest penetration will be achieved in the same poloidal plane as 
the recycling surface.
Consequently, toroidally axisymmetric 
(i.e., a single plasma cross section at a 
particular toroidal angle) neutral transport calculations 
can place an upper bound
on the neutral penetration distance.

In scoping problems such
as this, satisfactory results can be obtained by assuming that
the plasma parameters are constant on flux surfaces (of course,
if one does know the plasma variation along a flux surface, this
assumption can be relaxed).  The DEGAS 2
geometry can then be constructed out of flux surface shapes, e.g.,
from an equilibrium calculation of some sort. The basic approach is
to define a ``wall'' corresponding to each flux surface and then
connect adjacent walls to form the ``polygons'' required
by {\tt definegeometry2d} (see 
\href{definegeometry2d.pdf#page.1}{its documentation} for
detailed descriptions of these objects).  More specfically, each pair of
adjacent flux surfaces  is used to construct two polygons (e.g., 
a lower and upper half) so that
the interior of both is topologically well defined.

This objective can be in principle be achieved by manual manipulation
of the coordinate data.  However, the process is sufficiently
straightforward, not to mention tedious, that a code can be constructed that
generates the {\tt wallfile} needed by
{\tt definegeometry2d}, as well as the 
polygons created out of those walls that form the bulk of 
the {\tt definegeometry2d} input file.

The starting point in this case was a 3-D equilibrium
file generated by the VMEC code, containing Fourier representations
of the NCSX flux surfaces.   An example code 
distributed with the VMEC file was first modified to evaluate these
surfaces at a set of input toroidal angles (actually just two, with
one represented here).  A uniform poloidal angle mesh (an
arbitrary choice) was then defined so that a given flux
surface could be represented by a set of discrete points.
There were many more flux surfaces in the VMEC equilibrium than were
needed for this application.  The set of surfaces was
reduced (again arbitrarily) so that the maximum distance
between two adjacent surfaces was 5 cm.  A ``wall'' was 
created out of each of these surfaces, with the
constituent points being the $R$, $Z$ coordinates of the
surface at each of the points along the poloidal angle mesh.
Because the relative location of these walls was known,
the generation of the corresponding polygons was a 
simple matter.  
The code also defined walls and polygons filling
the gap between the VMEC surfaces and the vacuum vessel
and connecting the vacuum vessel to the universal cell.
Details such as these depend on the particular
application and are, thus,
not worth describing here.

Another crucial task that must be begun at this stage is
establishing the mechanism that will allow {\tt defineback}
to map the plasma parameters onto these
polygons.  The first, and generally essential, step to doing
this is to label each polygon with a ``stratum'' number.
Note that DEGAS 2's internal geometry specification, 
contained in the geometry
netCDF file, has no record of the polygons
used by {\tt definegeometry2d}; it only knows about zones.
The stratum labels are used in conjunction with the
geometry's ``sectors'', but they are of little use within
the plasma volume.  For this reason, {\tt definegeometry2d} also 
generates its \verb+polygon_nc_file+ based on the
2-D geometry class (see the  
\href{classes.pdf#page.132}{geometry2d class documentation}).
Via the \verb+g2_polygon_stratum+ and \verb+g2_polygon_zone+
arrays in this class,
the user can connect these  
zones back to the polygons
specified on input to {\tt definegeometry2d}.
In this example, this objective is accomplished 
by the subroutine in the
{\tt ncsxplasma.web} file compiled into 
{\tt defineback}.
Be aware that the polygons in the \verb+polygon_nc_file+
file are derived from, but not identical to, those
input to {\tt definegeometry2d}.  Specifically,
polygons broken up with {\tt definegeometry2d}'s
\verb+breakup_polygon+ will remain as individual 
polygons in the \verb+polygon_nc_file+.  But, 
those processed with \verb+trianglulate_polygon+ or
\verb+triangulate_to_zones+ are represented by the
resulting triangles within the \verb+polygon_nc_file+ and, thus,
correspond to multiple polygons there.  The stratum
labels are propagated to the triangle polygons in the
process, however, permitting them and the zone numbers
to be related to the original polygons.

In this example, the plasma profiles are specified as
a function of the square root of the normalized toroidal flux.
Consequently, the geometry setup code also writes out a file,
{\tt radiiz1li383},
containing the toroidal flux corresponding to each stratum.
Because each pair of adjacent flux surfaces is used to construct
two polygons, these two polygons are assigned the same stratum
label in the {\tt definegeometry2d} input file.
The names of the files containing the plasma profiles, the
effective stratum radii, and the \verb+polygon_nc_file+
are input to {\tt defineback} and passed to the routine 
in {\tt ncsxplasma.web}.

One will also need to be able to tell
{\tt defineback} where to place the neutral sources.  This usually
amounts to manually identifying  the ``strata'' corresponding to
those sources within the {\tt definegeometry2d} input file, and
then following a simple procedure for deducing the specific ``segments''
of those strata that are to be used.

To run this example, follow these steps:
\begin{enumerate}   
  \item {\bf Switch to the example directory.}  If the main 
DEGAS 2 directory sits
within your home directory:
\begin{verbatim}
cd ~/degas2/examples/NCSX
\end{verbatim}
  \item {\bf Copy the {\tt degas2.in} file into your working directory.}  
E.g., if you are
working on a 64-bit Linux machine,
\begin{verbatim}
cp degas2.in ../../LINUX64
cd ../../LINUX64
\end{verbatim}
  \item {\bf Prepare reference data.}
This example uses the standard reference data files.
\begin{verbatim}
gmake datasetup
./datasetup
\end{verbatim}
  \item {\bf Prepare problem data.}
\begin{verbatim}
gmake problemsetup
./problemsetup
\end{verbatim}
  \item {\bf Generate geometry file.}
\begin{verbatim}
gmake definegeometry2d 
./definegeometry2d ../examples/NCSX/dg2dinz1li383
\end{verbatim}
  \item {\bf Generate background file.} As was described above,
the detailed specification of the plasma parameters is done
by the \verb+get_n_t+ routine in {\tt ncsxplasma.web}.  It requires
as input the three files described previously (the effective radii for
each stratum, the plasma profiles, and the \verb+polygon_nc_file+), 
and the plasma density and temperature scrape-off lengths (the plasma
parameters are only given on closed flux surfaces).

The remainder of the {\tt defineback} input file 
contains the specification of the
neutral sources.  In this example, this is a recycling
source at a limiter.  The corresponding ``stratum'' was
established as part of the process of creating the input
files for {\tt definegeometry2d}.  
A general, and relatively quick, technique for 
determining the source segment numbers
is to use the ``\verb+*+'' option (see the 
\href{defineback.pdf#page.9}{defineback documentation}) to
select all of the segments associated with a stratum and then
manually identify those desired for the problem at hand.
The source strength is arbitrary in this problem, although
a physically realistic value was selected.
\begin{verbatim}
cd ../src
cp ncsxplasma.web usr2dplasma.web
cd ../LINUX64
gmake defineback
./defineback ../examples/NCSX/z1li383_db.in
\end{verbatim}
  \item {\bf Set the number of flights.}
The background file \verb+bk_ncsx.nc+ contains the specification of the 
number of flights to be run.  Open this file in an editor (Emacs
will be the most convenient) and find the line: 
\begin{verbatim}
 source_num_flights = 100 ;
\end{verbatim}
The results in this directory were obtained with 10,000 flights.
  \item {\bf Set up the tallies.}
\begin{verbatim}
gmake tallysetup
./tallysetup
\end{verbatim}
  \item {\bf Run the code.}
\begin{verbatim}
gmake flighttest
./flighttest
\end{verbatim}
  \item {\bf Check the text output.}
The file {\tt density.out} can be directly compared with the one in the 
\verb+examples/NCSX+ directory.  The numbers in here have only five 
digits of precision.  If an exact match is not obtained, something
is amiss.  These results were obtained with {\tt Triangle} V.\ 1.6.
  \item {\bf Check the binary output.}
You can use the matchout utility to do this.  The largest differences
should be $\sim 10^{-10}$ or smaller.
  \item {\bf Generate graphical output.}
For the 2-D simulation,
\begin{verbatim}
cp ../examples/NCSX/geometry.inp .
gmake geomtesta
./geomtesta
\end{verbatim}
\end{enumerate}

\section{Run Control Parameters}

Several run control parameters are contained in the 
\href{classes.pdf#page.26}{sources class}. Ideally, these would be
set by the user in a dialog box just prior to the run. For now,
they must be changed manually by editing the background netCDF file.
The following subsections describe the features, roughly in order of usefulness
to the average user:

\subsection{Time Dependence}
\label{TimeDep2}
\index{time dependence}

The use of time dependence in DEGAS 2 is described above
in Sec.~\ref{TimeDep}.  For completeness, we only note here
the principal controlling parameters in the
\href{classes.pdf#page.26}{sources class}.  The switch
{\tt so\_time\_dependent} will be {\tt FALSE} (0) for steady-state
runs and {\tt TRUE} (1) for the time dependent case.
For the latter, the time interval for the run will 
be from {\tt so\_time\_initial} to {\tt so\_time\_final};
both have units of seconds.
There is also a source-group dependent parameter
{\tt so\_time\_varn} that controls the time
variation of tahat group.

\subsection{Checkpoints}
\label{checkpoints}
\index{checkpoints}

By default, the DEGAS 2 output netCDF file is written only at the end of
the run. Each source group can be broken up into an integral number
of pieces with an intermediate output file being written after
each such checkpoint. The array {\tt so\_chkpt(grp)} contains the number
of checkpoints for source group {\tt grp}. If {\tt so\_chkpt(grp)} is 0,
no file is written. If 1, an intermediate output file will be written
at the end of the source group, and so on. 

The intermediate output file does not contain post-processed results and
cannot be used with post-processing utilities. The raw data there are
suitable only for a restart of the code.

\subsection{Restarting}
\index{restarting}

The two main uses of the restarting capability are 

\begin{enumerate}
  \item To complete an interrupted (and checkpointed) run without having
to restart from the beginning,
  \item To increase the number of flights in a run to improve statistics.
\end{enumerate}

In both cases, the user needs to change the value of the flag 
{\tt so\_restart} from its default value of {\tt FALSE} (represented
by 0 in the background netCDF file) to {\tt TRUE} (1). The code will
expect to find an existing output file, with the name specified in
the {\tt degas2.in} file. This can be either an intermediate
output file from a checkpoint or a completed one generated at the
end of a run. 

The restarted run will extend the number of flights to the value specified
in the background netCDF file (by {\tt so\_nflights}). A few different
situations may arise:

\begin{enumerate}
  \item If the number of flights is the same as those in the output file,
no additional flights are run. This feature permits the data in an
intermediate output file to be post-processed into a ``completed''
output file.
  \item If the previous checkpointed run was interrupted, the restarted
run resumes from the point of the last checkpoint and completes the
specified number of flights. The results should match those that would
have been obtained if the run were to have run to completion on the first
try.
  \item If the user wishes to increase the number of flights for 
the {\em last} source group of a completed run, the code will compute the
required number of flights. The results will be the same as if the
flights were all done at once.
  \item If the user wishes to increase the number of flights in {\em more
than one} source group, the code will extend the run accordingly. However,
the results {\em will not} exactly match those obtained in a single run of
the same length. The reason is that by default a run uses a continuous
random number chain for all source groups. In a restart, that chain is
resumed at its end. Hence, the first source groups will be using 
random numbers for these new flights different from those they would have
had in a single run. Note that the resulting run should still be
{\em statistically equivalent} to the single run. The ``spaced seeds''
option has been added to provide a workaround to this minor shortcoming
of the restart procedure.
\end{enumerate}

\subsection{Seed String}

Prior to version 2.6 of DEGAS 2, the initial random number seed was
hardwired in {\tt flighttest.web}. The user can now change the initial
seed (and, hence, the entire random number chain) via 
{\tt so\_seed\_string} in the background netCDF file. 

The preferred mechanism for specifying the initial seed is as a
character string. This is transformed into the appropriate integer
representation using routines in {\tt random.web} 
(see Sec.~\ref{Randomnumbers}). The default value is ``12''. Reasons
to use different values include
\begin{enumerate}
  \item An unexpected, but perhaps not statistically significant, result
has been obtained. A repeat of the run with a different random number
seed will aid in establishing the validity of the result,
  \item The DEGAS 2 code package contains several tools for 
statistically analyzing results (e.g., {\tt matchout} and
{\tt matcheir}). By changing the random number seed, the effectiveness of
these statistical comparisons can be evaluated.
\end{enumerate}

\subsection{Spaced Seeds}

If the user anticipates needing to extend a multiple source group run to
a larger number of flights (e.g., to achieve a desired variance) and 
insists that the results match those of a single run, the flag
{\tt so\_spaced\_seeds} should be set to {\tt TRUE} (1) in the initial 
run. With this option enabled, the code will set the initial seed
for each source group to be 100000000 (set by {\tt so\_seed\_spacing})
flights apart. Each source group
in a restarted run will then be able to ``pick up where it left off'',
up to a total of 100,000,000 flights.

To minimize the chances of compromising the integrity of the random
number chain, {\tt so\_spaced\_seeds} should be left at its default
{\tt FALSE} (0) value, unless the user truly needs to be able to extend
the run and match the single run results.

\subsection{Direct Sampling}

The default (random) method of sampling the initial particle source in
DEGAS 2 leads to a source distribution that differs from the ideal,
input distribution by a fraction of order $1/\sqrt{N}$ ($N$ being the
number of flights). Yet, if we were to take those same $N$ flights
and divide them up by hand amongst the source segments, we would 
obtain an error of order $1/N$. Such an approach is impractical in
general as it requires specifying the number of flights once and for
all at the beginning of the run (a restart would not be possible) and
the possibility of overlap in the initial source positions would arise
if the number of flights were large enough.

An equally effective alternative can be developed based on hashing
algorithms described by Knuth\cite{Knuth97}. The idea is to replace the
pseudo-random number $\xi$ usually used in the source sampling process with
\begin{equation}
\xi = (x + i / \phi) \bmod 1,
\end{equation}
where $x$ is a single, fixed random number, $i$ is the flight number,
and 
\begin{equation}
\phi^{-1} = \frac{\sqrt{5} - 1}{2}
\end{equation}
is the golden ratio. This value yields the desired distribution properties
(as can be verified using {\tt sourcetest}), but minimizes the possibility
of overlap by ensuring that $\xi$ is far from low order rational numbers.

The sampling method used is governed by the flag {\tt so\_sampling}. The
default value of {\tt so\_random} yields the usual random sampling procedure.
This is the recommended value for most applications. The new, direct 
sampling method is selected by setting {\tt so\_sampling} to 
{\tt so\_direct}. This method was installed to see if it would lead
to more efficient iterations in coupling DEGAS 2 to plasma transport
codes.

\section{Adding Reactions and PMI}
\label{addreactions}

\subsection{Overview of Data Format}

The same, basic approach is used to store all reaction- and PMI-related data
in DEGAS 2. First, the data for each reaction or PMI is stored in its own 
netCDF file. That file contains:

\begin{enumerate}
  \item Name of reaction,
  \item Number of dependent variables,
  \item Organization of  data for each dependent variable (``table'' or
``fit''),
  \item Rank of each dependent variable (number of independent variables),
  \item Character name for each dependent and independent variable,
  \item Number of values of each independent variable (for tabular
data), 
  \item Data table, a single one-dimensional array. Indices into this
array are computed on the fly based on the number of values used for
each independent variable.
\end{enumerate}
For more details, see the description of the 
\href{classes.pdf#page.70}{cross section class}.  % Start of the description section
The corresponding information for PMI is in the description of the
\href{classes.pdf#page.78}{PMI format class}.

The only general tool available for creating new data files is {\tt ratecalc}
(see Sec.~\ref{othersetup} and \href{ratecalc.pdf#page.1}{the ratecalc
file} itself). Example input files for {\tt ratecalc} can be found in the
{\tt data} directory. Two other routines {\tt reactionwrite} and
{\tt pmiwrite} (again see Sec.~\ref{othersetup}) read existing text files
and write the data out into netCDF files. These two routines can be
used as examples for generating specific data files.

At run time, all of the reaction data files specified in the problem
input file (see Sec.~\ref{probleminfile}) (and only those) are read
in. The reaction rates (for PMI, the corresponding quantity is termed
the ``yield'') and the descriptions of their independent
variables and other characteristics are compiled into one set of arrays.
The information is organized in a similar manner to that of the
original data files, but with an additional index corresponding to
the problem reaction number. Additional processing is done to replace
character strings with integer indices, to make for more efficient searches.
Again, all of the actual reaction rate
data values are stored as a single 1-D array with macros used to provide
simple access to individual entries for a particular reaction.
The reaction rate is used in computing the current mean free path of the
flight and in choosing amongst the flight's possible reactions when a 
collision is indicated.

All other data present in the reaction data files is compiled into a set
of ``handling'' arrays, indicating that their primary purpose is to
control the specification of collision product velocities for that
reaction. In addition to an index for problem reaction number, there is
one for dependent variable number. The collision routines will 
search these data arrays for required information and will stop code
execution (in {\tt DEBUG} mode, anyway) if they are not found. The flow
of data through the various DEGAS 2 classes is shown in 
Fig.~\ref{classeschart}.

\begin{figure}
\begin{center}
\leavevmode
\includegraphics[width=5in]{classes_chart}
% \includegraphics[width=5in,angle=180]{classes_chart}
\caption{Flow of atomic and surface physics data through various DEGAS 2
classes during preprocessing.}
\label{classeschart}
\end{center}
\end{figure}

The data format, evaluation routines, and scoring mechanisms are designed so
that specific quantities can be scored just by adding the appropriate
data to the data file and creating a corresponding tally with that
quantity as the dependent variable. 

One aspect unique to the PMI is that specification of the outgoing velocity
distribution necessarily involves a ``fit'' since standard FORTRAN will
not permit a function (in this case, the function which interpolates
the data tables) to return a vector.

\subsection{An Example: Bateman Format Data}

A particularly involved example of a PMI is the specification of
``Bateman format'' data\cite{Heif82,Bate80}. A few uses of this for
reflection processes are currently in the code. A large set of data
for PMI in an improved version of this format will soon be added.

For each incident
energy and polar angle pair, $(E_{\rm in}, \theta_{\rm in})$,  we have a 
conditional distribution for the product atom
\begin{equation}
P_{E_{\rm in},\theta_{\rm in}}(v,\alpha,\phi) v^{2} dv \: \sin \alpha d\alpha \: d \phi,
\end{equation}
where $\alpha$ and $\phi$ are the outgoing polar and azimuthal angles, 
respectively.

This distribution is sampled from three separate 1-D distributions. First,
the outgoing velocity $v_{0}$ is specified by
\begin{equation}
f^{1}_{E_{\rm in},\theta_{\rm in}}(v) = \int \int P_{E_{\rm in},\theta_{\rm in}}(v,\alpha,\phi)
\sin \alpha d\alpha \: d \phi.
\end{equation}

The outgoing polar angle $\alpha_{0}$ obeys the distribution
\begin{equation}
f^{2}_{E_{\rm in},\theta_{\rm in}}(\alpha) = \int
P_{E_{\rm in},\theta_{\rm in}}(v_{0},\alpha,\phi) d \phi.
\end{equation}

And, finally, the outgoing azimuthal angle $\phi_{0}$ is taken from
\begin{equation}
f^{3}_{E_{\rm in},\theta_{\rm in}}(\phi) = 
P_{E_{\rm in},\theta_{\rm in}}(v_{0},\alpha_{0},\phi).
\end{equation} 

The inverse cumulative distribution, $F(\omega) = G^{-1}(\omega)$, with
\begin{equation}
G(x) = \omega = \int\limits^{x}_{0} f(y) \: dy,
\end{equation}
is specified, say, 5 values of $0 < \omega < 1$. For example,
for a normally incident $\theta_{\rm in}) = 0$, 
$(E_{\rm in} = 1$ eV H atom being reflected off
of Fe, the data are

$R_{N}(E_{\rm in},\theta_{\rm in}) = $ 
\begin{verbatim}
  8.27750E-01
\end{verbatim}

$F^{1}(\xi) =$ 
\footnotesize
\begin{verbatim} 
  1.99146E+00  2.25513E+00  2.32691E+00  2.36544E+00  2.40806E+00
\end{verbatim}
\normalsize
Interpolation into this array with a random number $\xi$ yields 
$\sqrt{E_{\rm out}}$, and thus $v_{0}$.

$F^{2}(\eta,\xi) =$
\footnotesize
\begin{verbatim} 
  5.16532E-01  6.99576E-01  8.11503E-01  8.99906E-01  9.67300E-01
  3.70481E-01  6.11799E-01  7.63020E-01  8.72539E-01  9.61285E-01
  3.80010E-01  6.18834E-01  7.63825E-01  8.75375E-01  9.61256E-01
  3.57688E-01  5.95520E-01  7.51752E-01  8.59766E-01  9.57489E-01
  3.48497E-01  5.17881E-01  6.53833E-01  7.79846E-01  9.23956E-01
\end{verbatim}
\normalsize
A 2-D interpolation into this table with the same first random number $\xi$
and a second one $\eta$ gives $\cos \theta_{0}$.

$F^{3}(\zeta,\eta,\xi) =$ 
\footnotesize
\begin{verbatim}
 -9.54674E-01 -5.93616E-01 -4.55439E-03  6.11607E-01  9.54292E-01
 -9.39719E-01 -5.34726E-01  7.70975E-02  5.91124E-01  9.38723E-01
 -9.47749E-01 -5.99118E-01 -5.34413E-03  5.58645E-01  9.59176E-01
 -9.52932E-01 -5.09963E-01  1.11845E-01  6.68895E-01  9.59592E-01
 -9.63812E-01 -6.32861E-01  2.84474E-02  6.49311E-01  9.65024E-01

 -9.49486E-01 -5.27460E-01  4.71202E-02  5.40823E-01  9.43092E-01
 -9.18789E-01 -5.34890E-01  3.87104E-03  5.97719E-01  9.43635E-01
 -9.35546E-01 -5.66353E-01 -6.17420E-02  5.65376E-01  9.28863E-01
 -9.39133E-01 -5.88361E-01  4.03247E-02  5.36441E-01  9.41553E-01
 -9.42917E-01 -5.66390E-01  4.88304E-03  5.72368E-01  9.58254E-01

 -9.59566E-01 -6.35221E-01 -7.50576E-02  5.72627E-01  9.59374E-01
 -9.43053E-01 -6.23133E-01 -8.35145E-02  5.13894E-01  9.33565E-01
 -9.42753E-01 -6.40605E-01 -1.50300E-02  5.51493E-01  9.32335E-01
 -9.28625E-01 -5.03451E-01  3.69856E-02  6.04412E-01  9.51354E-01
 -9.47567E-01 -5.75171E-01  6.78900E-02  6.17675E-01  9.61149E-01

 -9.42125E-01 -6.25972E-01 -6.74072E-02  5.10579E-01  9.26111E-01
 -9.24288E-01 -5.62514E-01 -6.67559E-02  5.49679E-01  9.57450E-01
 -9.45958E-01 -5.88903E-01  6.07203E-02  5.98492E-01  9.45981E-01
 -9.64714E-01 -6.17340E-01 -2.49212E-02  5.61052E-01  9.49261E-01
 -9.59291E-01 -6.47473E-01 -8.40439E-02  4.58677E-01  9.33637E-01

 -9.43142E-01 -5.25077E-01  4.79891E-02  6.22436E-01  9.53027E-01
 -9.43214E-01 -5.90395E-01  2.32011E-02  5.42675E-01  9.35411E-01
 -9.40994E-01 -4.99927E-01  7.81558E-02  6.19613E-01  9.44418E-01
 -9.35018E-01 -5.38589E-01  8.60652E-02  6.47157E-01  9.51505E-01
 -9.49873E-01 -5.38587E-01  8.93784E-02  6.39298E-01  9.66431E-01
\end{verbatim}
\normalsize
A third random number $\zeta$ is used with the other two in a 3-D
interpolation of this table to obtain $\cos \phi_{0}$. The re-use of
these random numbers in interpolating these three separate dependent
variables means that DEGAS 2 has to treat the random numbers as
fixed independent variables, just like electron density or
ion temperature, rather than generating the random numbers on the
fly in the interpolation process. Hence, the appearance of 
\verb+1st_random_number+, \verb+2nd_random_number+, and
\verb+3rd_random_number+ in the independent variables list
(see the \href{classes.pdf#page.78}{PMI format class}).

So, each file for such PMI contains
\begin{enumerate}
  \item $R_{N}(E_{\rm in},\theta_{\rm in})$,
  \item $F^{1}(\xi,E_{\rm in},\theta_{\rm in})$,
  \item $F^{2}(\eta,\xi,E_{\rm in},\theta_{\rm in})$,
  \item $F^{3}(\zeta,\eta,\xi,E_{\rm in},\theta_{\rm in})$.
\end{enumerate}
Again, all of the data values are stored in a single 1-D array
with the array behavior specified by macros
(see the \href{classes.pdf#page.78}{PMI format class}).

\subsection{Reaction Processing Routines}
\label{reactionprocessing}

Need to create ``set products'' and ``products'' routines, as well as
collision and track-length routines if new type.

\section{Defining Radiation Detectors}

The ``detector'' class contains the objects needed to assemble
synthetic diagnostics for filterscopes, imaging cameras, and other
diagnostics that effectively integrate light emission along a 
chord through the problem volume.  In the detector class, each
such chord is referred to as a ``view''.  

All of the views in the problem are characterized by:
\begin{enumerate}
  \item Two points,
  \begin{itemize}
    \item The first is taken as the starting point,
    \item The second fixes direction. 
  \end{itemize}
  \item Angular halfwidth(s) (i.e., each view consists of a cone or
pyramid with apex at the starting point); the rectangular algorithm
requires two independent halfwidths, one for each perpendicular
direction,
  \item An averaging algorithm, used to simulate the finite spatial
resolution obtained with real detectors.
The three techniques, each having 2-D and 3-D variants, available are
  \begin{itemize}
    \item Uniform weighting (i.e., square),
    \item Rectangular weighting (independent angular half-widths in two
perpendicular directions, generalizing uniform weighting to non-square
pyramids),
    \item Circular weighting.
  \end{itemize}
  \item A binning variable,
  \begin{itemize}
    \item Only examples are ``none'' and ``wavelength,''
    \item To use ``wavelength'', would also need to specify a 
minimum, maximum, and the number of wavelength bins.
  \end{itemize}
\end{enumerate}
The views thus defined can be separated into ``groups''. 
E.g., each group might represent a different diagnostic.
The actual variables are defined and described in the
\href{classes.pdf#page.49}{detector class}.

The contribution from the light emitted in
each zone to a view is usually precomputed so that the
signal can be quickly determined
during or after the run from the volumetric emission.
The array used to store these data potentially has size
(number of zones) $\times$ (number of chords).
However, the code only retains non-zero entries,
resulting in an array of modest size
since each chord
only runs through a handful of zones.  Should the user find
processing of the chord signals too computationally
demanding or that even the compressed array is
inconveniently large, a subset of the chords can be
used during the main run of the code and the full array
generated only in postprocessing.  The latter
can be done in parallel, if need be.

The first of the three following subsections, Sec.~\ref{frags}, describes
the computation of these signal contributions.  The second subsection,
Sec.~\ref{detectorsetup}, outlines the structure of the subroutine
used to set up the detector views and groups of views.  The third
subsection, Sec.~\ref{postdetector}, describes one approach to
computing the detector signal in post-processing.

\subsection{Signal Computation}
\label{frags}

There are two situations:
\begin{enumerate}
  \item Symmetric (2-D) geometry in which the detector chord
is perpendicular to the ignorable coordinate.  E.g., the
chord would be in the $x$ - $z$ (planar symmetry) or $R$ - $z$
(cylindrical symmetry) plane.
  \item Symmetric geometry in which the detector chord has a
component along the ignorable component or a 3-D geometry.
For simplicity, we will refer to this as the ``3-D'' situation.
\end{enumerate}

In the 2-D case, an effective average over a 3-D cone is computed
by exploiting the assumed symmetry in the ignorable coordinate.
The 2-D viewing cone is divided into 30 subchords / degree. The code
expects, but does not insist (a warning will be printed), 
that the apex of the cone does not
fall inside a plasma or vacuum zone. Each subchord is tracked from the
apex to the start of a plasma or vacuum zone. From there, the distance
the subchord traverses through each zone is computed. The subchord is
terminated when it reaches a solid zone; this is the actual end point.
The resulting contribution made to the view by zone $j$, call it $f(z_{j})$,
is for $i = 2N+1$
uniformly weighted subchords
\begin{equation}
f(z_{j}) = \sum_{i=-N}^{N}  d_{i,j} w_{i} / V_{j},
\label{zonefrag}
\end{equation}
where $d_{i,j} = $ is the distance along subchord $i$ through zone $j$, 
and $V_{j}$ is the volume of zone $j$.  The relative weight of each
subchord is $w_{i}$ with
\begin{equation}
w_{i}^{\rm uniform} = \frac{1}{4 \pi (2N + 1)},
\end{equation}
and
\begin{equation}
w_{i}^{\rm circular} = \frac{1}{2 \pi^{2} N \sqrt{1 - (i/N)^2}}.
\end{equation}
% Cylindrical and beam algorithms are drafted but not yet implemented;
% their 2-D reduction is not yet worked out.
For the rectangular case, $N$ is computed from $\delta_1$ (the halfwidth
associated with $\hat{\delta}_{\perp,1}$, i.e., the vertical/in-plane
direction) exactly as for the uniform case, and
\begin{equation}
w_{i}^{\rm rect} = \frac{1}{4 \pi (2N + 1)} = w_{i}^{\rm uniform}.
\end{equation}
$\delta_2$ does not appear: for a chord confined to the symmetry plane
(no $\hat{y}$ component), $\hat{\delta}_{\perp,1}$ as defined in the 3-D
case coincides exactly with the 2-D case's averaging direction
($\hat{\delta}_{\perp,1} = \hat{\delta}_{\parallel}\times\hat{y}$ in this
limit), and the marginal distribution of a uniform rectangular pyramid
along one perpendicular axis is uniform over that axis regardless of the
other axis's extent. The uniform case is simply the $\delta_1 = \delta_2$
special case of this.
The units of $f$ are ${\rm m}^{-2}/{\rm st}$.

In the 3-D case, each chord is modeled as a rectangular based pyramid
(for rectangular or uniform weighting) or a cone (for circular
weighting), each emanating from the apex of the view (the starting
point) and subtending a solid angle on the order of $\delta^{2}$, where
$\delta$ is the specified halfwidth expressed in radians. More
precisely, the solid angles are
\begin{equation}
\Omega^{\rm uniform} = 4 \arctan \frac{\sin^{2} \delta}{\sqrt{1
+ 2 \sin^{2} \delta}} \rightarrow 4 \delta^{2},
\end{equation}
\begin{equation}
\Omega^{\rm circular} = 2 \pi (1 - \cos \delta) \rightarrow \pi
\delta^{2},
\end{equation}
and, for a rectangular pyramid with dihedral half-angles $\delta_1$ and
$\delta_2$,
\begin{equation}
\Omega^{\rm rect} = 4\arcsin(\sin \delta_1 \sin \delta_2) \rightarrow 4
\delta_1 \delta_2,
\end{equation}
where the arrows indicate the value in the $\delta \ll 1$ limit.

The subchords are used to perform a 2-D integration over this solid
angle. The direction of each subchord is defined relative to that
of the detector chord; call the corresponding unit vector
$\hat{\delta}_{\parallel}$. The principal perpendicular dimension,
$\hat{\delta}_{\perp,1}$, is chosen to be orthogonal to
$\hat{\delta}_{\parallel}$ and oriented in the vertical, $\hat{z}$,
direction; however, if $\hat{\delta}_{\parallel} = \hat{z}$ then the
principal direction is aligned with $\hat{x}$ instead. The secondary
perpendicular dimension is
$\hat{\delta}_{\perp,2} = \hat{\delta}_{\parallel} \times
\hat{\delta}_{\perp,1}$.
For views with symmetry about the parallel direction (uniform and
circular) both perpendicular directions are arbitrary. However, for the
rectangular case the user should take care in specifying the angles
$\delta_1$ and $\delta_2$ according to their calculated directions.

The velocity vector sent to the tracking algorithm in the uniform
case is
\begin{equation}
\vec{v}_{j,k}^{\rm uniform} = \hat{\delta}_{\parallel}
+ \frac{(j - 1/2) - n}{n} \sin \delta \hat{\delta}_{\perp,1}
+ \frac{(k - 1/2) - n}{n} \sin \delta \hat{\delta}_{\perp,2},
\end{equation}
where the subchords are labeled by the indices $j$ and $k$,
both ranging from 1 to $2n$; currently $n = 6$.
The weighting for each subchord is
\begin{equation}
w_{j,k}^{\rm uniform} = \frac{\sin^{2} \delta}{4 \pi \Omega^{\rm uniform}
n^{2}}\left| \vec{v}_{j,k}^{\rm uniform} \right|^{-3/2}.
\end{equation}

The weight assigned to each subchord in the rectangular case is instead
expressed as the fraction of the total solid angle, $\Omega^{\rm rect}$,
normalized by the total $4\pi$,
\begin{equation}
w_{j,k} = \frac{1}{4\pi}\frac{\Omega_{j,k}}{\Omega^{\rm rect}},
\end{equation}
where the solid angle of each subchord patch is
\begin{equation}
\Omega_{j,k}^{\rm rect} = \frac{\tan(\delta_1)\tan(\delta_2)}{n^2}
\left\Vert \vec{v}_{j,k}^{\rm rect} \right\Vert^{-3},
\end{equation}
with the velocity vector sent to the tracking algorithm
\begin{equation}
\vec{v}_{j,k}^{\rm rect} = \hat{\delta}_{\parallel}
+ \frac{(j - 1/2) - n}{n} \tan(\delta_1) \hat{\delta}_{\perp,1}
+ \frac{(k - 1/2) - n}{n} \tan(\delta_2) \hat{\delta}_{\perp,2}.
\end{equation}
Thus, the weighting for each subchord is
\begin{equation}
w_{j,k}^{\rm rect} = \frac{\tan(\delta_1)\tan(\delta_2)}{4 \pi
\Omega^{\rm rect} n^{2}}\left\Vert \vec{v}_{j,k}^{\rm rect} \right\Vert^{-3}.
\end{equation}
In the $\delta_1 = \delta_2 = \delta$ limit this reproduces the
square-based pyramid case above to leading order in $\delta$; the two
are not identical at finite $\delta$, since {\tt de\_algorithm\_uniform}
scales the subchord offsets by $\sin \delta$ rather than $\tan \delta$,
for historical reasons, so {\tt de\_algorithm\_uniform} and {\tt
de\_algorithm\_rectangular} remain distinct, independently-selectable
algorithms.

In the circular case, the azimuthal angle's $2 \pi$ radians are
discretized into $\phi_{j} = (j - 1/2) \pi / n$, with 
$j = 1 \rightarrow 2n$.  The polar angle is given by
$\theta_{k} = (k - 1/2) \delta / n$, with $k = 1 \rightarrow n$.
Currently, $n = 6$. In terms of these angles, the velocity vector
sent to the tracking algorithm is
\begin{equation}
\vec{v}_{j,k}^{\rm circular} = \hat{\delta}_{\parallel}
+ \sin \theta_{k} ( \cos \phi_{j} \hat{\delta}_{\perp,1}
+ \sin \phi_{j} \hat{\delta}_{\perp,2}),
\end{equation}
and the weight for each subchord is
\begin{equation}
w_{j,k}^{\rm circular} = \frac{ \delta }{4 \Omega^{\rm circular} n^{2}}
\sin \theta_{k}.
\end{equation}

The resulting contribution to the chord by each zone
is again given by Eq.~(\ref{zonefrag}) in both cases.

These zone contributions can then be used during execution of the main
code or during post-processing to compute, say, the 
H$_{\alpha}$ signal $S_{\alpha}$
in the absence of recombination  (see Sec.~\ref{AtomicPhysics}) as
\begin{equation}
S_{\alpha} = \sum_{j} N_{1}(z_{j}) \left[ \frac{N_{3}}{N_{1}} 
A_{23} E_{23} \right] f(z_{j}),
\end{equation}
where $N_{1}$ is the number of neutral H atoms in zone $j$. The
bracketed quantity is interpolated from the input data file, with $E_{23}$
being the energy of the Balmer-$\alpha$ transition, and other quantities
defined in Sec.~\ref{AtomicPhysics}.
Note, however, that the code is completely general regarding wavelength 
and origin of radiation. Other line emission rates can be scored in
an analogous way if the appropriate data are present in the input data
files. 

\subsection{The Detector Setup Subroutine}
\label{detectorsetup}

The subroutine used to set up the detector views and groups of views
should be named {\tt detector\_setup}; all DEGAS 2 geometry definition routines
({\tt boxgen}, {\tt readgeometry}, and {\tt definegeometry2d}) call this subroutine.
The default, null implementation is contained in {\tt def2ddetector.web}.
A user developed {\tt detector\_setup} replacement for this routine should
be placed in a file called {\tt usr2ddetector.web} in the {\tt src} directory.
Example implementations of this routine provided with the code are
{\tt btopdetector.web} (a 1-D array) and {\tt gpicamera.web} (a 2-D imaging
camera; see also its \href{gpicamera.pdf#page.1}{documentation}).

The suggested structure of {\tt detector\_setup} is:
\begin{enumerate}
  \item For each view, 
  \begin{itemize}
    \item Set the two points defining the viewing chord,
    \item Set the algorithm and half-width values.
    \item Use these as arguments to the {\tt detector\_view\_setup}
subroutine (in {\tt dediagsetup.web}).
    \item This returns the contribution of each zone in the
problem to that view's signal, $f(z_{j})$ in Sec.~\ref{frags}.  This
quantity is copied into the {\tt de\_zone\_frags} array.
  \end{itemize}
  \item Assemble these views into one or more groups, defining each with
a call to {\tt de\_grp\_init} (in {\tt dediagsetup.web}).
  \begin{itemize}
    \item If the detector group is to compute the spectrum of light
emission, the requisite parameters are also passed to 
{\tt de\_grp\_init}.
  \end{itemize}
\end{enumerate}

\subsection{Detector Computation in Post-Processing}
\label{postdetector}

The {\tt postdetector} (see also its \href{postdetector.pdf#page.1}{documentation})
code provides an example of detector signal
computation during post-processing.  
{\tt postdetector} provides additional
data, such as the plasma parameters, along the
same detector chords not available in the main output.
The resulting detector based data are output directly to
netCDF or HDF format files, simplifying subsequent
processing.
{\tt postdetector}
can also be used to compute the complete signal if a low resolution
detector signal was used with {\tt flighttest} or to 

As is noted in Sec.~\ref{HeGPI}, the version of {\tt postdetector.web} 
distributed with DEGAS 2 was developed for simulating the camera
in Gas Puff Imaging experiments.  Device and experiment specific
information (namely a specification of the ``target plane''; see
\href{gpicamera.pdf#page.1}{gpicamera.web}) appear here.  Moreover,
it explicitly calls the {\tt gpi\_views} subroutine invoked by
{\tt detector\_setup} in {\tt gpicamera.web}.  Both there and in 
{\tt postdetector}, the coordinates of a particular pixel are
passed to {\tt gpi\_views}; it returns the end points and contributions
of each zone to the signal, $f(z_{j})$.
In this way, the detailed information required to deterimine
the view endpoints can be confined to a single 
subroutine, {\tt gpi\_views}.

\section{Diagnostic Sectors}
\label{Diagnostic}

As was noted in Sec.~\ref{Geometry},
sectors can be defined for purely diagnostic use. The zones adjacent
to the sector surface can be of any type. A given sector may be used to
specify more than one diagnostic (e.g., one may track particle current,
another might tally particle energy).

Diagnostic sectors are defined in a manner closely paralleling that of 
detectors. Diagnostic sector groups are used to identify sectors of
similar functionality. The default sectors defined by DEGAS 2's
geometry setup codes (such as
{\tt definegeometry2d}, {\tt readgeometry} and {\tt boxgen})
are used to construct the following
default diagnostic groups:
\begin{description}
  \item[Wall and Target Counts] Includes all wall and target sectors. No
independent variable is associated with this diagnostic group; it is only
capable of summing a particular test particle property (which will be
specified by a corresponding tally, particle mass, for example).
  \item[Exit Counts] Includes all exit sectors. No
independent variable is associated with this diagnostic group either.
  \item[Wall and Target Energy Spectrum] Specifies the particle
energy as the independent variable. As in the ``Counts'' groups, some
as-yet-unspecified particle property will be tracked. The difference is
that the tally will be divided up into the energy bins defined for this
group.
  \item[Wall and Target Angle Spectrum] Similar to the ``Energy Spectrum''
group, but the independent variable is the angle of incidence of the particle.
\end{description}

The particle properties to be associated with these diagnostic
groups for the purpose of defining tallies are directionally
dependent. For instance, the particle current leaving and entering
a zone may be specified separately. Details will be given in 
Sec.~\ref{tallies}.

Additional diagnostic groups can be defined as needed; this is 
especially straightforward with {\tt definegeometry2d}. 

\section{Output File}
\label{outputformat}

The output netCDF file (\verb+outputfile+ in {\tt degas2.in}) contains 
the tally data in four arrays:
\begin{enumerate}
  \item Sorted by the neutral source group, prior to post-processing,
  \item Summed over neutral sources, prior to post-processing,
  \item Sorted by neutral sources, after post-processing,
  \item Summed over neutral sources, after post-processing.
\end{enumerate}
All have standard deviation data, but their interpretation is not
guaranteed for the latter two since the impact of post-processing
(especially the addition of post-processed scores to a tally) is not
accounted for. If the output file is the result of a checkpoint
dump (Sec.~\ref{checkpoints}), keep in mind that post-processing
has not been done. Navigating these large arrays is extremely
tedious and not recommended except as a last resort. Instead,
the user should utilize {\tt outputbrowser}.
There is a separate output array for data to be passed to
2-D fluid plasma code.

A couple of text output files, largely of historical origin, are generated.
The {\tt density.out} file has been described already in 
Sec.~\ref{analyticfluidbench}. The {\tt sources.out} and
{\tt testdata.out} files are unformatted files used to transfer
data to the UEDGE code. All of the information in them is 
accessible through {\tt outputbrowser}, so no further 
description of them will be given.

\section{Other Documentation}
\label{otherdoc}

The file {\tt classes.web} contains an extensive description of what
the code looks like internally. Most of the common variables (i.e.,
the contents of the {\tt hweb} header files) are documented here
as well. Try following this 
\href{classes.pdf#page.1}{link} if
you want to have a look at it now.

\chapter{EIRENE Benchmark}

\section{Introduction}

Since UEDGE was first coupled to EIRENE and since EIRENE is currently the most
widely used neutral transport code in the field, the first step in coupling 
DEGAS 2 is to perform a benchmark against EIRENE. A simple slab, 
``single-null'' geometry and plasma generated by UEDGE is used as input
to both codes. An \href{http://w3.pppl.gov/degas2/DE_Bench/}{initial effort}
was performed without recombination. Essential information from that 
documentation has been included here for completeness.

This chapter describes a second round of benchmark exercises in which
recombination is accounted for.
The plasma $T_{e,i} \sim 1$ eV near plate in this case. For reasons 
to be explained below, several numerical differences
between the codes are accentuated under these conditions. We will describe
here the isolation and elimination of most of these differences, enabling
us to get agreement of densities and plasma sources to within $5 \%$.
In the process we have developed a means for quantitatively comparing the
code results. 

% Note: can probably just incorporate useful info from previous
% benchmark: change to energy distribution, basic atomic physics
% comparison (maybe), and performance benchmark.

The following sections will: 
\begin{enumerate}
  \item Describe the geometry, boundary conditions, and UEDGE plasma,
  \item Qualitatively compare ``out of the box'' code results,
  \item Describe how code results are quantitatively compared,
  \item Explain several differences between the codes, eliminating or
minimizing them when possible,
  \item Characterize the level of agreement between the codes,
  \item Go through a performance benchmark and optimization of DEGAS 2.
\end{enumerate}

\section{Problem Description}

The geometry is intended to be a 2-D slab representation of a toroidally
symmetric single-null scrape-off layer. The resulting ``box'' is 1 m
long, extending from a mirror boundary at $Z = 0$ to a molybdenum target
plate at $Z = 1$ m. Only half of the scrape-off layer is being simulated;
hence, the appearance of the mirror. The radial width is 0.05 m.
An exit is used to represent the core boundary at $R = -0.01$ m
between $Z = 0$ and $Z = 0.75$ m. The outer wall at $R = 0.04$ m
is assumed to be made of molybdenum. The section representing the
wall adjacent to the private flux region, at $R = -0.01$ m
between $Z = 0.75$ m and $Z = 1$ m, is also taken to be made of
molybdenum.  Finally, a mirror surface at $Z = 0.75$ m extending from
$R = -0.01$ m to $R = 0$ creates a private flux region above it.
The length of the box in the third dimension is 1 m, although
periodic boundary conditions are established at both ends.

The UEDGE mesh has a nonuniform spacing with 64 zones in the $Z$
direction and 32 in $R$. The geometry and plasma data are provided to
DEGAS 2 in a file generated during UEDGE post-processing by the BASIS 
utility. The electron temperature and density computed by UEDGE
with its fluid neutral transport model are shown in Figs.~\ref{te}
and \ref{ne}.

\begin{figure}
% \epsfverbosetrue
\begin{center}
\leavevmode
\includegraphics[width=5in]{te}
\caption{Electron temperature near the target ($Z = 1$ m) as computed by 
UEDGE with its fluid neutral model. }
\label{te}
\end{center}
\end{figure}

\begin{figure}
% \epsfverbosetrue
\begin{center}
\leavevmode
\includegraphics[width=5in]{ne}
\caption{Electron density near the target ($Z = 1$ m) as computed by
UEDGE with its fluid neutral model.}
\label{ne}
\end{center}
\end{figure}

\section{``Out of the Box'' Results}
\label{outofthebox}

Figures~\ref{eirDden} and \ref{degas2Dden} show how the atomic deuterium
densities computed by EIRENE and DEGAS 2 using the UEDGE plasma compare
with roughly standard physics.  The word ``roughly'' is meant to imply
that some of the changes to be described later in this chapter have
been incorporated into these simulations. For EIRENE, the random
number problem (see Sec.~\ref{sampling}) has been remedied. For DEGAS 2,
the treatment of reflection used on the molybdenum surfaces uses
the EIRENE data, including the attempts to mimic EIRENE's extrapolation
behavior (see Secs.~\ref{erefl} and \ref{arefl}).

\begin{figure}
% \epsfverbosetrue
\begin{center}
\leavevmode
\includegraphics[width=5in]{eir_D_den}
\caption{Neutral deuterium atom density computed by EIRENE.}
\label{eirDden}
\end{center}
\end{figure}

\begin{figure}
% \epsfverbosetrue
\begin{center}
\leavevmode
\includegraphics[width=5in]{degas2_D_den}
\caption{Neutral deuterium atom density computed by DEGAS 2 using ``standard
physics''.}
\label{degas2Dden}
\end{center}
\end{figure}

The peak EIRENE density is about $3 \times 10^{20}$ m$^{-3}$. For DEGAS 2,
the peak is only about $2 \times 10^{20}$ m$^{-3}$. In Sec.~\ref{codediffs},
we will list the causes of this discrepancy and describe the techniques
we use to eliminate or minimize them.

\section{Statistical Basis for Comparisons}
\label{statcompare}

Because Monte Carlo codes provide only a statistical description of
a result, e.g., the neutral density at a particular point in space 
and an estimate of its error, comparisons between codes require some
care. Differences between codes smaller than the standard error in their
results cannot be discerned, as this section will demonstrate. 
We will describe a quantitative procedure for comparing code results
which can be used to tell if the differences cannot be statistical
in nature.

The
standard error decreases with $1/\sqrt{N}$, where $N$ is the number
of flights, and can in principle be made as small as desired. For the
problem at hand, we've used 80,000 flights for each of the source
groups, resulting in standard errors of a few percent near the plate
for simple tallies such as the neutral atom density. The differences
between EIRENE and DEGAS 2 remaining after the steps described in 
Sec.~\ref{codediffs} have been taken are slightly larger than this, so
raising the number of flights further would not be useful.

Both DEGAS 2 and EIRENE compute the recombination  contributions to
the source terms (e.g., subtracting recombination rate from the ion
particle source) directly and add them to the final Monte Carlo results
at the end of the run. DEGAS 2 does provide the option for computing
these contributions in the standard Monte Carlo way (as a check), but
the statistical error is increased substantially.
Comparison of the code results is, thus, further complicated because the
computed variance estimates cannot be used directly with the final output
results.
Rather, our statistical comparisons must be made {\em without} these
``post-processing'' contributions.

DEGAS 2 has two sets of output arrays:
\begin{enumerate}
   \item One written before post-processing and containing relative
standard deviations,
   \item One with the post-processing results and no variance data.
\end{enumerate}
The screen output generated by EIRENE likewise contains 
results without these post-processing
contributions and their relative standard deviations, provided
they were requested in the input file. These data along with the
first of the two DEGAS 2 output array sets will be analyzed with
the procedure outlined here.  

The principal basis our statistical comparison is the Central Limit 
Theorem\cite{Span69}. Say we have the density score in a particular zone 
from a run of
$K$ flights; call it $\mu_{K}$. This result is effectively a sample 
drawn from a parent population; call its mean $m$ and the 
standard deviation $\sigma$. The objective of the Monte Carlo method is 
to estimate this mean $m$.

The Central Limit Theorem then says that the relative error
\begin{equation}
\epsilon_{\rm rel} \equiv \frac{ | \mu_{K} - m | }{ \sigma_{K} },
\end{equation}
where $\sigma_{K} = \sigma / \sqrt{K}$ is the standard error,
has a Maxwellian distribution.
That is, the probability that $\epsilon_{\rm rel} < 1$ is 
$68.3 \%$, $< 2$ is $95.4 \%$ and $< 3$ is $99.7 \%$.

To apply this directlyto our situation, we would need to:
\begin{itemize}
  \item Do many runs of both codes, 
  \item Find values of $m$ for each code which are consistent with
those results,
  \item Compare the $m$ values between codes,
  \item Repeat the process for each zone.
\end{itemize}
 
Such a process is too lengthy to permit extensive comparisons
of different variables and input conditions. We propose using
instead a simpler, compromise procedure.
Say we have a density score from DEGAS 2 $\mu^{\rm D}_{K}$ and
one from EIRENE $\mu^{\rm E}_{K^{\prime}}$,
We begin by {\em postulating} that they have the same $m$,
If this is true, the random variable $z \equiv \mu^{\rm D}_{K}
- \mu^{\rm E}_{K^{\prime}}$ has mean $m_{z} = 0$,
Then, if each of the codes' standard deviation estimates
$s^{\rm D}_{K}$ and $s^{\rm E}_{K^{\prime}}$ are also consistent with
$\sigma$, one can show that the distribution of $z$ about its
mean $0$ will have a standard deviation
\begin{equation}
s_{z,K,{\rm eff}} = \sqrt{(s^{\rm D}_{K})^{2} 
+ (s^{\rm E}_{K^{\prime}})^{2}},
\end{equation}
This is just what one expects from propagation of errors.

We are not done yet in that applying this would still require
many runs of both codes.
However, we do have many zones in each run.
{\em If} the scores in each zone are uncorrelated with other
zones, we could treat the value of 
\begin{equation}
\epsilon_{z,{\rm rel}} = 
z / s_{z,K,{\rm eff}}
\label{epsilonrel}
\end{equation}
in each
zone as a separate sample and compile a distribution of 
$\epsilon_{z,{\rm rel}}$ by comparing the results in each zone.
If we find this distribution to be Maxwellian, we may infer that our 
postulates are correct, i.e., the codes are giving the same results.
If the distribution strongly deviates from a Maxwellian, we can suspect
that the codes do not agree.

The extent to which scores in two zones are correlated depends
on many factors. Since we do not require precise statistical agreement
between the codes (a few percent of systematic error is acceptable), we
will assume that the effect of these correlations is too small to prevent this
test from detecting significant (more than a few percent) differences.

One other consideration is that the Central Limit Theorem only
applies for ``large enough $K$''.
For a given zone, this effectively means a ``small enough
relative standard deviation''.
For these comparisons, both codes produce relative standard
deviations near the target of a few percent.
If we restrict the comparison to zones with relative standard
deviations below some value, say $\sigma_{\rm max} = 20 \%$, 
we should still have a large
number of zones to work with.

Note that the codes output the relative standard deviation,
$\sigma / (m \sqrt{K})$.
To get $\sigma^{\rm D}_{K}$ we multiply the 
output relative standard deviation by $\mu^{\rm D}_{K}$.

\section{Code Differences}
\label{codediffs}

\subsection{Atomic Physics}

The common ancestry of EIRENE, DEGAS, and DEGAS 2 is apparent in their
atomic physics data. All three
codes rely upon the molecular data in the Janev book\cite{Jane87} (see 
Sec.~\ref{h2rate}, however). The reaction rates
are explicitly the same in all three codes, as are the dissociation energies,
and the electron energy loss rates. 
Note also that since the 
scoring of momentum transfer to the plasma in EIRENE is relatively new, 
its implementation is incomplete: there is no momentum transfer for the
molecular dissociation processes. However, given the dominance of the other
momentum sources, this should not be a problem. DEGAS 2 tracks all three
components of the momentum transfer in all reactions.

This version of EIRENE (with the present input file) uses the reaction
rate for charge exchange in the Janev book\cite{Jane87}. 
No attempt is made to ensure
that the sampled background ion velocities are consistent with the 
charge exchange cross section\cite{Heif82}.
Furthermore, the momentum and energy transfer expressions are correct only
in the limit of $\langle \sigma v \rangle = \sigma v$ 
(see Sec.~\ref{estdiffs}). More recent
versions of the code certainly do better than this\cite{Reit93}.

The DEGAS 2 standard charge exchange data comes from
the
Janev-Smith database\cite{Jane93}, with 
consistently computed reaction rates, and
energy and momentum transfer rates\cite{Reit93}. The simulation described in
Sec.~\ref{outofthebox} used these data. For subsequent runs, data
equivalent to those in EIRENE will be substituted.

DEGAS 2 and EIRENE each rely on their own collisional 
radiative\cite{Reit93} codes for 
the multi-step electron-impact ionization and recombination of 
hydrogen data. 
These two codes have been compared several times in the past.
The basic
difference is that the EIRENE code, AMJUEL, 
is based upon the Johnson-Hinnov ionization cross sections\cite{John73}.
The data used in DEGAS 2 are described in Sec.~\ref{AtomicPhysics}.

Figures~\ref{hionize} and \ref{elossfig}
compare the rate data. The AMJUEL ionization rates are noticeably lower
for $T_{e} < 10$ eV. This is likely responsible for most of the
differences between the densities in Figs.~\ref{eirDden} and 
\ref{degas2Dden}. The energy loss rates agree. However, when normalized
to the ionization rate, to give the energy lost per ionization, the
differences of Fig.~\ref{hionize} are factored in. 

\begin{figure}
% \epsfverbosetrue
\begin{center}
\leavevmode
\includegraphics[width=5in]{Hionize_Comparison}
\caption{Comparison of effective collisional radiative electron impact
hydrogen ionization rate from EIRENE's
AMJUEL data file and from DEGAS 2's IRLS code (contained in the
publicly distributed file {\tt ehr2.dat}) at
densities of $10^{18}$ m$^{-3}$ and $10^{20}$ m$^{-3}$. The lower
figure shows the ratios of the rates computed by the two codes.}
\label{hionize}
\end{center}
\end{figure}

\begin{figure}
% \epsfverbosetrue
\begin{center}
\leavevmode
\includegraphics[width=5in]{E_Loss_Comparison}
\caption{Comparison of effective collisional radiative electron energy
loss rate from EIRENE's
AMJUEL data file and from DEGAS 2's IRLS code (contained in the
publicly distributed file {\tt ehr2.dat}) at
densities of $10^{18}$ m$^{-3}$ and $10^{20}$ m$^{-3}$. The upper figure
shows the energy lost per ionization; the lower gives the power
lost.}
\label{elossfig}
\end{center}
\end{figure}

Figure~\ref{hrecombine} shows that the recombination rates agree well. 
The reason is that the Johnson-Hinnov\cite{John73} 
recombination cross section data are used in both collisional radiative
codes. The energy exchange rates agree for $T_{e} < 10$ eV, but disagree
at higher temperatures. This quantity includes the average energy of
a radiatively recombining electron; the expression used to compute 
it involves\cite{Reit93} 
$d \langle \sigma v \rangle_{\rm recombination} / d T_{e}$. Apparently,
the AMJUEL values were obtained by differentiating a fit to
$\langle \sigma v \rangle_{\rm recombination}$. The values used in DEGAS 2,
however, were obtained by differentiating closed form expressions for
$\langle \sigma v \rangle_{\rm recombination}$ and should be correct.

\begin{figure}
% \epsfverbosetrue
\begin{center}
\leavevmode
\includegraphics[width=5in]{Recomb_Plots}
\caption{Comparison of effective collisional radiative electron recombination
rates from EIRENE's
AMJUEL data file and from DEGAS 2's IRLS code (contained in the
publicly distributed file {\tt ehr2.dat}) at
densities of $10^{16}$ m$^{-3}$ and $10^{22}$ m$^{-3}$. The upper figure
shows the effective recombination rates; the lower gives the power
lost.}
\label{hrecombine}
\end{center}
\end{figure}

As with charge exchange,
in Sec.~\ref{outofthebox} DEGAS 2 used the {\tt ehr2.dat} data. 
For subsequent runs, the AMJUEL data will be used in both codes.

\subsection{Source Sampling}
\label{sampling}

The ion current distribution to the target is sampled according to
the relative fractions of current striking each segment (i.e., each
radial zone). Likewise, the probability distribution used in sampling
the ``birth'' zone of a recombining ion is given by the fraction of the
total recombination occurring in that zone. Whether or not a code is
correctly sampling these distributions can be easily determined. For
DEGAS 2, the {\tt sourcetest} utility was written to do just this
(see Sec.~\ref{testroutines}). Cruder means were used to extract the
requisite data from EIRENE.

Figure~\ref{sourcesegments} shows that the two codes are able to match
the input distribution (to within the error bars). To make the comparison
quantitative, we use a $\chi^{2}$ test\cite{Bevi69}. The result is $p$, 
the probability that sampled distribution  matches the parent distribution.
For the 5000 samples used in generating Fig.~\ref{sourcesegments}, 
the DEGAS 2 results yield $p = 0.97$. For EIRENE, we get $p = 0.69$. 
Since $p$ in both cases is not too much smaller than 1, we conclude that
the input distribution is being adequately sampled.

\begin{figure}
% \epsfverbosetrue
\begin{center}
\leavevmode
\includegraphics[width=4in]{source_segments}
\caption{Sampled target flux distributions, as a function of radial zone,
for DEGAS 2 and EIRENE. The curve labeled ``Input'' is the ideal
result.}
\label{sourcesegments}
\end{center}
\end{figure}

We use 50,000 samples of the recombination source. Because of the
large number of zones, a simple plot analogous to Fig.~\ref{sourcesegments}
is unintelligible. Furthermore, not all of the problem zones are 
adequately sampled to permit the application of a $\chi^{2}$ test to the
whole problem space. Instead, we isolate the 648 zones 
with highest recombination rate (the lowest
temperature region, chosen somewhat arbitrarily). Applying a
$\chi^{2}$ test to samples taken in those zones yields for
DEGAS 2 $p = 0.94$. But, for EIRENE $p = 0.21$. Sensing that this is
small enough to cause concern, we repaired a known problem with the
EIRENE random number generator (as received, the code had random number
seeds which differed by unity; the fix is to make one of these integers
very different). The result is much improved: $p = 0.79$. All other
EIRENE runs described in this chapter utilize this bug fix.

\subsection{Energy Distribution of Reflected Atoms}
\label{erefl}

The EIRENE surface physics is more detailed than 
that of the old DEGAS and, hence, DEGAS 2. Consequently, we will
not make an overall comparison with different
surface physics data. Instead, we have translated the Bateman format data 
for reflection of D off of Mo taken from EIRENE's {\tt TRIM} input
file into a DEGAS 2 format. To be more precise, the Mo data obtained from
Fe data using a scaling argument\cite{Ecks91}.

These data prescribe the reflection coefficient, the outgoing energy, and
two outgoing angles as a function of the incident energy and polar angle.
The Bateman format is described in detail elsewhere\cite{Heif82}.
Briefly, these data specify these parameter values (e.g.,
the outgoing energy) at
intervals of the cumulative distribution function: 0.1, 0.3, 0.5, 0.7, 0.9.
One can also establish data points at each end of the distribution. 
For example,
a minimum energy might be the wall temperature; a maximum would be the
incident energy. As originally implemented in DEGAS 2, this 
additional information to
extrapolate beyond the ends of the data. Doing this carefully required
specifying additional cumulative distribution values at 0, 0.2, 0.4, 0.6,
0.8, and 1.0 (intermediate values being obtained by linear interpolation).

EIRENE, however, mindlessly extrapolates the data and enforces the minimum
and maximum values after sampling. The result is that the sampled maximum
and minimum may be noticeably different from what one expects.
The \href{http://w3.pppl.gov/~dstotler/DE_Bench/full_code.html}{first set of benchmark runs}
indicated that the difference in the
handling of the lowest reflected energy atoms was significant. The
reason is that the atom density near the target plate (i.e., where
ionization is lowest) scales like the inverse of the atom velocity,
and hence is impacted strongly by the lowest energy atoms. To 
eliminate this difference, we reset the lower bound on the energy 
so DEGAS 2 would mimic the EIRENE extrapolation. Figure~\ref{edistrevised}
shows that the original DEGAS 2  reflected energies for a normally 
incident 19 eV atom match the input data well
(the difference at the lower end is due to the use of logarithmic 
interpolation in DEGAS 2). The revised DEGAS 2 reflected energies differ
from the input, but better match the values obtained from EIRENE.

\begin{figure}
\begin{center}
\leavevmode
\includegraphics[width=5in]{E_dist_revised}
\caption{Comparison of the cumulative distribution of reflected energies
sampled in DEGAS 2 and EIRENE
with the input data obtained from EIRENE's {\tt TRIM} file.
The ``revised'' DEGAS 2 curve indicates that the code is now able to
mimic the simple extrapolation used in EIRENE. These data are for a
19 eV deuterium atom incident on molybdenum.}
\label{edistrevised}
\end{center}
\end{figure}

A second difference becomes significant with the present set of
benchmarks including recombination. The temperatures in this plasma
are sufficiently low that a significant number of atoms are striking the
target with energies below 1 eV. EIRENE assumes that such atoms are
always absorbed (and desorbed thermally as molecules in this simulation).
DEGAS 2 treats the data literally; the reflection coefficient is 
about 0.2 at 1 eV. For the initial comparison run (see Sec.~\ref{outofthebox}),
DEGAS 2 was run in this manner, although the energy and angular
extrapolation changes described above and in the next subsection
were implemented. For all subsequently discussed benchmark
simulations, DEGAS 2 will assume absorption for
incident atoms of less than 1 eV energy.

\subsection{Angular Distribution of Reflected Atoms}
\label{arefl}

The cosines of the angular distribution are sampled in a way analogous
to the energy.
In this case, the bounds of the sampled parameters are clearer still since
they are the cosines of the angles (i.e., 0 and 1). However, during this
benchmark exercise, a discrepancy in the near target deuterium density
was again connected with the distribution of neutral velocities in the
$Z$ direction. With the difference due to energy extrapolation (see
Sec.~\ref{erefl}) having been eliminated, the cause this time was
traced to the extrapolation of the cosine of the outgoing polar 
angle, $\cos (\theta_{\rm out})$. 

\begin{figure}
\begin{center}
\leavevmode
\includegraphics[width=5in]{costh}
\caption{Comparison of the cumulative distribution of reflected 
outgoing polar angles sampled in DEGAS 2 and EIRENE 
with the input data obtained from EIRENE's {\tt TRIM} file.
These distributions were compiled from all of the initial incident
{\em ions} (i.e., at the start of the flight) in a particular
zones. All of these ions strike the target normally with 6.88 eV.}
\label{costh}
\end{center}
\end{figure}

Figure~\ref{costh} shows the sampled cumulative distribution from two
DEGAS 2 runs and from an EIRENE run. These data were compiled from
all of the ions incident on the target in a particular zone (normal
incidence at 6.88 eV). The curve labeled ``original'' shows the DEGAS 2
treatment extending down to near $\cos(\theta_{\rm out}) \simeq 0$.
The ``revised'' curve demonstrates that our modifications to the DEGAS 2
data successfully mimic the EIRENE extrapolation behavior.


 \subsection{H$_{2}$ Dissociation Rate}
\label{h2rate}

Significant differences in the electron energy sink were traced to
a discrepancy in the H$_{2}$ dissociation rate (reaction 2.2.5
in the Janev book\cite{Jane87}) at the lowest electron temperatures,
$T_{e} \simeq 1$ eV. Unfortunately, both the plot and the fit 
for these data in the Janev book are {\em wrong}. One needs
to refer to the corresponding preprint even to get the correct
values for the fit coefficients. EIRENE has those values, but applies
them down to the smallest $T_{e}$ used in the code, 1 eV. However,
the fit is valid down to only 1.2 eV. The dissociation rate values
used in DEGAS 2 at these lowest temperatures were taken directly from
the tabular data provided with the original DEGAS code\cite{Heif82},
not from the fit. The point is that these values differ from those
computed by the fit (presumably due to the fit not being valid there).
No record explaining the origin of the DEGAS values, but we assume
that they were taken from the initial data compilation from which the
fit was derived. Regardless, the remedy for the objectives of this
benchmark is to replace these low $T_{e}$ dissociation rates with the
same values used in EIRENE. The results of Sec.~\ref{outofthebox}
have the DEGAS values; runs discussed in subsequent sections will
have the EIRENE values.

\subsection{Estimator Differences}
\label{estdiffs}

With DEGAS 2 run as in Sec.~\ref{outofthebox}, the relative standard
deviations for the electron and ion energy sources were substantially
larger than those for EIRENE (these differences were most noticeable
close to the target). The cause of the higher electron energy
source variance was that DEGAS 2 was using a collision estimator
for the molecular contributions to that quantity (for historical reasons), 
while EIRENE was using a track length estimator. Switching the estimator
used in DEGAS 2 required a minor change to {\tt tallysetup} (see
Sec.~\ref{maincode}). At the same time, we changed the scoring of
the electron impact ionization of deuterium contributions to the ion
energy and momentum sources from being done ``post process''
(i.e., computed from the neutral density at the end of the run) to
a track length estimator. This allowed the variance contributions from
this process to be accounted for in these scores and was essential for
comparing with the EIRENE values (see Sec.~\ref{statcompare}).

The main reason for the larger ion energy source variance in DEGAS 2,
however, was that charge exchange was being treated by the collision
estimator, rather than the track length estimator used in EIRENE. 
The original reason for choosing the collision estimator in DEGAS 2
was that the charge exchange data file did not have the required
integrals of the cross section needed to score the energy and momentum
exchanges\cite{Reit93}. We have now worked out the forms these
integrals must have to emulate the EIRENE assumption of constant
$\sigma v$ and incorporated them into the charge exchange data file.
% Can add here variants of these expressions once we've documented
% the right expressions elsewhere.

The runs of Sec.~\ref{outofthebox} do not employ these estimator
changes (and utilize a different charge exchange cross section
altogether, as was noted previously), while runs in subsequent
sections have these modifications incorporated.

\subsection{Other Differences}
\label{otherdiffs}

We suspect that there are other differences remaining which will prevent
the two codes from being brought into closer agreement. For example,
DEGAS 2 uses a logarithmic interpolation in incident energy for the
Bateman format reflection data. EIRENE assumes a linear interpolation. 
The interpolation schemes in DEGAS 2 presently require a uniform
spacing, either linearly or logarithmically, between data points 
(a nonuniform spacing option was tried at one point, but abandoned
as too difficult to maintain). Hence, the interpolation technique
used by EIRENE cannot be emulated without significant changes to the
code. On the other hand, the subtlety of some of the differences which have
been detailed in this section suggests that remaining nonstatistical
discrepancies in the code results are unlikely to be of physical
significance and are consistent with minor numerical differences
between the codes.

\section{Characterize Comparison}

The principal means of comparing the code results is now considered to
be the quantitative prescription given in Sec.~\ref{statcompare}.
We will, however, begin this section with plots to illustrate two
points.

\subsection{Density Differences}

In Fig.~\ref{ddendiff}, we show the relative error,
\begin{equation}
\epsilon_{\rm rel} = \frac{n^{\rm D}_{D} - n^{\rm E}_{E}}{
\sigma_{\rm rsd}^{\rm D} \sqrt{(n^{\rm D}_{D})^{2} + (n^{\rm E}_{E})^{2}}}.
\end{equation}
This expression is equivalent to Eq.~(\ref{epsilonrel}) with $\sigma_{\rm rsd}$
being the relative standard deviation. Within the {\tt geomtesta} 
post-processor (see Sec.~\ref{maincode}), we do not have access to the
EIRENE error estimates (see Sec.~\ref{statcompare}). However, in the case
of the deuterium density, the EIRENE and DEGAS 2 relative standard
deviations are sufficiently similar that we can treat them as equal for
this purpose. Both are on the order of a one to a few percent near
the target. The scale in Fig.~\ref{ddendiff} indicates the number
of standard errors between the DEGAS 2 and EIRENE density values.
Note that this is a signed quantity. The white zones indicate areas
with differences greater than 3 times the standard error; the black 
areas indicate likewise that the lower bound has been exceeded.

\begin{figure}
% \epsfverbosetrue
\begin{center}
\leavevmode
\includegraphics[width=5in]{ddendiff}
\caption{Difference between DEGAS 2 and EIRENE deuterium density relative
to the standard error estimate.}
\label{ddendiff}
\end{center}
\end{figure}

The dominance of the green areas is indicative of the nearly complete
statistical agreement (see Sec.~\ref{quantcompare}). The red, orange,
and yellow zones near the target likely point to a remaining systematic
discrepancy between the codes, as suggested in Sec.~\ref{otherdiffs}.

\subsection{Comparison of Momentum Source}

As indicated by Figs.~\ref{dmomsrc} and \ref{eirmomsrc}, the deuterium
ion parallel momentum sources computed by the two codes are qualitatively
(i.e., visually) similar. However, the statistical errors are extremely 
high (Fig.~\ref{momrsd}), presumably caused by the tight coupling with
the background ions through charge exchange (i.e., the net momentum
source due to charge exchange is the difference of two nearly equal
numbers).

\begin{figure}
% \epsfverbosetrue
\begin{center}
\leavevmode
\includegraphics[width=5in]{d_mom_src}
\caption{Source of deuterium ion parallel momentum computed by DEGAS 2.}
\label{dmomsrc}
\end{center}
\end{figure}

\begin{figure}
% \epsfverbosetrue
\begin{center}
\leavevmode
\includegraphics[width=5in]{eir_mom_src}
\caption{Source of deuterium ion parallel momentum computed by EIRENE.}
\label{eirmomsrc}
\end{center}
\end{figure}

\begin{figure}
% \epsfverbosetrue
\begin{center}
\leavevmode
\includegraphics[width=5in]{mom_rsd}
\caption{Relative standard deviation for the deuterium ion momentum source
computed by DEGAS 2.}
\label{momrsd}
\end{center}
\end{figure}

The lack of regions with small relative standard deviations makes comparison
of the code results (see Sec.~\ref{statcompare}) problematic since we 
require a large number of (independent, ideally) zones with ``good
statistics'' (so that the Central Limit Theorem applies) for the analysis.
As will be seen in Sec.~\ref{quantcompare}, relaxing one of these two
constraints permits a comparison to be made. In each case, the results
are consistent with the two codes being in agreement. The deuterium ion 
energy source likewise has significant relative standard deviations, 
although smaller than those shown in Fig.~\ref{momrsd}. The consequences
for the stability of a coupling to a plasma transport code may be
more serious, however. Future work will assess and remedy this situation,
if necessary.

\subsection{Results of Quantitative Comparison}
\label{quantcompare}

The procedure outlined in Sec.~\ref{statcompare} must be amended slightly
to accomodate the remaining few-percent systematic differences in some
of the code results. The relative error is given instead by
\begin{equation}
\epsilon_{\rm rel} = \frac{\mu^{\rm D}_{K} - \mu^{\rm E}_{K^{\prime}}}
{\max{\left( \sqrt{(s^{\rm D}_{K})^{2} + (s^{\rm E}_{K^{\prime}})^{2}},
\sigma_{\rm min} \frac{\mu^{\rm D}_{K} + \mu^{\rm E}_{K^{\prime}}}{2}\right)}},
\end{equation}
where again the $s_{K}$ refers to the standard error for a simulation
of $K$ flights (i.e., $s_{K} = \sigma_{\rm rsd} \mu_{K}$). The 
additional parameter, $\sigma_{\rm min}$ is associated with the
statistical error (expressed as a fraction, like the relative
standard deviation). 

For all of these except the ion energy and momentum
source, $\sigma_{\rm max} = 20 \%$ and $\sigma_{\rm min}$ is given.
For those two, $\sigma_{\rm min} = 0$ and $\sigma_{\rm max}$ is given.
In the former case, the relative standard deviation is required to be
$< 20 \%$ to ensure that the Central Limit Theorem is applicable and
$\sigma_{\rm min}$ represents the systematic error. Of the tested
values of $\sigma_{\rm min}$ the one which yields a set of percentages
most closely matching the Maxwellian ideal ($68 \%$, $95 \%$, $99.7 \%$)
is taken as our estimate of the systematic error. For the ion energy
and momentum sources, we are unable to get ``good statistics'' with
$\sigma_{\rm max} = 20 \%$; hence, larger values are tested. The
relative standard deviations are too large for any systematic error
to be detected, so we use $\sigma_{\rm min} = 0$.

\begin{tabular}{|c||c|c|c||c|c|c|}
\hline
 Source:       & \multicolumn{3}{|c||}{Plate} & \multicolumn{3}{|c|}{Recombination} \\ \hline  % Do better?
        & $1 \sigma$ & $2 \sigma$ & $3 \sigma$ & $1 \sigma$ & $2 \sigma$ & $3 \sigma$ \\ \hline
\multicolumn{7}{|l|}{D Density} \\ \hline
$\sigma_{\rm min}=7 \%$ & $82 \%$ & $99.2 \%$ & $100  \%$ & $90 \%$ & $99.4 \%$ & $99.9 \%$ \\ \hline
$\sigma_{\rm min}=5 \%$ & $75 \%$ & $95   \%$ & $99.8 \%$ & $82 \%$ & $98   \%$ & $99.9 \%$ \\ \hline
$\sigma_{\rm min}=0 \%$ & $56 \%$ & $86   \%$ & $96   \%$ & $52 \%$ & $88   \%$ & $99   \%$ \\ \hline
\multicolumn{7}{|l|}{D$_{2}$ Density} \\ \hline
$\sigma_{\rm min}=0 \%$ & $64 \%$ & $92 \%$ & $98   \%$ & $59 \%$ & $92 \%$ & $99.2 \%$ \\ \hline
\multicolumn{7}{|l|}{D$_{2}^{+}$ Density} \\ \hline
$\sigma_{\rm min}=0 \%$ & $63 \%$ & $92 \%$ & $100  \%$ & $68 \%$ & $97  \%$ & $100  \%$ \\ \hline
\multicolumn{7}{|l|}{D$^{+}$ Ion Source Rate} \\ \hline
$\sigma_{\rm min}=7 \%$ & $74 \%$ & $99.6 \%$ & $100  \%$ & $83 \%$ & $99.3 \%$ & $100  \%$ \\ \hline
$\sigma_{\rm min}=5 \%$ & $57 \%$ & $94   \%$ & $100  \%$ & $63 \%$ & $98   \%$ & $99.8 \%$ \\ \hline
$\sigma_{\rm min}=0 \%$ & $38 \%$ & $73   \%$ & $93   \%$ & $38 \%$ & $66   \%$ & $85   \%$ \\ \hline
\multicolumn{7}{|l|}{Electron Energy Source Rate} \\ \hline
$\sigma_{\rm min}=7 \%$ & $92 \%$ & $99 \%$ & $99.1 \%$ & $93 \%$ & $99 \%$ & $99.3 \%$ \\ \hline
$\sigma_{\rm min}=5 \%$ & $79 \%$ & $98 \%$ & $99.1 \%$ & $86 \%$ & $98 \%$ & $99.1 \%$ \\ \hline
$\sigma_{\rm min}=0 \%$ & $38 \%$ & $72 \%$ & $91   \%$ & $45 \%$ & $78 \%$ & $94   \%$ \\ \hline
\multicolumn{7}{|l|}{D$^{+}$ Ion Energy Source Rate} \\ \hline
$\sigma_{\rm max}=70 \%$ & $69 \%$ & $94 \%$ & $99   \%$ & $68 \%$ & $96 \%$ & $99.7 \%$ \\ \hline
$\sigma_{\rm max}=50 \%$ & $69 \%$ & $94 \%$ & $99.1 \%$ & $65 \%$ & $97 \%$ & $100  \%$ \\ \hline
$\sigma_{\rm max}=20 \%$ & $68 \%$ & $94 \%$ & $100  \%$ & $61 \%$ & $98 \%$ & $100  \%$ \\ \hline
\multicolumn{7}{|l|}{D$^{+}$ Ion Parallel Momentum Source Rate} \\ \hline
$\sigma_{\rm max}=70 \%$ & $72 \%$ & $88 \%$ & $99.2 \%$ & $75 \%$ & $92  \%$ & $99.4 \%$ \\ \hline
$\sigma_{\rm max}=50 \%$ & $74 \%$ & $97 \%$ & $100 \%$ & $76 \%$ & $97  \%$ & $99   \%$ \\ \hline
$\sigma_{\rm max}=20 \%$ & $58 \%$ & $89 \%$ & $100 \%$ & $0  \%$ & $0   \%$ & $0    \%$ \\ \hline
\end{tabular}

The two molecular densities fare the best, given nearly Maxwellian percentages
without allowing for any systematic error. The atom density,  electron
energy source (which is given by the atom density times a function of $T_{e}$,
apart from a small contribution due to molecules), and ion
source rate are all roughly consistent
with a $5 \%$ systematic error. Again, the ion energy and momentum sources
are too noisy to permit an estimate of the systematic error. However, these
numbers indicate that the two codes do agree within the estimated
statistical errors.

\section{Performance Benchmark and Optimization}
\label{performance}

Although efficiency was kept in mind during all
stages of the design process, no effort was made to optimize the
performance of DEGAS 2 along the way.
In this section, we describe the benchmark of the code's performance against 
EIRENE. Speed bottlenecks were discovered through profiling. 
Improvements eliminating those bottlenecks.were implemented.
Some involved
significant code revisions; others were
effected by simple changes to source or
input files. The end result was that the run time for DEGAS 2 was 
roughly equal to
that of EIRENE within the variations normally experienced in a time-sharing
environment.

The implemented code revisions were motivated by profiling the code to
determine which subroutines were the most time-consuming. At one point,
assignments and comparisons of string variables were using significant
fractions of the run time. The strings responsible were 
being used to 
describe auxiliary data associated with the atomic physics reactions.
The code was modified to compile a list of all of these strings at the
beginning of the run and to use an integer pointer into that list for
the run-time comparisons and assignments. Because these code changes
were pervasive, their impact on the run time could not be documented
as carefully as the other improvements noted below. Roughly, however,
they resulted in a reduction of about 10 seconds per 
1000 flights.

\subsection{Starting Point}

The starting assumptions were the same as those associated with the
code as run for the initial EIRENE benchmark (specifically, the 
DEGAS 2 version with the CVS tag \verb+V1_8a+).
The ``EIRENE physics'' set was used to minimize differences due to the 
number of collisions and to permit direct comparisons of the variances.

For each code configuration, runs were performed at 1000, 2000, and
4000 flights. The time for the 1000 flight case was subtracted from
the other two and the results used to estimate the incremental
time required to track 1000 flights. The idea is that production
runs would be substantially longer than these (with relative standard
deviations on the order of $1 \%$) so that the overhead associated with
input and output would be negligible. In a few cases, production length
runs (e.g., 80,000 flights) were done; for those, the run time was just
taken to be the time consumed divided by 80. 

These runs were performed on a Sun UltraSPARC 2 running SunOS 5.5.1
and V. 4.2 of f77 with \verb+-O4+ optimization. Again, because this computer 
is a time-sharing system, these run times are only approximate. The 
estimated error is about $10 \%$. The baseline run time for
DEGAS 2 was: 136 seconds per 1000 flights.

\subsection{Charge Exchange Rejection}

The charge exchange cross section depends on the relative ion-neutral
velocity. However, collisions are decided upon using a cross section
averaged over a Maxwellian distribution. In the baseline, the ion
collision partners were chosen so as to be consistent with the actual
cross section using a rejection technique, sampling several ions 
before finding one which is satisfactory.
The version of EIRENE used in this benchmark does not do this 
(neither did the original DEGAS). 

For the
first change, charge exchange rejection was turned off. A cursory
examination of the results showed no significant impact on the neutral
density. However, a stronger effect may have occurred elsewhere (e.g.,
energy transferred to background species). Run time after disabling
charge exchange rejection: 119 seconds per 1000 flights.

\subsection{Reduce Number of Scores}

One impressive feature of EIRENE is how thoroughly its default operation
has been pared down to the minimum necessary for coupling to the fluid
plasma codes. In particular, no data on the variances are kept. DEGAS 2 was
written with the philosophy that the mean value for a score is meaningless
without a corresponding variance, and the two data values are kept
together. For this comparison, a short list of variances was requested
in the EIRENE input file. These will be needed later.

As presently distributed, DEGAS 2 by default sets up about 14 scores or
tallies. This list was reduced to 7 at this step (later, 3 more will
be eliminated leaving the neutral density and the 3 scores representing
the particle, momentum, and energy transfer to the background species).
Run time after cutting the number of scores to 7: 95 seconds per
1000 flights.

\subsection{Compression of Scores}

The first offender in the profiling exercise was the routine responsible
for compiling the scores accumulated during a single flight into the
global total. As originally written, both the total and incremental
scoring arrays were full-sized (roughly the
number of zones times the number of scores times the number of
background species). However, each flight visits only a small fraction
of the problem space, and this routine was spending a significant
amount of time adding $0$ to $0$.

These scoring arrays were modified to contain only the non-zero scores,
with an array of pointers mapping their contents back to the corresponding
locations in the full-sized arrays (which were still used for the
final output stage). Run time after implementation of compressed
scoring: 49 seconds per 1000 flights.

\subsection{Other Changes with Minimal Effects}

Further reducing the number of scores from  7 to 4 
had little impact on
the run, probably because of the use of
compression at this point.

DEGAS 2 (and EIRENE) employ a 10 term quadratic representation of all
surfaces in the geometry. 
In some problems, such as this one, only the linear terms
are needed. The inclusion of the quadratic terms was made a compile-time
option so that they could be turned off for fully linear geometries.
Earlier tests of the impact of this change showed a reduction of
about 4 seconds per 1000 flights. However, in this more systematic
series of trials, the improvement could not be quantified with certainty.

\subsection{Variance-Altering Changes}

Thus far, all of the changes made led to the same final results as
the run with ``EIRENE physics''.
Most of the subsequent changes,
while reducing the run time, also result in increases
in the variance of the results. In a given Monte Carlo calculation,
the variance is inversely proportional to the number of flights and, hence,
to the run time. So, the performance Figure of Merit (FOM) is the
variance times the run time.  An attempt will be made
below to quantify this FOM, but we can also compare qualitatively the
variance over the most relevant portion of the problem space (near the
target plate) via the relative standard deviation $\sigma_{\rm rsd}$. 
Figure~\ref{densityrsd} shows $\sigma_{\rm rsd}$ for the deuterium density
in the baseline case.

\begin{figure}
\begin{center}
\leavevmode
\includegraphics[width=5in]{spD_density_rsd_ex}
\caption{Relative standard deviation of the atomic deuterium density in the
baseline run of the performance benchmark.}
\label{densityrsd}
\end{center}
\end{figure}

Below we will use $\sigma_{\rm rsd}$ for the source rate of D$^{+}$ 
due to ionization
of neutrals. Although not computed for this baseline run, this
$\sigma_{\rm rsd}$ should
be comparable to that shown in Fig.~\ref{densityrsd} except for the 
first one or two zones adjacent to the plate
in which the contributions from D$_{2}$ are significant.

\subsection{Removal of Suppressed Ionization}

The simplest (``analog'') Monte Carlo simulation kills off neutrals
at their first ionizing collision. However, this makes finding the
solution more than one mean free path from the source difficult. The
predominant non-analog improvement is to assign a weight to each neutral
flight (e.g., initially 1) and reduce that weight at each step to reflect
the amount of ionization which should have occurred along the way. 
As a result, smaller variances are obtained in low probability regions of
the problem. Because each flight will thus be tracked for more steps, 
a run using suppressed ionization will take longer. Whether or not
that time is well spent depends on the problem at hand.

These EIRENE runs do not use suppressed ionization. For the purposes
of comparison, we turned it off in DEGAS 2 as well. The run time
without suppressed ionization dropped to 15 seconds per 1000 flights.
Figure~\ref{ionsourceisrsd} shows the resulting $\sigma_{\rm rsd}$.

\begin{figure}
\begin{center}
\leavevmode
\includegraphics[width=5in]{D__Ion_Source_is_rsd}
\caption{Relative standard deviation of the D$^{+}$ ion source {\em without}
suppressed ionization.}
\label{ionsourceisrsd}
\end{center}
\end{figure}

\subsection{Collision Estimator}

Some of the scores in DEGAS 2, e.g., the neutral density and the D$^{+}$ source
rate, are computed via the track-length estimator. The other scores
are compiled using data only at collisions (though most of these can also
be done by track-length). Since the collision routines get executed
regardless of whether or not the resulting data are used in the scores, we
felt that it would be interesting to switch all of the estimators to 
collision (except density) and eliminate the overhead associated with
generating the track-length scores. Of course, doing this increases the
variance as is shown in Fig.~\ref{ionsourcecolrsd}. The run with using
the collision estimator (again without suppressed absorption) needed
10 seconds per 1000 flights.

% NOTE: it's not clear from my original text which options were used
% in this section and the next - are all of these time reductions essentially
% additive?

\begin{figure}
\begin{center}
\leavevmode
\includegraphics[width=5in]{D__Ion_Source_col_rsd}
\caption{Relative standard deviation of the D$^{+}$ ion source using 
a collision estimator and {\em without}
suppressed ionization.}
\label{ionsourcecolrsd}
\end{center}
\end{figure}

\subsection{Russian Roulette for Molecules}

Two of the seven reactions involving D$_{2}$ and D$_{2}^{+}$ in 
DEGAS 2 result in two D
atoms. DEGAS 2 normally tracks both of these. EIRENE, however, like the
original DEGAS, chooses one of the two to follow and ``kills'' off the other.
This is a simple example of a general nonanalog Monte Carlo technique known
as ``Russian roulette''. Whether the use (or non-use) of this technique 
improves the variance again depends on the problem. For the purposes of
this benchmark, Russian roulette was added to DEGAS 2's molecular 
dissociation routine. Those scores switched to collision estimator in the
previous section were reverted to the track-length estimator for this
configuration. Figure~\ref{ionsourcerrrsd} shows the corresponding
$\sigma_{\rm rsd}$. This run used only 8 seconds per 1000 flights.

\begin{figure}
\begin{center}
\leavevmode
\includegraphics[width=5in]{D__Ion_Source_rr_rsd}
\caption{Relative standard deviation of the D$^{+}$ ion source using 
Russian roulette for molecules.}
\label{ionsourcerrrsd}
\end{center}
\end{figure}

\subsection{Figures of Merit}

We now have four different configurations with which to evaluate this
figure of merit, variance times run time. Presently, there is no single
``answer'' in these simulations which we can use to compute the FOM. 
Somewhat arbitrarily, we have selected a region of the problem spanning
the width (radius) of this geometry and encompassing most of the integrated
D$^{+}$ source (the selection was made by eyeballing the plot; $83 \%$ of the
integrated source was included). The $\sigma_{\rm rsd}$ for each
configuration was divided by that of the baseline and integrated over this
region. We then defined the FOM as the product of the 1000
flight run times (quoted above) and the square (to get the variance) of
this ratio.

\begin{tabular}{|c|c|c|c|}
\hline
Configuration & Seconds / 1000 flights & $\sigma_{\rm rsd}$ ratio & FOM \\ \hline 
Baseline      &             49         &     1.0               &  49 \\ \hline
No Suppressed Ionization &  15         &     1.9               &  54 \\ \hline
Collision Estimator      &  10         &     4.3               & 185 \\ \hline
D$_{2}$ Russian Roulette &   8         &     2.3               &  41 \\ \hline
\end{tabular}

Clearly, using the collision estimator is not a good idea. The effectiveness
of suppressed ionization could go either way. However, doing Russian
roulette on the molecular products (this run was done without ionization
suppression) looks to be the overall winner. Probably not by accident, this
configuration is the closest to the default mode of operation for EIRENE.

\subsection{EIRENE Performance}

Of course, the whole point of this portion of the benchmark is to compare
the performance of DEGAS 2 against that of EIRENE. Apart from the addition
of the variance output to EIRENE, no other modifications have been made to
its default mode of operation. However, the compiler optimization was
changed from the \verb+-O3+ value specified with the IPP-Garching version
to the \verb+-O4+ which appears to work well for DEGAS 2 with the Sun
FORTRAN 77 compiler. Over several runs of length 1000 to 10000 flights,
the incremental time for 1000 flights is 12 seconds. 
The relative standard deviations computed by EIRENE were consistent with
those of the corresponding DEGAS 2 configuration (Russian roulette), 
although a direct comparison was not made here (see
Sec.~\ref{estdiffs}). 

Within normal variations, the two codes are thus running at the same
speed. The particular DEGAS 2 configuration should be chosen according to
the needs of the problem at hand. Keep in mind that the original objective of
DEGAS 2 was to be faster than DEGAS; the hope was that the code would
yield performance comparable to that of EIRENE while retaining
flexibility. That objective has been met. DEGAS 2 is also able
to take full advantage of the substantial 
benefits of parallel processing.

Note, however, that EIRENE's computation of variances is less 
than optimal. Disabling them in the input file and turning off
unneeded tallies reduces the incremental run time to 3 seconds
per 1000 flights. At this point, EIRENE is substantially faster than 
DEGAS 2 since turning off
its variance computation (which has been optimized) will not have a significant
impact. The performance comparison between the two codes will
be revisited as part of a later benchmark utilizing a
more realistic divertor geometry.

The advantages of dynamic memory allocation of all
run-time arrays, another design feature of DEGAS 2, have also been
demonstrated during this benchmark. The run-time sizes were: DEGAS 2 -
7 MB, EIRENE - 142 MB. By reducing the dimensioning parameters in EIRENE to 
values more
appropriate for this input file, the size of the code was reduced from
142 MB to 55 MB.

\chapter{Troubleshooting}

\section{Help!}
\label{help}

George Wilkie \\
Princeton Plasma Physics Laboratory \\
P. O. Box 451 \\
Princeton NJ 08543-0451\\
(609) 243-2629 \\
\href{mailto:gwilkie@pppl.gov}{E-Mail: gwilkie@pppl.gov}

Daren Stotler \\
\href{mailto:dstotler@pppl.gov}{E-mail: dstotler@pppl.gov}

\begin{thebibliography}{99}
\bibitem{Tend87}The development of analytic, numerical, and Monte Carlo  
methods of solving
the time-independent Boltzmann equation describing neutral kinetics is
reviewed in: \\ M. Tendler and D. Heifetz, {\em Fusion Tech.} {\bf 11}, 289 
(1987).
\bibitem{Span69} The ``bible'' of Monte Carlo transport techniques is: \\
J. Spanier and E. M. Gelbard, ``MonteCarlo Principles and Neutron
Transport Problems'' (Addison-Wesley Publishing Company, Reading
Massachusetts, 1969). \\
Other references from the neutron transport field include: \\
J. M. Hammersley and D. C. Handscomb, ``Monte Carlo Methods'' (Metuchen,
London, 1964). \\
L. L. Carter and E. D. Cashwell, ``Particle-Transport Simulation with the
Monte Carlo Method'' (National Technical Information Service, Springfield,
Virginia, 1975). \\
Ivan Lux and Laszlo Koblinger, ``Monte Carlo Particle Transport Methods:
Neutron and Photon Calculations'' (CRC Press, Boca Raton, 1991).
\bibitem{Heif82} The Algorithm and applications of DEGAS are discussed in: \\
D. Heifetz, D. Post, M. Petravic, J. Weisheit, and G. Bateman,
{\em J. Comp. Phys.} {\bf 46} 309 (1982). \\
The Proceedings from a NATO-sponsored school in Val-Morin, Canada (1984)
are an invaluable reference for scrape-off layer physics. In addition
to again describing the algorithms and applications of DEGAS, Heifetz
covers in his article some of the basics of neutral transport in a plasma:
D. B. Heifetz, in {\em Physics of Plasma-Wall Interactions in Controlled 
Fusion}, D. Post and R. Behrisch, Eds., (Plenum, New York, 1986), p. 695.
\bibitem{Stot92} D. P. Stotler et al., {\em J. Nucl. Mater.}
{\bf 196-198}, 894 (1992).
\bibitem{Reit95} D.~Reiter et al., {\em J. Nucl. Mater.} {\bf 220--222}
987 (1995).
\bibitem{Bate61} D. R. Bates, A. E. Kingston, and R. W. P. McWhirter,
{\em Proc. R. Soc. A} {\bf 267}, 297 (1962).
\bibitem{McWh65} R. W. P. McWhirter, in {\em Plasma Diagnostic Techniques},
R. H. Huddlestone and S. L. Leonard, Eds. (Academic, New York, 1965).
\bibitem{FWEB} \verb+http://w3.pppl.gov/~krommes/fweb_toc.html+
(\href{http://w3.pppl.gov/~krommes/fweb_toc.html}{FWEB home page})
\bibitem{nCDF} \verb+http://www.unidata.ucar.edu/software/netcdf/+
(\href{http://www.unidata.ucar.edu/software/netcdf/}{netCDF home page})
\bibitem{HDF} \verb+http://www.hdfgroup.org/+
(\href{http://www.hdfgroup.org/}{HDF home page})
\bibitem{Silo} \verb+https://wci.llnl.gov/codes/silo/index.html+
(\href{https://wci.llnl.gov/codes/silo/index.html}{Silo home page})
\bibitem{Reit92} ``The EIRENE Code User Manual'' (D. Reiter,
Forschungszentrum J\"{u}lich, 2017), 
\href{http://www.eirene.de/index.html}{EIRENE home page}.
information.
\bibitem{Case67} K. M. Case
and P. F. Zweifel, {\em Linear Transport Theory} (Addison-Wesley, Reading, MA,
1967).
\bibitem{Marc96} R. Marchand and M. Dumberry, {\em Comp. Phys. Comm.}
{\bf 96}, 232 (1996).
\bibitem{Stot96} D. P. Stotler, A. Yu. Pigarov, C. F. F. Karney, 
S. I. Krasheninnikov, B. LaBombard, B. Lipschultz, G. M. McCracken,
A. Niemczewski, J. A. Snipes, J. L. Terry, and R. A. Vesey, in {\em 
Fusion Energy 1996} (Proc. 16th Int. Conf. Montreal, 1996), Vol. 2 
(IAEA, Vienna, 1997) p. 633 (see also the corresponding
\href{http://bp.pppl.gov/pub_report/1997/PPPL-3221-abs.html}{PPPL Report}).
  \bibitem{Jane87} R. K. Janev, W. D. Langer, K. Evans, Jr., and D. E. Post, Jr., {\em Elementary
Processes in Hydrogen-Helium Plasmas} (Springer-Verlag, Berlin, 1987).
\bibitem{Reit93} D. Reiter, in {\em Atomic and Plasma-Material Interaction
Processes in Controlled Thermonuclear Fusion}, R. K. Janev and H. W.
Drawin, Eds., (Elsevier, New York, 1993), p. 243.
\bibitem{Jane93} R. K. Janev and J. J. Smith, {\em Atomic and Plasma-Material
Interaction Data for Fusion} (Supplement to the journal Nuclear
Fusion) {\bf 4} (1993).
\bibitem{John73} L. C. Johnson and E. Hinnov, {\em J. Quant. Spectrosc.
Radiat. Transfer} {\bf 13}, 333 (1973).
\bibitem{Weis75} J. C. Weisheit, {\em J. Phys. B: Atom. Molec. Phys.} {\bf 8},
2556 (1975).
\bibitem{Goto2003} M. Goto, {\em J. Quant. Spectrosc. Radiat.
Transfer} {\bf 76}, 331 (2003).
\bibitem{Fuji79} T. Fujimoto,  {\em J. Quant. Spectrosc. Radiat. 
Transfer} {\bf 21}, 439 (1979).
\bibitem{Kanz00} R. J. Kanzleiter, D. P. Stotler, C. F. F. Karney, and 
D. Steiner {\em Phys. Plasmas} {\bf 7}, 5064 (2000).
\bibitem{Krst98} P. S. Krstic and D. R. Schultz, {\em Atomic and 
Plasma-Material Data for Fusion} {\bf 8}, 1 (1998).
\bibitem{Bach95} P. Bachmann and D. Reiter, {\em Contrib. Plasma Phys.} 
{\bf 35}, 45 (1995).
\bibitem{Abou91} H. H. Abou-Gabal and G. A. Emmert, {\em Nucl. Fusion}
{\bf 31}, 407 (1991).
\bibitem{Dalg58} A. Dalgarno, {\em Phil. Trans. R. Soc. London, Ser. A}
{\bf 250}, 426 (1958).
\bibitem{Hols52} T. Holstein, {\em J. Phys. Chem.} {\bf 56}, 832 (1952).
\bibitem{Wade79} J. M. Wadehra, {\em Phys. Rev. A} {\bf 20}, 1859 (1979).
\bibitem{Barn90}C. F. Barnett, {\em Atomic Data for Fusion}, Oak Ridge
National Laboratory Report ORNL-6086, Vol. 1 (1990).
\bibitem{Newm82} J. H. Newman, J. D. Cogan, D. L. Ziegler et al., {\em Phys.
Rev. A} {\bf 25}, 2976 (1982).
\bibitem{Revi71} A. C. Reviere, {\em Nucl. Fusion} {\bf 11}, 363 (1971). 
\bibitem{OHan89}J. F. O'Hanlon, {\em A User's Guide to Vacuum Technology}
(John Wiley \& Sons, New York, 1989).
\bibitem{Niem95} A. Niemczewski, Ph. D. Thesis, Massachusetts Institute
of Technology, Plasma Fusion Center Report PFC/RR-95-8 (1995).
\bibitem{Bhat54} P. L. Bhatnagar, E. P. Gross, and M. Krook, {\em Phys.
Rev.} {\bf 94}, 511 (1954).
\bibitem{Wela54} P. Welander, {\em Arkiv Fysik} {\bf 7}, 507 (1954).
\bibitem{Mayxx} Chr. May, Ph. D. Thesis. 
\bibitem{Reit97} D. Reiter, Chr. May,
M. Baelmans, and P. B\"{o}rner, {\em J. Nucl. Mater.} {\bf 241}, 342 (1997).
\bibitem{Kanz99} R. J. Kanzleiter, Ph. D. Thesis, Rensselaer Polytechnic
Institute (1999).
\bibitem{Chap70} S. Chapman and T. G. Cowling, {\em The Mathematical
Theory of Non-Uniform Gases} (Cambridge University Press, 1970).
\bibitem{Mors63} T. F. Morse, {\em Phys. Fluids} {\bf 6}, 1420 (1963).
\bibitem{Reid87} R. C. Reid, J. M. Prausnitz, B. E. Poling, {\em The
Properties of Liquids and Gases} (McGraw-Hill Book Company, New York, 1987).
\bibitem{CRC} {\em CRC Handbook of Chemistry and Physics, 75th Edition,
1913--1995}, D. R. Lide, Ed., p. 6-239 (CRC Press, Boca Raton, 1994).
\bibitem{Chan2008} C. S. Chang, S. Klasky, J. Cummings et al., 
{\em J. Phys.: Conf. Ser.} {\bf 125}, 012040 (2008).
\bibitem{Krst99} P. S. Krstic and D. R. Schultz, {\em Phys. Rev. A}
{\bf 60}, 2118 (1999).
\bibitem{ADAS} H. P. Summers, The ADAS User Manual, Version
2.6, 2002.  \href{http://www.adas.ac.uk/}{http://www.adas.ac.uk/}
\bibitem{Stot2017} D. P. Stotler, D. J. Battaglia, R. Hager et al.,
{\em Nucl. Mater. Energy} {\bf 12}, 1130 (2017).
\bibitem{Will62} D. R. Willis, {\em Phys. Fluids} {\bf 5}, 127 (1962).
\bibitem{Cerc69} C. Cercignani, {\em J. Stat. Phys.} {\bf 1}, 297 (1969).
\bibitem{Cerc75} C. Cercignani, {\em Theory and Application of the Boltzmann
Equation} (Scottish Academic Press, New York, 1975).
\bibitem{Bevi69} P. R. Bevington, {\em Data Reduction and Error Analysis
for the Physical Sciences} (McGraw-Hill Book Company, New York, 1969).
\bibitem{Ecks91} W. Eckstein, {\em Computer Simulation of Ion-Solid
Interactions} (Springer-Verlag, New York, 1991).
\bibitem{Bate80} G. Bateman, ``Distribution of Neutrals Scattered Off 
A Wall'', \href{http://w3.pppl.gov/degas2/Bateman_distribution.pdf}{PPPL Applied Physics Division Report \# 1 (1980)}.
\bibitem{Macm67} D. B. Macmillan, {\em SIAM J. Appl. Math.} {\bf 15}, 264 
(1967).
\bibitem{Span66} J. Spanier, {\em SIAM J. Appl. Math.} {\bf 14}, 702 (1966).
\bibitem{Bran2014} S. Brandt, {\em Data Analysis: Statistical and
Computational Methods for Scientists and Engineers} (Springer International
Publishing, Switzerland, 2014), 
\href{https://link.springer.com/book/10.1007%2F978-3-319-03762-2}{online}.
\bibitem{MCNP} X-5 Monte Carlo Team, ``MCNP - A General Monte Carlo N-Particle Transport
Code, Version 5, Volume I: Overview and Theory'', Los Alamos National
Laboratory Report \href{https://laws.lanl.gov/vhosts/mcnp.lanl.gov/pdf_files/la-ur-03-1987.pdf}{LA-UR-03-1987} 
(2008).
\bibitem{Knuth97} D. E. Knuth, {\em The Art of Computer Programming} 
(Addison-Wesley, Reading, MA, 1997).
\bibitem{Pitc00} C. S. Pitcher, C. J. Boswell, J. A. Goetz et al., {\em Phys.
Plasmas} {\bf 7}, 1894 (2000).
\bibitem{Pitc01} C. S. Pitcher, C. J. Boswell, T. Chung
et al., {\em J. Nucl. Mater.} {\bf 290-293}, 812 (2001).
\bibitem{Stot01} D. P. Stotler, C. S. Pitcher, C. J. Boswell et al.,
{\em J. Nucl. Mater.} {\bf 290-293}, 967 (2001).
\bibitem{Stot2007} D. P. Stotler,
J. Boedo, B. LeBlanc, R. J. Maqueda, and S. J. Zweben, 
{\em J. Nucl. Mater.} {\bf 363-365}, 686 (2007).
\bibitem{Stot2004} D. P. Stotler, D. A. D'Ippolito, B. LeBlanc, 
R. J. Maqueda, J. R. Myra, S. A. Sabbagh,
and S. J. Zweben, {\em Contrib. Plasma Phys.} {\bf 44}, 294 (2004).
\bibitem{Zarn2001} M. C. Zarnstorff, L. A. Berry, A. Brooks et al.,
{\em Plasma Phys. Control. Fusion} {\bf 43}, A237 (2001). 
\end{thebibliography}

%%tth:\printindex

\begin{theindex}

  \item adaswrite, \hyperpage{66}
  \item adsorption, \hyperpage{72}
  \item Allocation, memory, \hyperpage{3}
  \item Atomic physics, \hyperpage{29}

  \indexspace

  \item background species, \hyperpage{72}
  \item BGK, \hyperpage{48}
  \item Boltzmann equation, \hyperpage{1}
  \item boxgen, \hyperpage{66}

  \indexspace

  \item Carre, \hyperpage{69}, \hyperpage{99}
  \item Charge exchange, \hyperpage{38}
  \item chargex, \hyperpage{71}
  \item checkpoints, \hyperpage{109}
  \item collisional-radiative, \hyperpage{29}
  \item CVS, \hyperpage{3}

  \indexspace

  \item dataexam, \hyperpage{67}
  \item datamatch, \hyperpage{68}
  \item datasetup, \hyperpage{66}
  \item defineback, \hyperpage{28}, \hyperpage{63}, \hyperpage{99}
  \item defineback input, \hyperpage{73}
  \item definegeometry2d, \hyperpage{61}, \hyperpage{63}, 
		\hyperpage{99}
  \item definegeometry2d input, \hyperpage{73}
  \item DEGAS, \hyperpage{2}
  \item desorption, \hyperpage{72}
  \item DG, \hyperpage{69}, \hyperpage{99}
  \item Differential cross section, \hyperpage{38}
  \item Diffusion, \hyperpage{48}
  \item Diffusion cross section, \hyperpage{38}
  \item directories, \hyperpage{62}
  \item dissoc, \hyperpage{71}
  \item dissoc\_rec, \hyperpage{71}
  \item Distribution, probability, \hyperpage{5}

  \indexspace

  \item efit2dg2d, \hyperpage{66}
  \item elastic, \hyperpage{71}
  \item Elastic scattering, \hyperpage{38}
  \item elements\_infile, \hyperpage{71}
  \item elementsfile, \hyperpage{71}
  \item Emacs, \hyperpage{59}
  \item emacs, \hyperpage{59}

  \indexspace

  \item flighttest, \hyperpage{63}
  \item ftangle, \hyperpage{60}
  \item fweave, \hyperpage{60}
  \item FWEB, \hyperpage{3}, \hyperpage{60}

  \indexspace

  \item Generators, random number, \hyperpage{5}
  \item generic, \hyperpage{71}
  \item generic PMI, \hyperpage{72}
  \item generic reactions, \hyperpage{71}
  \item geometry.inp, \hyperpage{74}
  \item geomtesta, \hyperpage{63}
  \item gmake, \hyperpage{59}

  \indexspace

  \item HDF, \hyperpage{3}, \hyperpage{62}
  \item HDF5, \hyperpage{62}
  \item HOME, \hyperpage{62}

  \indexspace

  \item ionization, \hyperpage{29}
  \item ionize, \hyperpage{71}
  \item ionize\_suppress, \hyperpage{71}

  \indexspace

  \item LaTeX, \hyperpage{60}

  \indexspace

  \item Macros, \hyperpage{3}
  \item make, \hyperpage{59}
  \item Makefile, \hyperpage{59}
  \item Makefile.local, \hyperpage{62}
  \item matcheir, \hyperpage{68}
  \item matchout, \hyperpage{68}
  \item materials\_infile, \hyperpage{72}
  \item materialsfile, \hyperpage{72}
  \item Maxwell molecules, \hyperpage{48}
  \item Monte Carlo, \hyperpage{5}

  \indexspace

  \item ncdum, \hyperpage{59}
  \item ncdump-filter, \hyperpage{59}
  \item ncgen, \hyperpage{59}
  \item ncgen-filter, \hyperpage{59}
  \item NetCDF, \hyperpage{3}
  \item netCDF, \hyperpage{59}
  \item Neutral Transport, \hyperpage{1}
  \item Neutral-ion scattering, \hyperpage{38}
  \item Neutral-neutral scattering, \hyperpage{48}
  \item non-generic PMI, \hyperpage{72}
  \item non-generic reactions, \hyperpage{71}

  \indexspace

  \item Object oriented programming, \hyperpage{3}
  \item outputscript, \hyperpage{74}

  \indexspace

  \item PMI, \hyperpage{72}
  \item pmi\_infile, \hyperpage{72}
  \item pmitest, \hyperpage{67}
  \item pmiwrite, \hyperpage{66}
  \item problem\_infile, \hyperpage{72}
  \item problemfile, \hyperpage{72}
  \item problemsetup, \hyperpage{63}
  \item Processing, parallel, \hyperpage{3}

  \indexspace

  \item Random, \hyperpage{5}
  \item randomtest, \hyperpage{67}
  \item ratecalc, \hyperpage{66}
  \item rates, atomic physics, \hyperpage{29}
  \item reaction type, \hyperpage{71}
  \item reaction\_infile, \hyperpage{71}
  \item reactionfile, \hyperpage{71}
  \item reactiontest, \hyperpage{67}
  \item reactionwrite, \hyperpage{66}
  \item readbackground, \hyperpage{63}
  \item readbackground input, \hyperpage{73}
  \item readgeometry, \hyperpage{63}
  \item readgeometry input, \hyperpage{73}
  \item recombination, \hyperpage{29}, \hyperpage{71}
  \item reference level, \hyperpage{72}
  \item reflection, \hyperpage{72}
  \item restarting, \hyperpage{110}

  \indexspace

  \item Scattering angle, \hyperpage{38}
  \item Scoring, \hyperpage{5}
  \item Silo, \hyperpage{62}
  \item snapshot, \hyperpage{28}
  \item sourcetest, \hyperpage{67}
  \item species, \hyperpage{71}
  \item species\_infile, \hyperpage{71}
  \item speciesfile, \hyperpage{71}
  \item Spin exchange, \hyperpage{38}
  \item subset data, \hyperpage{72}
  \item symbolic name, \hyperpage{70}
  \item sysdeptest, \hyperpage{67}

  \indexspace

  \item Tags, \hyperpage{59}
  \item tally, \hyperpage{74}
  \item tally\_infile, \hyperpage{74}
  \item tallyfile, \hyperpage{74}
  \item tallysetup, \hyperpage{63}, \hyperpage{74}
  \item test species, \hyperpage{72}
  \item TeX, \hyperpage{60}
  \item Theory, transport, \hyperpage{1}
  \item time dependence, \hyperpage{28}, \hyperpage{109}
  \item Total cross section, \hyperpage{38}
  \item Tracking, \hyperpage{5}
  \item Transport codes, \hyperpage{1}
  \item Triangle, \hyperpage{61}

  \indexspace

  \item ucd\_plot, \hyperpage{63}

  \indexspace

  \item Viscosity, \hyperpage{48}
  \item Viscosity cross section, \hyperpage{38}

\end{theindex}


\tableofcontents

\end{document}
